% Reconstitution chronologique de l'histoire des fonds océaniques — SVT 1ère TI — Cours signé OnBuch+
% Programme officiel MINESEC. Compiler avec : tectonic NCT10-chronologie-fonds-oceaniques.tex
\newcommand{\DOCMATIERE}{SVT}
\newcommand{\DOCNIVEAU}{1ère TI}
\newcommand{\DOCMODULE}{Module 3 : L'Éducation à l'Environnement et au Développement Durable}
\newcommand{\DOCLECON}{NCT10}
\newcommand{\DOCTITRE}{Reconstitution chronologique de l'histoire des fonds océaniques}
% OnBuch+ — préambule « cours signés » ENRICHI (compilé avec tectonic / XeLaTeX)
% Système visuel « Pop » d'OnBuch+ : fond crème (jamais blanc), bordures encre
% épaisses, ombres « dures » décalées, titres Archivo Black en capitales.
% Version MATHÉMATIQUES : graphes (pgfplots/tikz), boîtes pédagogiques
% (définition, propriété, méthode, attention, le savais-tu, exercice résolu,
% exercice, TP, bilan), page de garde et signature OnBuch+ ancrée en bas.
%
% Macros attendues dans main.tex AVANT \input{preamble} :
%   \DOCMATIERE  \DOCNIVEAU  \DOCMODULE  \DOCLECON (numéro)  \DOCTITRE
\documentclass[11pt]{article}

\usepackage{fontspec}
\usepackage[french]{babel}
\usepackage[a4paper,margin=2cm,top=3.4cm,bottom=2.8cm,headheight=1.4cm,footskip=1.3cm]{geometry}
\usepackage{amsmath,amssymb,amsfonts}
\usepackage{mathtools} % \xmapsto, \coloneqq... souvent utilisés par les modèles
\usepackage{cancel} % simplifications barrées (bilans de forces)
% \abs{z} (module, valeur absolue) et \norm{u} (norme), utilisés par les modèles
\providecommand{\abs}{}\let\abs\relax\DeclarePairedDelimiter{\abs}{\lvert}{\rvert}
\providecommand{\norm}{}\let\norm\relax\DeclarePairedDelimiter{\norm}{\lVert}{\rVert}
\usepackage[table]{xcolor} % [table] : fournit aussi \rowcolors (alternance de lignes)
\usepackage{pagecolor}
\usepackage{tikz}
\usetikzlibrary{arrows.meta,positioning,calc,shapes.geometric,shapes.misc,shapes.arrows,decorations.pathmorphing,decorations.pathreplacing,decorations.markings,patterns,shadows,mindmap,trees,fit,backgrounds,matrix,intersections,angles,quotes}
\usepackage{pgfplots}
\usepackage[version=4]{mhchem} % équations nucléaires \ce{^{235}_{92}U}
\usepackage[european,siunitx]{circuitikz} % schémas électriques
\pgfplotsset{compat=1.18}
\usepgfplotslibrary{fillbetween,groupplots} % aires entre courbes (calcul intégral)
\usepackage{enumitem}
\usepackage{fancyhdr}
% [most] charge toutes les bibliothèques tcolorbox (listings, minted, raster,
% documentation...) et gonfle inutilement la mémoire TeX sur un document long
% (~300 boîtes) : on ne charge que ce qui est réellement utilisé.
\usepackage[skins,breakable]{tcolorbox}
\usepackage{graphicx}
\usepackage{colortbl}
\usepackage{booktabs}
\usepackage{tabularx}
\usepackage{diagbox} % case d'en-tête diagonale des tableaux à double entrée
% Colonnes extensibles centrée / gauche / droite (C, L, R), souvent utilisées par les modèles
\newcolumntype{C}{>{\centering\arraybackslash}X}
\newcolumntype{L}{>{\raggedright\arraybackslash}X}
\newcolumntype{R}{>{\raggedleft\arraybackslash}X}
\usepackage{array}
\usepackage{multicol}
\usepackage{siunitx}
% parse-numbers=false : accepte aussi \SI{2,5 \times 10^{-3}}{...}
\sisetup{locale=FR,output-decimal-marker={,},group-minimum-digits=4,parse-numbers=false}
\DeclareSIUnit\ohm{\ensuremath{\symup{\Omega}}} % U+2126 absent de la police
\DeclareSIUnit\division{div} % réglages d'oscilloscope (ms/div, V/div)
\DeclareSIUnit\tr{tr} % tours (vitesse de rotation, ex. \SI{1500}{\tr\per\minute})
\DeclareSIUnit\ch{ch} % cheval-vapeur (puissance moteur, unité usuelle hors SI)
\DeclareSIUnit\franccfa{FCFA} % franc CFA (coûts, contexte camerounais)
\DeclareSIUnit\dioptre{\ensuremath{\delta}} % vergence d'une lentille (dioptrie)
\usepackage{qrcode}
% \degree : commande fréquemment utilisée par les modèles pour un angle ou une
% température en dehors de \SI{}{} (non fournie par les paquets chargés).
\providecommand{\degree}{^{\circ}}
% \qed (fin de démonstration), utilisé par les modèles sans amsthm
\providecommand{\qed}{\hfill$\square$}
% \barre : double barre des valeurs interdites dans un tableau de variations (tkz-tab)
\providecommand{\barre}{\ensuremath{\|}}
% environnement proof (amsthm non chargé) utilisé par les modèles
\ifdefined\proof\else\newenvironment{proof}[1][Démonstration]{\par\noindent\textit{#1.}\ }{\qed\par}\fi
\usepackage{lastpage}
\usepackage{hyperref}
\hypersetup{hidelinks}

% ============================================================
% Palette « Pop » OnBuch+ — mêmes hex que la classe P (plus_ui.dart)
% ============================================================
\definecolor{poporange}{HTML}{F59321}
\definecolor{popdark}{HTML}{E07A0C}
\definecolor{popcream}{HTML}{FFF6E8}
\definecolor{popink}{HTML}{1C1714}
\definecolor{poppurple}{HTML}{7A5AE0}
\definecolor{popgreen}{HTML}{2E9E6A}
\definecolor{poppink}{HTML}{E0457B}
\definecolor{popblue}{HTML}{2F7DE1}
\definecolor{popgold}{HTML}{C98A12}
\definecolor{popmuted}{HTML}{8A7F76}
\definecolor{popline}{HTML}{EADFD2}
\definecolor{popgray}{HTML}{8A7F76}
\colorlet{popgrey}{popgray}
% Alias tolérants (le modèle nomme parfois les couleurs différemment)
\colorlet{popwhite}{white}
\colorlet{popblack}{popink}
\colorlet{popred}{poppink}
\colorlet{popyellow}{popgold}
% Teintes claires pour fonds de tableaux / zones
\colorlet{popblueL}{popblue!10!white}
\colorlet{poporangeL}{poporange!14!white}
\colorlet{popgreenL}{popgreen!12!white}
\colorlet{poppurpleL}{poppurple!12!white}
\colorlet{poppinkL}{poppink!12!white}
\colorlet{popgoldL}{popgold!14!white}

\pagecolor{popcream}

% ============================================================
% Typographie
% ============================================================
\defaultfontfeatures{Ligatures=TeX}
\setmainfont{PlusJakartaSans-Medium.ttf}[Path=../../../fonts/,BoldFont=PlusJakartaSans-Bold.ttf]
% Maths en sans-serif (Fira Math) avec chiffres et lettres droites de Plus
% Jakarta, pour que les formules se fondent dans le texte.
\usepackage[math-style=ISO,bold-style=ISO]{unicode-math}
\setmathfont{FiraMath-Regular.otf}[Path=../../../fonts/]
\setmathfont{PlusJakartaSans-Medium.ttf}[Path=../../../fonts/,range={up/num,up/{Latin,latin},\mathup}]
\usepackage{icomma} % virgule décimale sans espace en mode math
% Caractères mathématiques Unicode absents des polices de texte (Plus Jakarta,
% Archivo Black) : rendus par leur équivalent mathématique, en texte comme en
% formule — sinon ils s'affichent en carré vide (ex. « Arithmétique dans ℤ »).
\usepackage{newunicodechar}
\newunicodechar{ℝ}{\ensuremath{\mathbb{R}}}
\newunicodechar{ℤ}{\ensuremath{\mathbb{Z}}}
\newunicodechar{ℂ}{\ensuremath{\mathbb{C}}}
\newunicodechar{ℕ}{\ensuremath{\mathbb{N}}}
\newunicodechar{ℚ}{\ensuremath{\mathbb{Q}}}
\newunicodechar{𝔻}{\ensuremath{\mathbb{D}}}
\newunicodechar{α}{\ensuremath{\alpha}}
\newunicodechar{β}{\ensuremath{\beta}}
\newunicodechar{γ}{\ensuremath{\gamma}}
\newunicodechar{δ}{\ensuremath{\delta}}
\newunicodechar{ε}{\ensuremath{\varepsilon}}
\newunicodechar{θ}{\ensuremath{\theta}}
\newunicodechar{λ}{\ensuremath{\lambda}}
\newunicodechar{μ}{\ensuremath{\mu}}
\newunicodechar{σ}{\ensuremath{\sigma}}
\newunicodechar{τ}{\ensuremath{\tau}}
\newunicodechar{φ}{\ensuremath{\varphi}}
\newunicodechar{ψ}{\ensuremath{\psi}}
\newunicodechar{ω}{\ensuremath{\omega}}
\newunicodechar{Σ}{\ensuremath{\Sigma}}
% majuscules produites par \MakeUppercase dans les titres (α-aminés -> Α)
\newunicodechar{Α}{\ensuremath{\alpha}}
\newunicodechar{Β}{\ensuremath{\beta}}
\newunicodechar{Γ}{\ensuremath{\Gamma}}
\newunicodechar{Δ}{\ensuremath{\Delta}}
\newunicodechar{Ε}{\ensuremath{\varepsilon}}
\newunicodechar{Θ}{\ensuremath{\Theta}}
\newunicodechar{Λ}{\ensuremath{\Lambda}}
\newunicodechar{Μ}{\ensuremath{\mu}}
\newunicodechar{Τ}{\ensuremath{\tau}}
\newunicodechar{Φ}{\ensuremath{\Phi}}
\newunicodechar{Ψ}{\ensuremath{\Psi}}
\newunicodechar{Ω}{\ensuremath{\Omega}}
\newunicodechar{↦}{\ensuremath{\mapsto}}
\newunicodechar{∈}{\ensuremath{\in}}
\newunicodechar{∉}{\ensuremath{\notin}}
\newunicodechar{∧}{\ensuremath{\wedge}}
\newunicodechar{⇔}{\ensuremath{\Leftrightarrow}}
\newunicodechar{⟺}{\ensuremath{\Longleftrightarrow}}
\newunicodechar{⇒}{\ensuremath{\Rightarrow}}
\newunicodechar{⟹}{\ensuremath{\Longrightarrow}}
\newunicodechar{∘}{\ensuremath{\circ}}
\newunicodechar{∪}{\ensuremath{\cup}}
\newunicodechar{∩}{\ensuremath{\cap}}
\newunicodechar{≡}{\ensuremath{\equiv}}
\newunicodechar{⊂}{\ensuremath{\subset}}
\newunicodechar{⊥}{\ensuremath{\perp}}
\newunicodechar{∃}{\ensuremath{\exists}}
\newunicodechar{∀}{\ensuremath{\forall}}
\newunicodechar{∛}{\ensuremath{\sqrt[3]{\,}}}
\newunicodechar{∜}{\ensuremath{\sqrt[4]{\,}}}
\newunicodechar{⃗}{\ensuremath{{}^{\to}}}
\newunicodechar{ⁿ}{\ifmmode^{n}\else\textsuperscript{\ensuremath{n}}\fi}
\newunicodechar{ˣ}{\ifmmode^{x}\else\textsuperscript{\ensuremath{x}}\fi}
\newunicodechar{ᵇ}{\ifmmode^{b}\else\textsuperscript{\ensuremath{b}}\fi}
\newunicodechar{ʳ}{\ifmmode^{r}\else\textsuperscript{\ensuremath{r}}\fi}
\newunicodechar{ᵘ}{\ifmmode^{u}\else\textsuperscript{\ensuremath{u}}\fi}
\newunicodechar{ʸ}{\ifmmode^{y}\else\textsuperscript{\ensuremath{y}}\fi}
\newunicodechar{ᵃ}{\ifmmode^{a}\else\textsuperscript{\ensuremath{a}}\fi}
\newunicodechar{ᵉ}{\ifmmode^{e}\else\textsuperscript{\ensuremath{e}}\fi}
\newunicodechar{ᵏ}{\ifmmode^{k}\else\textsuperscript{\ensuremath{k}}\fi}
\newunicodechar{ᵖ}{\ifmmode^{p}\else\textsuperscript{\ensuremath{p}}\fi}
\newunicodechar{ᴺ}{\ifmmode^{N}\else\textsuperscript{\ensuremath{N}}\fi}
\newunicodechar{ᶜ}{\ifmmode^{c}\else\textsuperscript{\ensuremath{c}}\fi}
\newunicodechar{ʰ}{\ifmmode^{h}\else\textsuperscript{\ensuremath{h}}\fi}
\newunicodechar{ᵠ}{\ifmmode^{q}\else\textsuperscript{\ensuremath{q}}\fi}
\newunicodechar{ₙ}{\ifmmode_{n}\else\textsubscript{\ensuremath{n}}\fi}
\newunicodechar{ₐ}{\ifmmode_{a}\else\textsubscript{\ensuremath{a}}\fi}
\newunicodechar{ᵢ}{\ifmmode_{i}\else\textsubscript{\ensuremath{i}}\fi}
\newunicodechar{ᵣ}{\ifmmode_{r}\else\textsubscript{\ensuremath{r}}\fi}
\newunicodechar{ₖ}{\ifmmode_{k}\else\textsubscript{\ensuremath{k}}\fi}
\newunicodechar{ᵦ}{\ifmmode_{\beta}\else\textsubscript{\ensuremath{\beta}}\fi}
\newunicodechar{ₓ}{\ifmmode_{x}\else\textsubscript{\ensuremath{x}}\fi}
\newfontfamily\popdisplay{ArchivoBlack-Regular.ttf}[Path=../../../fonts/]
\newfontfamily\poplogo{SpaceGrotesk-Bold.ttf}[Path=../../../fonts/]
\newfontfamily\popsemi{PlusJakartaSans-SemiBold.ttf}[Path=../../../fonts/]
\newfontfamily\popxbold{PlusJakartaSans-ExtraBold.ttf}[Path=../../../fonts/]
\newfontfamily\popxboldls{PlusJakartaSans-ExtraBold.ttf}[Path=../../../fonts/,LetterSpace=3.0]

\setlist[enumerate]{leftmargin=1.7em,itemsep=5pt,topsep=4pt}
\setlist[itemize]{leftmargin=1.7em,itemsep=4pt,topsep=4pt,label={\color{poporange}\textbullet}}
\setlength{\parindent}{0pt}
\setlength{\parskip}{5pt}
\renewcommand{\arraystretch}{1.25}

% pgfplots : style Pop par défaut
\pgfplotsset{
  every axis/.append style={axis line style={popink,line width=1pt},tick style={popink},
    grid style={popline,line width=0.5pt},label style={font=\small},tick label style={font=\footnotesize}},
  popcurve/.style={poporange,line width=1.8pt,no markers,smooth},
}
% Même style disponible dans un \draw TikZ simple (hors pgfplots)
\tikzset{popcurve/.style={poporange,line width=1.8pt}}
\tikzset{popgrid/.style={popline,very thin}}
\colorlet{popcurve}{poporange} % le nom du style est parfois utilisé comme couleur

% ============================================================
% Pill / squiggle
% ============================================================
\newcommand{\poppill}[3]{%
  \tikz[baseline=-0.55ex]\node[draw=popink,line width=1.2pt,rounded corners=2.6pt,%
    fill=#2,inner xsep=6.5pt,inner ysep=3pt]{%
    \popxboldls\fontsize{7}{7}\selectfont\color{#3}\MakeUppercase{#1}};%
}
\newcommand{\popsquiggle}[1]{%
  \tikz[baseline=-0.4ex]{
    \draw[#1,line width=2.6pt,line cap=round]
      plot [smooth] coordinates {(0,0.09) (0.34,0.01) (0.68,0.21) (1.02,0.01) (1.36,0.10)};
  }%
}
% Mot-clé surligné (marqueur orange sous le texte)
\newcommand{\cle}[1]{{\popxbold\color{popdark}#1}}

% ============================================================
% En-tête / pied
% ============================================================
\pagestyle{fancy}
\fancyhf{}
\renewcommand{\headrulewidth}{0pt}
\renewcommand{\footrulewidth}{0pt}
\lhead{\raisebox{-0.15ex}{\poplogo\fontsize{13}{13}\selectfont\color{popink}OnBuch\color{poporange}+}}
\rhead{\poppill{\DOCMATIERE\ \textbullet\ \DOCNIVEAU\ \textbullet\ Leçon \DOCLECON}{white}{popink}}
\lfoot{\popxbold\fontsize{7.6}{7.6}\selectfont\color{popmuted}\MakeUppercase{Cours signé OnBuch+}\ \textbullet\ {\popsemi\DOCTITRE}}
\rfoot{\tikz[baseline=-0.6ex]\node[rounded corners=8.5pt,fill=popink,minimum height=17pt,minimum width=17pt,inner xsep=5pt,inner sep=0]{\popxbold\fontsize{8}{8}\selectfont\color{popcream}\thepage\,/\,\pageref*{LastPage}};}

% ============================================================
% Titres
% ============================================================
\newcommand{\chaptitre}[2]{%
  \begin{tcolorbox}[enhanced,colback=poporange,colframe=popink,boxrule=2.6pt,arc=11pt,
      drop shadow=popink,left=15pt,right=15pt,top=12pt,bottom=11pt,before skip=3pt,after skip=18pt]
    \poppill{#2}{popink}{popcream}\par\vspace{9pt}
    {\raggedright\hyphenpenalty=10000\exhyphenpenalty=10000\popdisplay\fontsize{22}{24}\selectfont\color{popcream}\MakeUppercase{#1}\par}
  \end{tcolorbox}
}
% Section numérotée : pastille orange avec le numéro + titre Archivo
\newcounter{popsec}
\newcommand{\coursec}[1]{%
  \stepcounter{popsec}\setcounter{popsub}{0}%
  \par\vspace{16pt}\noindent
  \begin{minipage}{\linewidth}
  \tikz[baseline=-0.7ex]\node[draw=popink,line width=1.6pt,fill=poporange,rounded corners=4pt,minimum size=22pt,inner sep=0]{\popdisplay\fontsize{12}{12}\selectfont\color{popcream}\thepopsec};%
  \hspace{8pt}{\popdisplay\fontsize{15}{17}\selectfont\color{popink}\MakeUppercase{#1}}\par
  \vspace{4pt}\noindent{\color{poporange}\rule{36pt}{3.6pt}}
  \end{minipage}\par\nopagebreak\vspace{8pt}
}
\newcounter{popsub}
\newcommand{\courssub}[1]{%
  \stepcounter{popsub}%
  \par\vspace{9pt}\noindent{\popxbold\fontsize{11.5}{13}\selectfont\color{popink}{\color{poporange}\thepopsec.\thepopsub}\hspace{6pt}#1}\par\nopagebreak\vspace{4pt}
}

% ============================================================
% Boîtes pédagogiques — même recette : fond blanc, bordure épaisse colorée,
% ombre dure encre, badge plein. Titre optionnel [#1].
% ============================================================
\tcbset{popbase/.style={breakable,enhanced,colback=white,boxrule=2.3pt,arc=9pt,
  shadow={3pt}{-3pt}{0pt}{popink},left=14pt,right=14pt,top=11pt,bottom=11pt,before skip=11pt,after skip=15pt,
  fontupper=\popsemi\fontsize{10.6}{14.5}\selectfont\color{popink},
  fontlower=\popsemi\fontsize{10.6}{14.5}\selectfont\color{popink}}}
% Coche dessinée (le glyphe ✓ n'existe pas dans les polices du document)
\newcommand{\popcheck}[1]{\tikz[baseline=-0.5ex]\draw[#1,line width=1.6pt,line cap=round,line join=round](-0.13,0.0)--(-0.04,-0.09)--(0.13,0.1);}
\providecommand{\coche}{\popcheck{popgreen}}
\providecommand{\resultat}[1]{\par\smallskip\noindent\fcolorbox{popgreen}{popgreenL}{\parbox{\dimexpr\linewidth-2\fboxsep-2\fboxrule\relax}{\textbf{Résultat.} #1}}\par\smallskip}
\newcommand{\popboxhead}[3]{\poppill{#1}{#2}{white}\ifx\relax#3\relax\else\hspace{6pt}{\popxbold\fontsize{10.6}{12}\selectfont\color{popink}#3}\fi\par\vspace{7pt}}

\newtcolorbox{aretenir}[1][]{popbase,colframe=popgreen,before upper={\popboxhead{À retenir}{popgreen}{#1}}}
\newtcolorbox{exemplebox}[1][]{popbase,colframe=poporange,before upper={\popboxhead{Exemple}{poporange}{#1}}}
\newtcolorbox{definition}[1][]{popbase,colframe=popblue,before upper={\popboxhead{Définition}{popblue}{#1}}}
\newtcolorbox{propriete}[1][]{popbase,colframe=poppurple,before upper={\popboxhead{Propriété}{poppurple}{#1}}}
\newtcolorbox{methode}[1][]{popbase,colframe=popgold,before upper={\popboxhead{Méthode}{popgold}{#1}}}
\newtcolorbox{attention}[1][]{popbase,colframe=poppink,before upper={\popboxhead{Attention}{poppink}{#1}}}
\newtcolorbox{experience}[1][]{popbase,colframe=popblue,colback=popblueL,before upper={\popboxhead{Expérience}{popblue}{#1}}}
% « Le savais-tu ? » : carte violette pleine (contexte, culture, Cameroun)
\newtcolorbox{savaistu}[1][]{popbase,colback=poppurple,colframe=popink,
  fontupper=\popsemi\fontsize{10.4}{14.2}\selectfont\color{white},
  before upper={\poppill{Le savais-tu\,?}{white}{poppurple}\ifx\relax#1\relax\else\hspace{6pt}{\popxbold\color{white}#1}\fi\par\vspace{7pt}}}
% Exercice résolu : énoncé en haut, \tcblower puis la solution sur fond crème
\newcounter{popexo}\newcounter{popexr}
\newtcolorbox{exoresolu}[1][]{popbase,colframe=poporange,colbacklower=popcream,
  segmentation style={popink,line width=1.2pt,dashed},
  before upper={\stepcounter{popexr}\popboxhead{Exercice résolu \thepopexr}{poporange}{#1}},
  before lower={\poppill{Solution}{popink}{popcream}\par\vspace{6pt}}}
% Exercice (non résolu en place) — la correction est dans la partie Corrigés
\newtcolorbox{exercice}[1][]{popbase,colframe=popink,
  before upper={\stepcounter{popexo}\popboxhead{Exercice \thepopexo}{popink}{#1}}}
% Corrigé
\newtcolorbox{corrige}[1][]{popbase,colframe=popgreen,colback=popgreenL,
  before upper={\popboxhead{Corrigé}{popgreen}{#1}}}
% Objectifs / prérequis / bilan
\newtcolorbox{objectifs}{popbase,colframe=popink,colback=poporangeL,
  before upper={\popboxhead{Objectifs de la leçon}{popink}{}}}
\newtcolorbox{prerequis}{popbase,colframe=popink,colback=popblueL,
  before upper={\popboxhead{Prérequis}{popblue}{}}}
\newtcolorbox{bilan}{popbase,colframe=popink,colback=white,boxrule=2.8pt,
  before upper={\popboxhead{Fiche bilan}{poporange}{L'essentiel à connaître pour l'examen}}}
% Figure encadrée avec légende
\newtcolorbox{popfigure}[1][]{enhanced,breakable=false,colback=white,colframe=popline,boxrule=1.4pt,arc=8pt,
  left=10pt,right=10pt,top=10pt,bottom=8pt,before skip=10pt,after skip=12pt,halign=center,
  lower separated=false,halign lower=center,
  fontlower=\popsemi\fontsize{9}{11.5}\selectfont\color{popmuted},
  #1}
\newcommand{\legende}[1]{\par\vspace{6pt}{\popsemi\fontsize{9}{11.5}\selectfont\color{popmuted}#1}}

% ============================================================
% Signature OnBuch+
% ============================================================
\newcommand{\poplogoinline}[1]{{\poplogo\fontsize{#1}{#1}\selectfont\color{popink}OnBuch\color{poporange}+}}
\newcommand{\signatureonbuch}{%
  \begin{tcolorbox}[enhanced,colback=popink,colframe=popink,boxrule=2.2pt,arc=10pt,
      drop shadow=poporange,left=14pt,right=14pt,top=10pt,bottom=10pt,before skip=0pt,after skip=0pt]
    \tikz[baseline=-0.7ex]\node[circle,fill=popgreen,draw=popcream,line width=1.3pt,minimum size=22pt,inner sep=0]{\popcheck{white}};%
    \hspace{9pt}\begin{minipage}[c]{0.62\linewidth}
      {\popxbold\fontsize{12}{13}\selectfont\color{popcream}Signé\ \textbullet\ {\poplogo OnBuch\color{poporange}+}}\par\vspace{2pt}
      {\raggedright\popsemi\fontsize{8}{10.5}\selectfont\color{popline}\MakeUppercase{Équipe pédagogique OnBuch+}\\\MakeUppercase{Conforme au programme officiel MINESEC}\par}
    \end{minipage}\hfill
    {\poplogo\fontsize{22}{22}\selectfont\color{popcream}OnBuch\color{poporange}+}
  \end{tcolorbox}%
}

% ============================================================
% Page de garde
% #1 titre · #2 matière · #3 niveau · #4 module · #5 n° leçon · #6 durée ·
% #7 description / objectifs (texte ou itemize)
% ============================================================
\newcommand{\pagegarde}[7]{%
  \thispagestyle{empty}
  \newgeometry{a4paper,margin=1.7cm,top=1.6cm,bottom=1.4cm}%
  \begin{tikzpicture}[remember picture,overlay]
    \fill[poporange!18] ([xshift=-3.2cm,yshift=-3.6cm]current page.north east) circle (4.6cm);
    \fill[poppurple!14] ([xshift=2.4cm,yshift=7.6cm]current page.south west) circle (3.8cm);
  \end{tikzpicture}%
  {\poplogo\fontsize{30}{30}\selectfont\color{popink}OnBuch\color{poporange}+}\hfill
  \raisebox{0.6ex}{\poppill{Cours signé \textbullet\ Premium}{popink}{popcream}}\par
  \vspace{1pt}\popsquiggle{poppurple}\par
  \vspace{8pt}
  \begin{tcolorbox}[enhanced,colback=poporange,colframe=popink,boxrule=2.8pt,arc=13pt,
      drop shadow=popink,left=18pt,right=18pt,top=12pt,bottom=13pt,after skip=10pt]
    \poppill{#2 \textbullet\ #3}{popink}{popcream}\hspace{5pt}\poppill{#4}{white}{popink}\par\vspace{8pt}
    {\popdisplay\fontsize{14}{14}\selectfont\color{popink}LEÇON #5}\par\vspace{4pt}
    {\raggedright\hyphenpenalty=10000\exhyphenpenalty=10000\popdisplay\fontsize{25}{27}\selectfont\color{popcream}\MakeUppercase{#1}\par}
    \vspace{8pt}
    {\popxbold\fontsize{9.5}{9.5}\selectfont\color{popink}\MakeUppercase{Durée conseillée : #6}}
  \end{tcolorbox}
  \begin{tcolorbox}[enhanced,colback=white,colframe=popink,boxrule=2.2pt,arc=10pt,
      drop shadow=popink,left=16pt,right=16pt,top=10pt,bottom=10pt,after skip=8pt]
    \poppill{Dans cette leçon}{poppurple}{white}\par\vspace{5pt}
    \popsemi\fontsize{9.3}{12.6}\selectfont\color{popink}#7
  \end{tcolorbox}
  \noindent\begin{minipage}[t]{0.487\linewidth}
    \begin{tcolorbox}[enhanced,colback=popgreen,colframe=popink,boxrule=2.2pt,arc=10pt,
        drop shadow=popink,left=11pt,right=11pt,top=10pt,bottom=10pt,height=3.1cm,valign=center]
      \tcbox[colback=white,colframe=popink,boxrule=1.3pt,arc=5pt,boxsep=0pt,nobeforeafter,
        left=3pt,right=3pt,top=3pt,bottom=3pt,valign=center]{\qrcode[height=1.9cm]{https://play.google.com/store/apps/details?id=com.luvvix.onbuch&hl=fr}}%
      \hspace{8pt}\begin{minipage}[c]{0.5\linewidth}
        {\popdisplay\fontsize{11}{12.5}\selectfont\color{white}\MakeUppercase{Télécharge OnBuch}}\par\vspace{4pt}
        {\raggedright\popsemi\fontsize{8.4}{11}\selectfont\color{white}Gratuit \textbullet\ Google Play\\Tuteur IA, annales, concours, résultats\par}
      \end{minipage}
    \end{tcolorbox}
  \end{minipage}\hfill
  \begin{minipage}[t]{0.487\linewidth}
    \begin{tcolorbox}[enhanced,colback=poppurple,colframe=popink,boxrule=2.2pt,arc=10pt,
        drop shadow=popink,left=11pt,right=11pt,top=10pt,bottom=10pt,height=3.1cm,valign=center]
      \tikz[baseline=-0.7ex]\node[circle,fill=white,draw=popink,line width=1.3pt,minimum size=1.6cm,inner sep=0]{\popdisplay\fontsize{24}{24}\selectfont\color{poppurple}+};%
      \hspace{8pt}\begin{minipage}[c]{0.6\linewidth}
        {\popdisplay\fontsize{11}{12.5}\selectfont\color{white}\MakeUppercase{Passe à OnBuch+}}\par\vspace{4pt}
        {\raggedright\popsemi\fontsize{8.4}{11}\selectfont\color{white}Tous les cours signés, Tuteur IA illimité, fiches PDF — onglet OnBuch+ de l'app\par}
      \end{minipage}
    \end{tcolorbox}
  \end{minipage}
  \vspace{8pt}
  \popcontenu
  \vfill
  \signatureonbuch
  \restoregeometry
  \newpage
}

% Rangée « Ce cours contient » (page de garde)
\newcommand{\poptile}[3]{% #1 couleur #2 gros texte #3 légende
  \begin{tcolorbox}[enhanced,colback=#1,colframe=popink,boxrule=2pt,arc=9pt,shadow={2.5pt}{-2.5pt}{0pt}{popink},
      left=6pt,right=6pt,top=9pt,bottom=9pt,halign=center,valign=center,height=1.75cm,width=\linewidth,nobeforeafter]
    {\popdisplay\fontsize{12.5}{13}\selectfont\color{white}\MakeUppercase{#2}}\par\vspace{3pt}
    {\centering\popsemi\fontsize{7.8}{9.5}\selectfont\color{white}#3\par}
  \end{tcolorbox}}
\newcommand{\popcontenu}{%
  \par\noindent{\popxboldls\fontsize{7.5}{7.5}\selectfont\color{popmuted}CE COURS SIGNÉ CONTIENT}\par\vspace{7pt}
  \noindent\begin{minipage}[t]{0.235\linewidth}\poptile{popblue}{Cours}{complet, illustré, pas à pas}\end{minipage}\hfill
  \begin{minipage}[t]{0.235\linewidth}\poptile{poporange}{Méthodes}{et exercices résolus}\end{minipage}\hfill
  \begin{minipage}[t]{0.235\linewidth}\poptile{poppink}{Exercices}{type examen + corrigés}\end{minipage}\hfill
  \begin{minipage}[t]{0.235\linewidth}\poptile{popgreen}{Fiche bilan}{l'essentiel à retenir}\end{minipage}\par}

% Sommaire
\newtcolorbox{sommaire}{enhanced,colback=white,colframe=popink,boxrule=2.2pt,arc=10pt,
  drop shadow=popink,left=18pt,right=18pt,top=13pt,bottom=13pt,before skip=6pt,after skip=20pt,
  before upper={\poppill{Sommaire}{poppurple}{white}\par\vspace{9pt}\popsemi\fontsize{10.6}{14.5}\selectfont\color{popink}}}

% Signature de fin de cours — poussée tout en bas de la dernière page
\newcommand{\findecours}{%
  \par\vfill\nopagebreak
  \begin{center}\popsquiggle{poporange}\end{center}\vspace{4pt}
  \signatureonbuch
}
\AtBeginDocument{\def\llbracket{\mathopen{[\mkern-3.5mu[}}\def\rrbracket{\mathclose{]\mkern-3.5mu]}}} % intervalles d'entiers (glyphe ⟦ absent de la police)
% Glyphes absents de Fira Math : redéfinis à partir de glyphes disponibles
\AtBeginDocument{%
  \def\setminus{\mathbin{\backslash}}\let\smallsetminus\setminus
  \def\vdots{\mathord{\vbox{\baselineskip3.2pt\lineskiplimit0pt\kern4pt\hbox{.}\hbox{.}\hbox{.}}}}%
  \def\ddots{\mathinner{\mkern1mu\raise6.5pt\hbox{.}\mkern2mu\raise3.6pt\hbox{.}\mkern2mu\raise0.7pt\hbox{.}\mkern1mu}}%
  \def\triangle{\mathord{\scalebox{0.9}{$\symup{\Delta}$}}}%
  \def\nabla{\mathord{\raisebox{\depth}{\scalebox{1}[-1]{$\symup{\Delta}$}}}}%
  \def\checkmark{\popcheck{popdark}}%
  \def\star{\mathbin{\ast}}\def\bigstar{\mathord{\ast}}%
  \def\diamond{\mathbin{\scalebox{0.6}{$\Diamond$}}}%
  \def\bigcup{\mathop{\mathchoice{\raisebox{-0.3em}{\scalebox{1.7}{$\cup$}}}{\scalebox{1.25}{$\cup$}}{\cup}{\cup}}\displaylimits}%
  \def\bigcap{\mathop{\mathchoice{\raisebox{-0.3em}{\scalebox{1.7}{$\cap$}}}{\scalebox{1.25}{$\cap$}}{\cap}{\cap}}\displaylimits}%
  \def\lesseqqgtr{\mathrel{\lessgtr}}\def\bigoplus{\mathop{\oplus}\displaylimits}%
}
\providecommand{\ssi}{si et seulement si} % abréviation fréquente
\AtBeginDocument{\providecommand{\wideparen}{\overparen}} % arc de cercle

%% FIGFIX-BEGIN v4 — correctifs globaux des figures (docs/onbuch-generation/tools/figfix)
\usepackage{adjustbox}\usepackage{etoolbox}
% toute figure TikZ/pgfplots plus large que la ligne est réduite à la largeur disponible
\BeforeBeginEnvironment{tikzpicture}{\begin{adjustbox}{max width=\linewidth}}
\AfterEndEnvironment{tikzpicture}{\end{adjustbox}}
% légendes pgfplots lisibles par-dessus les courbes
\pgfplotsset{every axis/.append style={legend style={fill=white,fill opacity=0.92,text opacity=1,draw=black!15,inner sep=3pt}}}
% étiquettes d'arêtes (syntaxe edge["..."]) avec fond blanc
\tikzset{every edge quotes/.append style={fill=white,inner sep=1.2pt}}
%% FIGFIX-END
\begin{document}

\pagegarde{Reconstitution chronologique de l'histoire des fonds océaniques}{SVT}{1ère TI}{Module 3 : L'Éducation à l'Environnement et au Développement Durable}{NCT10}{5 h}{Cette leçon montre comment l'âge des fonds océaniques et les anomalies magnétiques permettent de reconstituer la chronologie de l'expansion océanique à partir des dorsales. Elle introduit la notion de tectonique des plaques en reliant l'expansion des océans au mouvement des lithosphères.
\vspace{4pt}
\begin{multicols}{2}\small
\begin{itemize}
\item À la fin de cette leçon, je sais décrire la structure d'une dorsale océanique et le mécanisme d'expansion des fonds océaniques.
\item À la fin de cette leçon, je sais exploiter une carte des âges des fonds océaniques pour reconstituer une chronologie d'ouverture d'un océan.
\item À la fin de cette leçon, je sais interpréter des anomalies magnétiques symétriques de part et d'autre d'une dorsale.
\item À la fin de cette leçon, je sais calculer la vitesse d'expansion d'un océan à partir d'âges et de distances.
\item À la fin de cette leçon, je sais relier l'expansion océanique au mouvement des plaques lithosphériques.
\item À la fin de cette leçon, je sais expliquer pourquoi les fonds océaniques sont jeunes et pourquoi leur âge augmente en s'éloignant de la dorsale.
\end{itemize}\end{multicols}}

\begin{sommaire}
\begin{multicols}{2}
\begin{enumerate}
\item Les dorsales et l'expansion des fonds océaniques
\item L'âge des fonds océaniques : un marqueur chronologique
\item Les anomalies magnétiques : une bande enregistreuse
\item Reconstitution chronologique : synthèse des marqueurs
\item De l'expansion océanique à la tectonique des plaques
\item Activité d'intégration
\item Méthodes et exercices résolus
\item Exercices
\item Corrigés des exercices
\item Fiche bilan
\end{enumerate}
\end{multicols}
\end{sommaire}

\begin{objectifs}
\begin{itemize}
\item À la fin de cette leçon, je sais décrire la structure d'une dorsale océanique et le mécanisme d'expansion des fonds océaniques.
\item À la fin de cette leçon, je sais exploiter une carte des âges des fonds océaniques pour reconstituer une chronologie d'ouverture d'un océan.
\item À la fin de cette leçon, je sais interpréter des anomalies magnétiques symétriques de part et d'autre d'une dorsale.
\item À la fin de cette leçon, je sais calculer la vitesse d'expansion d'un océan à partir d'âges et de distances.
\item À la fin de cette leçon, je sais relier l'expansion océanique au mouvement des plaques lithosphériques.
\item À la fin de cette leçon, je sais expliquer pourquoi les fonds océaniques sont jeunes et pourquoi leur âge augmente en s'éloignant de la dorsale.
\item À la fin de cette leçon, je sais construire et exploiter un profil magnétique simplifié perpendiculaire à une dorsale.
\end{itemize}
\end{objectifs}

\begin{prerequis}
\begin{itemize}
\item Structure interne du globe : croûte, manteau, lithosphère, asthénosphère (notions de Seconde).
\item Roches magmatiques : basalte et gabbro, roches des fonds océaniques.
\item Aimantation des minéraux ferromagnétiques et notion de champ magnétique terrestre.
\item Datation absolue des roches (principe de la radiochronologie, vu en début d'année).
\item Lecture de cartes et calculs de proportionnalité (échelles, vitesses).
\end{itemize}
\end{prerequis}

\newpage

\coursec{Les dorsales et l'expansion des fonds océaniques}

En 1960, des navires scientifiques sillonnent l'océan Atlantique et mesurent le magnétisme des roches du fond océanique. Le résultat est stupéfiant : de part et d'autre d'une immense chaîne de montagnes sous-marines, les roches forment des bandes magnétiques parfaitement symétriques, comme les deux moitiés d'une partition musicale dépliée. Autre énigme : alors que la Terre a 4,6 milliards d'années et que les continents portent des roches de plus de 3 milliards d'années, aucun fond océanique n'a plus de 200 millions d'années. Et les mesures les plus précises montrent que l'Atlantique s'élargit chaque année de quelques centimètres, éloignant lentement l'Afrique de l'Amérique du Sud. Un géologue camerounais étudiant la marge atlantique au large de Kribi peut-il dater les fonds marins sans y prélever une seule roche ?

\textbf{Problématique :} comment se forment les fonds océaniques, et comment reconstituer leur histoire chronologique ? Pour y répondre, nous devons d'abord comprendre ce qu'est une dorsale océanique et découvrir le mécanisme fondamental qui s'y déroule : l'\cle{expansion des fonds océaniques}.

\courssub{Qu'est-ce qu'une dorsale océanique ?}

\textbf{Une observation qui intrigue.} Lorsque les premiers océanographes dressent des cartes précises de la profondeur des océans (bathymétrie) au milieu du XX\textsuperscript{e} siècle, ils s'attendent à trouver des plaines plates et monotones. Surprise : au milieu de l'Atlantique s'élève une gigantesque chaîne de montagnes sous-marines, longue de plus de \SI{15000}{km}, large de 1000 à \SI{2000}{km}, dont les sommets culminent parfois à moins de \SI{2000}{m} sous la surface. Certains sommets émergent même : l'Islande, les Açores, l'île de l'Ascension sont des morceaux de cette chaîne qui dépassent du niveau de la mer. Au sommet de la chaîne, les sondeurs révèlent une dépression étroite et profonde, une véritable vallée encaissée qui court le long de l'axe : le \cle{rift} axial.

\begin{definition}[La dorsale océanique]
Une \cle{dorsale océanique} est une chaîne de montagnes sous-marines au relief accidenté, dont l'axe est occupé par une dépression appelée \cle{rift} central (ou rift axial). C'est au niveau de ce rift que se forme en permanence de la nouvelle croûte océanique, par cristallisation du magma remonté de l'asthénosphère.
\end{definition}

\begin{attention}[Ne pas confondre]
Une dorsale océanique n'est \textbf{pas} un plissement comme la chaîne des Monts Oku ou les monts Bamboutos à l'Ouest camerounais. Les montagnes continentales résultent de la compression et du plissement de roches anciennes ; la dorsale, elle, est un relief \emph{d'origine volcanique}, constitué de roches \emph{jeunes} qui viennent de se former. C'est une zone d'\textbf{ouverture}, pas de collision.
\end{attention}

\courssub{Le réseau mondial des dorsales}

Les dorsales ne sont pas un accident isolé : elles forment un réseau continu d'environ \SI{60000}{km} de long, la plus grande chaîne de montagnes de la planète, presque entièrement cachée sous les océans. On distingue trois dorsales majeures :

\begin{itemize}
\item la \cle{dorsale médio-atlantique}, qui parcourt l'océan Atlantique du nord au sud, en épousant remarquablement le tracé des côtes africaines et sud-américaines — elle passe donc dans le prolongement du golfe de Guinée, non loin des côtes camerounaises ;
\item la \cle{dorsale est-Pacifique}, qui remonte le long de la côte ouest de l'Amérique du Sud ; son expansion y est rapide et son relief plus adouci ;
\item la \cle{dorsale médio-indienne}, qui sépare l'Afrique de l'Australie et se prolonge vers la mer Rouge et le golfe d'Aden.
\end{itemize}

\begin{popfigure}
\begin{center}
\begin{tikzpicture}[scale=1]
% Océan de fond
\fill[popblueL] (-7.5,-2.6) rectangle (7.5,2.6);
% Continents schématiques
\fill[popgreenL,draw=popgreen!60!black,thick] (-7.2,2.3) .. controls (-6.2,2.6) and (-5.4,1.8) .. (-5.0,1.2) .. controls (-4.6,0.6) and (-4.4,0.2) .. (-4.2,-0.4) .. controls (-4.0,-1.2) and (-4.6,-2.0) .. (-5.2,-2.4) .. controls (-6.0,-2.5) and (-6.8,-1.5) .. (-7.2,-0.5) .. controls (-7.5,0.6) and (-7.5,1.6) .. (-7.2,2.3) -- cycle; % Amériques simplifiées
\fill[popgreenL,draw=popgreen!60!black,thick] (-1.6,2.4) .. controls (-0.6,2.5) and (0.2,1.8) .. (0.4,1.0) .. controls (0.6,0.2) and (0.2,-0.6) .. (-0.2,-1.4) .. controls (-0.6,-2.2) and (-1.2,-2.4) .. (-1.4,-1.6) .. controls (-1.7,-0.6) and (-2.0,0.6) .. (-1.9,1.4) .. controls (-1.85,2.0) and (-1.75,2.2) .. (-1.6,2.4) -- cycle; % Afrique
\fill[popgreenL,draw=popgreen!60!black,thick] (1.2,2.4) .. controls (2.6,2.5) and (4.2,2.0) .. (5.0,1.2) .. controls (5.6,0.6) and (5.4,-0.2) .. (4.8,-0.8) .. controls (4.0,-1.5) and (2.8,-1.2) .. (2.0,-0.6) .. controls (1.2,0.0) and (0.9,1.2) .. (1.2,2.4) -- cycle; % Eurasie
\fill[popgreenL,draw=popgreen!60!black,thick] (4.4,-1.6) .. controls (5.2,-1.2) and (6.0,-1.4) .. (6.4,-1.9) .. controls (6.6,-2.4) and (5.8,-2.5) .. (5.0,-2.4) .. controls (4.4,-2.3) and (4.2,-2.0) .. (4.4,-1.6) -- cycle; % Australie
% Dorsale médio-atlantique
\draw[poporange,very thick,decorate,decoration={snake,amplitude=1.2pt,segment length=6pt}] (-2.9,2.5) .. controls (-3.3,1.6) and (-2.7,0.8) .. (-3.0,0.0) .. controls (-3.3,-0.9) and (-2.8,-1.8) .. (-3.1,-2.5);
% Dorsale est-Pacifique
\draw[poporange,very thick,decorate,decoration={snake,amplitude=1.2pt,segment length=6pt}] (-6.6,1.8) .. controls (-6.2,0.8) and (-6.5,-0.4) .. (-6.1,-1.6) .. controls (-5.9,-2.1) and (-6.2,-2.4) .. (-6.3,-2.5);
% Dorsale médio-indienne
\draw[poporange,very thick,decorate,decoration={snake,amplitude=1.2pt,segment length=6pt}] (2.2,0.6) .. controls (2.8,-0.2) and (2.4,-1.2) .. (3.2,-1.8) .. controls (3.8,-2.2) and (4.6,-2.3) .. (5.2,-2.4);
% Légendes pointées
\draw[popdark,thin] (-3.0,0.0) -- (-1.2,0.6) node[right,font=\small\itshape]{Dorsale médio-atlantique};
\draw[popdark,thin] (-6.3,-0.6) -- (-4.9,-1.9) node[right,font=\small\itshape]{Dorsale est-Pacifique};
\draw[popdark,thin] (2.8,-1.0) -- (1.6,-2.2) node[left,font=\small\itshape]{Dorsale médio-indienne};
\node[font=\small\itshape,popdark] at (-0.7,1.5) {Afrique};
\node[font=\small\itshape,popdark] at (-5.6,1.0) {Amériques};
\node[font=\small\itshape,popdark] at (3.2,1.6) {Eurasie};
\node[font=\small\itshape,popdark] at (5.4,-2.0) {Australie};
\end{tikzpicture}
\end{center}
\legende{Carte simplifiée du monde montrant le réseau des dorsales océaniques (tracés orange ondulés) : dorsale médio-atlantique au milieu de l'Atlantique, dorsale est-Pacifique le long de l'Amérique du Sud, dorsale médio-indienne entre l'Afrique et l'Australie.}
\end{popfigure}

\begin{attention}[Un détail qui compte]
Le nom d'une dorsale indique sa position \emph{dans l'océan}, pas sur le globe : la dorsale médio-atlantique est bien au milieu de l'Atlantique, mais la dorsale est-Pacifique n'est \emph{pas} au milieu du Pacifique — elle est décalée vers l'est. Au Probatoire, situe chaque dorsale par rapport à son océan.
\end{attention}

\courssub{La mise en évidence de l'expansion océanique}

\textbf{La démarche d'investigation.} Suivons le raisonnement des océanographes des années 1960, étape par étape.

\begin{enumerate}
\item \textbf{Observation.} Les campagnes de forages du programme \emph{Deep Sea Drilling Project} (DSDP, à partir de 1968, navire \emph{Glomar Challenger}) prélèvent des carottes de sédiments et de basaltes à différentes distances de l'axe de la dorsale médio-atlantique. Résultat : juste au niveau du rift, les basaltes sont pratiquement dépourvus de sédiments et leur âge est quasi nul ; plus on s'éloigne de l'axe, plus la couverture sédimentaire s'épaissit et plus les basaltes sous-jacents sont vieux.
\item \textbf{Hypothèse.} Le plancher océanique se formerait au niveau de la dorsale, puis s'en écarterait progressivement, comme deux tapis roulants partant du rift vers chaque continent.
\item \textbf{Vérification par les données.} Les âges mesurés de part et d'autre de l'axe sont \emph{symétriques} : à \SI{500}{km} à l'ouest et à \SI{500}{km} à l'est de la dorsale, on trouve des basaltes du même âge. Et l'âge croît régulièrement avec la distance.
\item \textbf{Interprétation.} Si le plancher s'écarte à vitesse constante, alors l'âge d'une roche est proportionnel à sa distance à l'axe : c'est exactement ce que montrent les forages.
\item \textbf{Conclusion.} Les fonds océaniques s'écartent de la dorsale de part et d'autre : c'est l'\cle{expansion océanique} (ou expansion des fonds océaniques).
\end{enumerate}

\begin{propriete}[Âge du plancher océanique et distance à la dorsale]
L'âge du plancher océanique (basaltes de la croûte et premiers sédiments déposés dessus) \textbf{croît régulièrement et symétriquement avec la distance à l'axe de la dorsale}. Cette propriété, admise, repose sur les résultats des campagnes de forages du Deep Sea Drilling Project. Conséquence directe : la bande de croûte la plus jeune (âge nul, 0 Ma) se trouve toujours au niveau du rift axial.
\end{propriete}

\begin{methode}[Repérer une dorsale sur une carte des âges des fonds]
Face à une carte des âges océaniques au Probatoire, procède ainsi :
\begin{enumerate}
\item Repère la ligne des fonds les plus jeunes, d'âge 0 Ma (souvent colorée en rouge ou en jaune vif).
\item Vérifie que cette ligne est bordée de part et d'autre par des bandes d'âges \emph{croissants} et \emph{symétriques} (mêmes âges à égale distance).
\item Conclus : cette ligne des âges nuls est l'axe de la dorsale ; les bandes symétriques traduisent l'expansion océanique.
\item Si deux bandes d'âges identiques se font face, mesure la distance qui les sépare de l'axe : elle doit être égale des deux côtés.
\end{enumerate}
\end{methode}

\courssub{Le mécanisme de l'expansion : fabriquer du plancher océanique}

Que se passe-t-il concrètement sous le rift ? Sous la dorsale, la pression exercée sur les roches profondes du manteau diminue localement (décompression), car le plancher s'écarte et « soulage » le poids. Or, à température élevée, une baisse de pression suffit à faire fondre partiellement les roches du manteau supérieur : c'est la \cle{fusion partielle} de la péridotite asthénosphérique. Le magma basaltique ainsi produit, moins dense que les roches environnantes, \textbf{remonte} vers la surface par les fissures du rift.

Arrivé au contact de l'eau de mer froide (environ 2 à \SI{4}{\celsius} au fond des océans), le magma se refroidit brutalement et \textbf{cristallise} :
\begin{itemize}
\item en surface, il forme des \cle{basaltes}, souvent en forme de coussins empilés appelés \emph{pillow-lavas} (laves en oreillers), témoins d'un refroidissement brutal sous l'eau ;
\item en profondeur, la cristallisation lente donne des \cle{gabbros}, roches de même composition chimique que les basaltes mais à gros cristaux.
\end{itemize}

Cet empilement basaltes-gabbros constitue la \cle{croûte océanique} nouvellement formée, d'une épaisseur moyenne d'environ 6 à \SI{7}{km}. Mais l'apport de magma ne s'arrête jamais : chaque nouvelle injection de magma au niveau du rift \textbf{pousse} la croûte déjà formée de part et d'autre, comme une fermeture éclair que l'on ouvrirait en insérant continuellement de nouvelles dents au milieu. Le plancher océanique s'écarte ainsi de l'axe à une vitesse quasi constante, de l'ordre de 1 à \SI{10}{cm/an} selon les dorsales.

\begin{popfigure}
\begin{center}
\begin{tikzpicture}[scale=1.05]
% Asthénosphère
\fill[poporangeL] (-6.4,-2.6) rectangle (6.4,-0.9);
\node[popdark,font=\small\itshape] at (0,-2.35) {Asthénosphère (manteau, roches ductiles)};
% Lithosphère / croûte océanique : blocs symétriques
\fill[popblue!35] (-6.4,-0.9) rectangle (-0.55,0.1);
\fill[popblue!35] (0.55,-0.9) rectangle (6.4,0.1);
% Strates d'âge (bandes de plus en plus vieilles)
\foreach \x/\c in {0.9/popblue!55, 1.9/popblue!45, 3.1/popblue!35, 4.5/popblue!25}{
  \fill[\c] (\x,-0.9) rectangle (\x+0.8,0.1);
  \fill[\c] (-\x-0.8,-0.9) rectangle (-\x,0.1);}
% Rift : dépression axiale
\fill[popblue!10] (-0.55,-0.55) -- (0.55,-0.55) -- (0.55,0.1) -- (-0.55,0.1) -- cycle;
\draw[popdark,thick] (-0.55,0.1) -- (-0.55,-0.55) -- (0.55,-0.55) -- (0.55,0.1);
% Magma remontant
\fill[popred!70!poporange] (-0.28,-2.6) -- (-0.28,-0.9) -- (-0.55,-0.55) -- (0.55,-0.55) -- (0.28,-0.9) -- (0.28,-2.6) -- cycle;
\draw[popred!70!poporange,->,very thick] (0,-2.2) -- (0,-0.75);
% Flèches de divergence
\draw[popgreen!60!black,->,very thick] (-1.1,0.35) -- (-2.9,0.35);
\draw[popgreen!60!black,->,very thick] (1.1,0.35) -- (2.9,0.35);
\node[popgreen!50!black,font=\small\bfseries] at (-2.0,0.62) {$v$};
\node[popgreen!50!black,font=\small\bfseries] at (2.0,0.62) {$v$};
% Âges
\node[font=\scriptsize,popdark] at (0,0.28) {0 Ma};
\node[font=\scriptsize,popdark] at (1.3,0.28) {jeune};
\node[font=\scriptsize,popdark] at (-1.3,0.28) {jeune};
\node[font=\scriptsize,popdark] at (5.6,0.28) {plus âgé};
\node[font=\scriptsize,popdark] at (-5.6,0.28) {plus âgé};
% Légendes pointées
\draw[popdark,thin] (0,-0.35) -- (2.6,-0.5) node[right,font=\small\itshape]{Rift axial};
\draw[popdark,thin] (0.15,-1.6) -- (2.6,-1.7) node[right,font=\small\itshape]{Remontée de magma basaltique};
\draw[popdark,thin] (-4.0,-0.4) -- (-4.0,1.1) node[above,font=\small\itshape]{Croûte océanique (basaltes + gabbros)};
\draw[popdark,thin] (-2.0,0.35) -- (-2.0,1.5) node[above,font=\small\itshape]{Flanc ouest : s'écarte à la vitesse $v$};
\end{tikzpicture}
\end{center}
\legende{Coupe schématique d'une dorsale océanique. Le magma issu de la fusion partielle de l'asthénosphère remonte au niveau du rift axial et cristallise en basaltes et gabbros, formant une croûte océanique neuve (âge 0 Ma à l'axe). Les apports successifs poussent la croûte de part et d'autre : les deux flancs divergent à la même vitesse $v$, et l'âge des roches croît symétriquement en s'éloignant de l'axe.}
\end{popfigure}

\courssub{Taux d'expansion et vitesse d'ouverture : un vocabulaire précis}

Puisque la croûte s'écarte des \emph{deux} côtés de la dorsale, il faut distinguer rigoureusement deux grandeurs :

\begin{definition}[Taux d'expansion et vitesse d'ouverture]
Le \cle{taux d'expansion} (ou demi-vitesse d'expansion) est la vitesse d'écartement d'\textbf{un seul flanc} par rapport à l'axe de la dorsale ; il s'exprime en cm/an. La \cle{vitesse d'ouverture totale} de l'océan est la somme des écartements des deux flancs : pour une expansion symétrique, elle vaut le \textbf{double} du taux d'expansion.
\[ v_{\text{ouverture totale}} = v_{\text{flanc gauche}} + v_{\text{flanc droit}} = 2 \times v_{\text{taux}} \quad \text{(expansion symétrique)} \]
\end{definition}

\begin{exemplebox}[La dorsale médio-atlantique]
La dorsale médio-atlantique s'écarte d'environ \SI{2}{cm/an} \textbf{par flanc}. La vitesse d'ouverture totale de l'océan Atlantique est donc d'environ $2 + 2 = \SI{4}{cm/an}$. En un siècle, l'Atlantique gagne ainsi environ 4 mètres de largeur ; en 100 millions d'années, cela représente \SI{4000}{km} — l'ordre de grandeur de la largeur actuelle de l'Atlantique, ce qui confirme que cet océan est né de l'ouverture progressive entre l'Afrique et l'Amérique du Sud.
\end{exemplebox}

\begin{attention}[L'erreur classique du Probatoire]
Ne confonds jamais le \textbf{taux d'expansion} (un seul flanc, en cm/an) avec la \textbf{vitesse d'ouverture totale} (les deux flancs, donc le double). Lis attentivement l'énoncé : si l'on te donne la distance entre deux roches de même âge situées \emph{de part et d'autre} de l'axe, la vitesse calculée à partir de cette distance totale est une vitesse d'ouverture ; le taux d'expansion s'obtient en la divisant par deux. Inversement, si l'énoncé donne « le plancher s'écarte de la dorsale de 2 cm/an », il s'agit du taux par flanc.
\end{attention}

\begin{exoresolu}[Calculer un taux d'expansion à partir d'âges]
Énoncé. Des forages réalisés dans l'Atlantique Sud montrent qu'un basalte situé à \SI{900}{km} à l'ouest de l'axe de la dorsale médio-atlantique est âgé de 60 Ma, et qu'un basalte situé à \SI{900}{km} à l'est de l'axe a le même âge. (1) Calculer le taux d'expansion de chaque flanc. (2) En déduire la vitesse d'ouverture totale de l'Atlantique à cette latitude.
\tcblower
Solution détaillée.
\begin{enumerate}
\item Le basalte ouest s'est formé à l'axe il y a 60 Ma, puis a parcouru \SI{900}{km}. Le taux d'expansion du flanc ouest est :
\[ v = \dfrac{d}{t} = \dfrac{\SI{900}{km}}{\SI{60}{Ma}} = \SI{15}{km/Ma} \]
Or $\SI{15}{km/Ma} = \dfrac{15 \times 10^{5}\ \text{cm}}{10^{6}\ \text{ans}} = \SI{1,5}{cm/an}$. La symétrie des âges montre que le flanc est a le même taux : \SI{1,5}{cm/an}.
\item La vitesse d'ouverture totale est la somme des deux flancs :
\[ v_{\text{totale}} = 1{,}5 + 1{,}5 = \SI{3}{cm/an} \]
Conclusion : à cette latitude, l'Atlantique s'élargit d'environ \SI{3}{cm/an}, chaque flanc s'écartant de la dorsale de \SI{1,5}{cm/an}.
\end{enumerate}
\end{exoresolu}

\courssub{Une conséquence majeure : des océans jeunes sur une Terre vieille}

L'expansion océanique explique enfin l'énigme de départ. Si la croûte océanique se forme en permanence aux dorsales et s'en écarte, elle doit aussi \emph{disparaître} quelque part (nous verrons plus tard qu'elle plonge dans les zones de subduction) : le plancher océanique est donc \textbf{renouvelé en permanence}. C'est pourquoi :

\begin{aretenir}[Jeunesse des fonds océaniques]
Aucun fond océanique n'a plus d'environ \textbf{200 Ma} (0,2 Ga), alors que les continents portent des roches de plus de \textbf{3 Ga} (certains cratons africains, comme ceux du socle camerounais, dépassent 2,5 Ga). Les continents sont des archives anciennes et conservées ; les fonds océaniques sont un tapis roulant jeune, fabriqué aux dorsales puis recyclé. La carte des âges océaniques est donc une véritable \emph{chronologie} de l'histoire récente de la Terre.
\end{aretenir}

\begin{savaistu}[L'Atlantique grandit comme ton ongle]
L'Atlantique s'élargit d'environ \SI{4}{cm/an}, soit à peu près la vitesse de pousse de tes ongles. Cela semble dérisoire, mais à l'échelle des temps géologiques, c'est colossal : depuis l'époque des dinosaures (il y a environ 150 Ma), l'Atlantique a gagné plusieurs milliers de kilomètres de largeur. L'Afrique — et donc le Cameroun — s'éloigne de l'Amérique du Sud d'environ 4 mètres par siècle. Au large de Kribi, la marge continentale camerounaise porte d'ailleurs les traces de cette ouverture : c'est en étudiant les âges des fonds au large du golfe de Guinée que les géologues retracent la séparation de l'Afrique et de l'Amérique du Sud.
\end{savaistu}

\begin{exercice}[Pour t'entraîner]
Sur une carte des âges de l'océan Indien, on repère une bande de fonds âgée de 40 Ma située à \SI{800}{km} de part et d'autre de l'axe de la dorsale médio-indienne. (1) Comment appelle-t-on la structure située à l'axe ? (2) Calculer le taux d'expansion d'un flanc, en cm/an. (3) Calculer la vitesse d'ouverture totale de l'océan à cette latitude. (4) Expliquer pourquoi on ne trouve nulle part dans cet océan de fonds plus âgés que 200 Ma.
\end{exercice}

\begin{corrige}[Exercice : dorsale médio-indienne]
\begin{enumerate}
\item À l'axe se trouve une \textbf{dorsale océanique}, marquée par un rift central où se forme la croûte neuve (âge 0 Ma).
\item Taux d'expansion d'un flanc : $v = \dfrac{d}{t} = \dfrac{\SI{800}{km}}{\SI{40}{Ma}} = \SI{20}{km/Ma} = \SI{2}{cm/an}$.
\item Vitesse d'ouverture totale : $2 \times 2 = \SI{4}{cm/an}$ (somme des deux flancs).
\item La croûte océanique se forme en permanence à la dorsale, s'en écarte, puis est recyclée (elle disparaît dans les zones de subduction). Le plancher océanique est donc entièrement renouvelé en moins de 200 Ma, contrairement aux continents qui conservent des roches de plus de 3 Ga.
\end{enumerate}
\end{corrige}

En résumé, les dorsales sont des chaînes de montagnes sous-marines où le magma asthénosphérique cristallise en croûte océanique neuve ; poussée par les apports suivants, cette croûte s'écarte symétriquement de l'axe à vitesse quasi constante, ce qui explique la jeunesse des fonds océaniques et fournit un premier marqueur chronologique : l'âge. Mais comment dater précisément ces fonds sans forer partout ? C'est là qu'intervient un enregistreur naturel remarquable : le magnétisme des roches, que nous étudierons dans la suite de la leçon.

\coursec{L'âge des fonds océaniques : un marqueur chronologique}

\courssub{Transition depuis l'expansion des dorsales}
Dans la section précédente, nous avons mis en évidence que les dorsales médio‑océaniques sont des zones de création continue de croûte océanique par refroidissement de magma basaltique. Cette création permanente repousse la croûte déjà formée vers l'extérieur, de part et d'autre de l'axe. Pour transformer cette observation qualitative en une histoire datée et quantifiée, il faut pouvoir \emph{mesurer l'âge} des roches qui composent le fond océanique. C'est l'objet de cette section.

\courssub{Principes de datation des fonds océaniques}

\begin{experience}[Datation des basaltes et des sédiments de couverture (Programmes DSDP/ODP/IODP)]
\textbf{Matériel :} Échantillons de basalte prélevés par forage océanique (programmes internationaux DSDP, ODP, IODP) à différentes distances orthogonales de l'axe d'une dorsale ; carottes de sédiments recouvrant le basalte (sédiments pélagiques).\\
\textbf{Protocole :}
\begin{enumerate}
\item \textbf{Radiochronologie du basalte :} Mesurer les rapports isotopiques \ce{^{40}Ar/^{39}Ar} (méthode \ce{K/Ar} modernisée) ou \ce{^{238}U/^{206}Pb} dans le verre basalte frais (marge de verre du coussin) pour obtenir un âge radiométrique absolu (âge de cristallisation).
\item \textbf{Biostratigraphie des sédiments de base :} Identifier le premier sédiment déposé directement sur le basalte (niveau basal, souvent un argile rouge ou un calcaire nannofossilifère) et dater ses microfossiles (foraminifères planctoniques, nannofossiles calcaires) par biostratigraphie (échelles de zones biochronologiques).
\item \textbf{Profil perpendiculaire :} Répéter l'opération sur plusieurs sites alignés le long d'un profil perpendiculaire à l'axe de la dorsale (flanc Est et flanc Ouest).
\end{enumerate}
\textbf{Observations :} Les âges radiométriques des basaltes et les âges biostratigraphiques des sédiments de base coïncident dans les incertitudes de mesure. L'âge augmente régulièrement quand on s'éloigne de l'axe de la dorsale, de part et d'autre.
\textbf{Interprétation :} Le basalte se solidifie à l'axe de la dorsale (âge ≈ 0 Ma) ; les sédiments pélagiques commencent à s'accumuler immédiatement après la formation de la croûte. La croûte océanique s'éloigne de l'axe en conservant son âge de formation.
\textbf{Conclusion :} L'âge du fond océanique en un point donné est l'âge de sa formation à la dorsale. Ce principe fondamental permet de transformer une carte de distances (par rapport à l'axe) en une carte d'âges absolus.
\end{experience}

\begin{definition}[Isochrone]
Une \cle{isochrone} est une ligne imaginaire tracée sur une carte des fonds océaniques qui relie tous les points ayant le \textbf{même âge} de formation. Les isochrones sont \textbf{parallèles à l'axe de la dorsale} et leur espacement (distance entre deux isochrones successives) reflète la vitesse d'expansion (taux d'accrétion) sur la période considérée.
\end{definition}

\courssub{La carte des âges : code couleur, symétrie et constat fondamental}

\begin{popfigure}
\begin{tikzpicture}[scale=0.75]
% Dorsale (axe vertical)
\draw[popblue, line width=2.5pt] (0,-9) -- (0,9) node[midway, right=4pt, popblue, rotate=90, anchor=south] {\textbf{Dorsale médio-océanique (0 Ma)}};
% Échelle : 1 cm = 500 km. Taux demi-dorsale 2 cm/an = 20 km/Ma.
% Bandes de 20 Ma d'intervalle -> largeur 400 km = 0.8 cm sur figure.
\foreach \age/\col in {
  20/popgreenL, 40/popgreen, 60/poporangeL, 80/poporange, 
  100/poppinkL, 120/poppink, 140/poppurpleL, 160/poppurple, 180/popblueL}{
  \pgfmathsetmacro{\dist}{\age * 20 / 500} % distance en cm sur figure
  \pgfmathsetmacro{\prevdist}{(\age-20) * 20 / 500}
  % Flanc Est (x > 0)
  \draw[fill=\col, opacity=0.6, draw=popline] (\prevdist,-9) -- (\dist,-9) -- (\dist,9) -- (\prevdist,9) -- cycle;
  % Flanc Ouest (x < 0)
  \draw[fill=\col, opacity=0.6, draw=popline] (-\prevdist,-9) -- (-\dist,-9) -- (-\dist,9) -- (-\prevdist,9) -- cycle;
  % Légendes âges (côté droit)
  \node[popdark, font=\small\bfseries] at (\dist+0.3, 0) {\age~Ma};
  \node[popdark, font=\small\bfseries] at (-\dist-0.3, 0) {\age~Ma};
}
% Flèches d'expansion
\draw[<->, popdark, thick, >=Stealth] (-4.5,-9.5) -- (4.5,-9.5) node[midway, below=2pt] {Distance à la dorsale (km)};
\draw[->, poporange, thick, >=Stealth] (0,-10.5) -- (2.5,-10.5) node[midway, below=2pt] {Écartement Est};
\draw[->, poporange, thick, >=Stealth] (0,-10.5) -- (-2.5,-10.5) node[midway, below=2pt] {Écartement Ouest};
% Barre d'échelle
\draw[|-|, popdark, thick] (0,-11.5) -- (3.6,-11.5) node[midway, below=2pt] {1800 km (≈ 90 Ma à 2 cm/an)};
\end{tikzpicture}
\legende{Carte simplifiée en vue de plan (Atlantique Nord type) : isochrones (lignes limites des bandes colorées) symétriques par rapport à l'axe de la dorsale (ligne bleue centrale). Chaque bande colorée correspond à un intervalle de 20 Ma. L'égalité de largeur des bandes (pour un intervalle de temps constant) témoigne d'un taux d'expansion constant. Les âges indiqués (20, 40, ..., 180 Ma) sont les âges aux limites des bandes.}
\end{popfigure}

\begin{propriete}[Relation âge–distance et symétrie des isochrones]
Sur une carte des âges des fonds océaniques (ou un profil perpendiculaire à la dorsale) :
\begin{itemize}
\item L'âge est \textbf{nul à l'axe de la dorsale} (formation continue actuelle).
\item L'âge augmente de part et d'autre jusqu'à un \textbf{maximum de 180–200 Ma} (Jurassique supérieur) au contact des marges continentales passives (ex. : marges africaine et brésilienne pour l'Atlantique Sud) ou des fosses de subduction (ex. : Pacifique Ouest).
\item Les isochrones sont \textbf{rigoureusement symétriques} par rapport à l'axe : un point daté de 60 Ma se trouve à la même distance de la dorsale sur le flanc est et sur le flanc ouest.
\end{itemize}
Cette symétrie bilatérale est la preuve directe que les deux flancs s'écartent \textbf{régulièrement et à la même vitesse} (taux d'expansion demi-dorsale constant) depuis l'ouverture de l'océan.
\end{propriete}

\courssub{Reconstitution chronologique et calcul du taux d'expansion}

\begin{methode}[Reconstituer l'histoire d'ouverture d'un océan à partir d'une carte d'âges]
\begin{enumerate}
\item \textbf{Repérer l'axe de la dorsale} (isochrone 0 Ma, souvent matérialisée par une vallée axiale ou une anomalie magnétique centrale).
\item \textbf{Identifier les isochrones symétriques} de part et d'autre (ex. : 10, 20, 40, 60, 80 Ma...).
\item \textbf{Ordonner chronologiquement} les événements : du plus récent (dorsale, 0 Ma) au plus ancien (marges/fosses, ~180 Ma).
\item \textbf{Mesurer la distance orthogonale} $d$ (en km) entre l'axe et une isochrone d'âge connu $t$ (en Ma), sur un seul flanc.
\item \textbf{Calculer le taux d'expansion demi-dorsale} $v_{\text{demi}} = \dfrac{d}{t}$ (conversion impérative : $d$ en cm, $t$ en années $\rightarrow$ $v$ en cm/an).
\item \textbf{Déduire le taux d'expansion complet} (vitesse relative des deux plaques) : $v_{\text{complet}} = 2 \times v_{\text{demi}}$.
\end{enumerate}
\end{methode}

\begin{propriete}[Démonstration guidée de la relation $v = d/t$]
On suppose un taux d'expansion demi-dorsale constant $v$ (en cm/an) sur l'intervalle de temps considéré. Un fragment de croûte formé il y a $t$ années (âge $t$) à l'axe de la dorsale s'est déplacé à la vitesse $v$ pendant la durée $t$. La distance $d$ parcourue (distance actuelle à l'axe) est donc :
\[ d = v \times t \quad \Rightarrow \quad v = \frac{d}{t} \]
La linéarité de la relation âge–distance (points alignés sur un graphique $d = f(t)$) valide \emph{a posteriori} l'hypothèse de vitesse constante sur l'intervalle considéré. Si la courbe s'infléchit, le taux a varié.
\end{propriete}

\begin{exoresolu}[Calcul du taux d'expansion à partir de données de forage]
\textbf{Énoncé :} Un forage réalisé sur le flanc est de la dorsale médio‑atlantique (site ODP légendaire) a prélevé un basalte de socle daté par la méthode \ce{^{40}Ar/^{39}Ar} à \num{60} Ma. La distance orthogonale entre le site de forage et l'axe actuel de la dorsale est de \num{1200} km. Déterminer le taux d'expansion demi‑dorsale et le taux d'expansion complet.
\tcblower
\textbf{Solution détaillée :}
\begin{align*}
\text{Données :} \quad & d = \num{1200}\ \text{km} = \num{1,2e8}\ \text{cm} \quad (\text{car } 1\ \text{km} = 10^5\ \text{cm}) \\
& t = \num{60}\ \text{Ma} = \num{6,0e7}\ \text{an} \quad (\text{car } 1\ \text{Ma} = 10^6\ \text{an}) \\[0.5em]
\text{Taux demi-dorsale :} \quad & v_{\text{demi}} = \frac{d}{t} = \frac{\num{1,2e8}\ \text{cm}}{\num{6,0e7}\ \text{an}} = \num{2,0}\ \text{cm/an} \\[0.5em]
\text{Taux complet (séparation des deux plaques) :} \quad & v_{\text{complet}} = 2 \times v_{\text{demi}} = \num{4,0}\ \text{cm/an}
\end{align*}
\textbf{Conclusion :} Le taux d'expansion demi‑dorsale est de \textbf{2,0 cm/an}. Le taux d'expansion complet vaut \textbf{4,0 cm/an}, valeur caractéristique d'une dorsale à expansion lente (Atlantique Nord actuel).
\end{exoresolu}

\begin{attention}[Piège fréquent : distance totale vs distance à l'axe]
Lors du calcul de $v_{\text{demi}}$, \textbf{n'utilisez jamais la distance entre deux points symétriques} (ex. : \num{2400} km entre les deux points datés de 60 Ma de part et d'autre de la dorsale). Cette distance totale correspond à $2d$ et donnerait un taux deux fois trop grand (\num{4} cm/an au lieu de \num{2} cm/an). La formule $v = d/t$ utilise la distance \textbf{d'un seul flanc à l'axe} ($d$). Vérifiez toujours : « Distance mesurée = distance à la dorsale (un flanc) » et non « distance entre deux points symétriques (deux flancs) ».
\end{attention}

\courssub{Comparaison des taux d'expansion : Atlantique (lent) vs Pacifique (rapide)}

\begin{popfigure}
\begin{tikzpicture}
\begin{axis}[
  width=\linewidth,
  height=8.5cm,
  xlabel={Distance à l'axe de la dorsale (km)},
  ylabel={Âge du fond océanique (Ma)},
  xmin=0, xmax=3000,
  ymin=0, ymax=200,
  grid=major,
  grid style={popline!50},
  legend pos=north west,
  tick label style={font=\small},
  label style={font=\small},
  legend style={font=\small},
]
% Atlantique : v_demi = 2 cm/an = 20 km/Ma -> pente = 1/20 = 0.05 Ma/km
\addplot[popblue, thick, domain=0:200, samples=2] ({x*20}, x);
\addlegendentry{Atlantique Nord ($v_{\text{demi}} \approx \SI{2}{cm/an}$)};
% Pacifique (Dorsale Est-Pacifique) : v_demi = 8 cm/an = 80 km/Ma -> pente = 1/80 = 0.0125 Ma/km
\addplot[poporange, thick, domain=0:200, samples=2] ({x*80}, x);
\addlegendentry{Pacifique Est ($v_{\text{demi}} \approx \SI{8}{cm/an}$)};
% Points expérimentaux simulés (DSDP/ODP)
\addplot[only marks, mark=*, mark size=2.5pt, popblue] coordinates {
  (0,0) (240,12) (480,24) (720,36) (960,48) (1200,60) (1440,72) (1680,84) (1920,96)
};
\addplot[only marks, mark=square*, mark size=2.5pt, poporange] coordinates {
  (0,0) (640,8) (1280,16) (1920,24) (2560,32)
};
% Ligne identité pour référence (non tracée)
\end{axis}
\end{tikzpicture}
\legende{Graphique âge–distance (profil perpendiculaire à la dorsale) pour deux types de dorsales. La pente de la droite ($\Delta \text{âge} / \Delta \text{distance}$) est l'inverse du taux d'expansion demi-dorsale. L'Atlantique (bleu) a une pente forte ($\approx 0,05$ Ma/km) $\rightarrow$ expansion lente ($\approx \SI{2}{cm/an}$). Le Pacifique (orange) a une pente faible ($\approx 0,0125$ Ma/km) $\rightarrow$ expansion rapide ($\approx \SI{8}{cm/an}$). Les points symbolisent des données de forages réels.}
\end{popfigure}

\begin{savaistu}[Pourquoi aucun fond océanique n'excède 200 Ma ?]
La croûte océanique suit un cycle sans fin : elle naît aux dorsales (accrétion), vieillit en s'éloignant (refroidissement, épaississement, sédimentation), et meurt dans les zones de subduction (fosses océaniques) où elle plonge dans le manteau pour y être recyclée. La durée de ce cycle (âge maximal de la croûte océanique) est d'environ \textbf{180 à 200 millions d'années} (Jurassique supérieur). Les roches océaniques plus anciennes ont toutes été subductées. Seuls des fragments arrachés lors de collisions (ophiolites, ex. : complexe de l'Ophiolite de l'Ouest camerounais, ou chaînes alpines) ou les marges passives très anciennes préservent des témoins plus vieux, mais ils ne font plus partie du fond océanique \emph{actif}.
\end{savaistu}

\courssub{Synthèse de la section}
\begin{aretenir}[Points clés à retenir pour le Probatoire]
\begin{itemize}
\item L'âge du fond océanique se détermine par deux méthodes indépendantes et concordantes : \textbf{radiochronologie des basaltes} (\ce{^{40}Ar/^{39}Ar}, \ce{U/Pb}) et \textbf{biostratigraphie des sédiments de base} (foraminifères, nannofossiles).
\item Les \textbf{isochrones} (lignes d'âge égal) sont parallèles à la dorsale et \textbf{symétriques} par rapport à son axe.
\item L'âge est nul à l'axe, maximal (\textbf{180–200 Ma}) aux marges passives ou fosses de subduction.
\item La relation fondamentale $d = v \times t$ (donc $v = d/t$) permet de calculer le \textbf{taux d'expansion demi-dorsale} ($v_{\text{demi}}$) à partir d'un couple (distance $d$ à l'axe, âge $t$). \textbf{Attention :} $d$ = distance un flanc / axe.
\item Le taux d'expansion complet (vitesse relative des plaques) vaut $2 \times v_{\text{demi}}$.
\item Deux régimes majeurs : \textbf{Atlantique} (lent, $\sim \SI{2}{cm/an}$ demi-dorsale) et \textbf{Pacifique} (rapide, $\sim \SI{8}{cm/an}$ à $\SI{10}{cm/an}$ demi-dorsale).
\end{itemize}
\end{aretenir}

\coursec{Les anomalies magnétiques : une bande enregistreuse}

Dans la section précédente, nous avons vu que l'âge des fonds océaniques augmente régulièrement lorsqu'on s'éloigne de l'axe des dorsales : la croûte la plus jeune se situe au niveau de la dorsale, la plus ancienne près des marges continentales. Mais comment les géologues ont-ils pu établir cette chronologie avec autant de précision, alors que les forages océaniques sont coûteux et rares ? La réponse tient en un phénomène remarquable : les fonds océaniques ont \emph{enregistré}, comme sur une bande magnétique, l'histoire des inversions du champ magnétique terrestre. C'est ce magnétisme fossile que nous allons maintenant étudier.

\courssub{Le champ magnétique terrestre et ses inversions}

\textbf{Une observation de départ.} Une boussole indique le nord grâce au champ magnétique terrestre. Ce champ, engendré par les mouvements du fer liquide dans le noyau externe (effet dynamo), n'est pourtant pas éternellement stable : l'étude des laves anciennes montre qu'au cours des temps géologiques, les pôles magnétiques nord et sud se sont \cle{inversés} à de nombreuses reprises.

\begin{definition}[Polarité du champ magnétique terrestre]
Le champ magnétique terrestre possède deux états de \cle{polarité} :
\begin{itemize}
\item la \cle{polarité normale} : le pôle magnétique nord se situe à proximité du pôle géographique nord (situation actuelle) ;
\item la \cle{polarité inversée} : le pôle magnétique nord se situe à proximité du pôle géographique sud (une boussole pointerait alors vers le sud !).
\end{itemize}
Les \cle{inversions magnétiques} sont les passages d'une polarité à l'autre. Elles se produisent de manière \textbf{irrégulière} : aucune périodicité ne gouverne leur survenue.
\end{definition}

\begin{savaistu}[Un champ magnétique capricieux]
Le champ magnétique terrestre s'est inversé des \textbf{centaines de fois} au cours des 160 derniers millions d'années. La dernière inversion, dite de \cle{Brunhes--Matuyama}, date d'il y a environ \num{780000} ans (\SI{0,78}{Ma}). Certains épisodes de polarité stable ont duré des dizaines de millions d'années, d'autres quelques dizaines de milliers d'années seulement. Actuellement, l'intensité du champ magnétique diminue lentement : certains géophysiciens se demandent si une nouvelle inversion ne se prépare pas, mais nulle inquiétude, le processus prend des milliers d'années !
\end{savaistu}

\courssub{L'aimantation thermorémanente : comment une roche enregistre le champ magnétique}

\textbf{Le problème.} Comment peut-on connaître la polarité du champ magnétique il y a 50 millions d'années, alors qu'aucun instrument de mesure n'existait ? La réponse est dans les roches elles-mêmes, et plus précisément dans les \cle{basaltes} qui constituent la croûte océanique.

Les basaltes contiennent des minéraux \cle{ferromagnétiques}, notamment la \cle{magnétite} (\ce{Fe3O4}). Ces minéraux se comportent comme de minuscules aimants microscopiques. Le mécanisme d'enregistrement est le suivant :

\begin{enumerate}
\item À la dorsale, le magma basaltique est émis à très haute température (environ \SI{1200}{\celsius}). À cette température, l'agitation thermique désoriente sans cesse les moments magnétiques des minéraux : la lave n'a pas d'aimantation propre.
\item En refroidissant, la lave se solidifie. Lorsque la température descend en dessous de la \cle{température de Curie} (environ \SI{580}{\celsius} pour la magnétite), les moments magnétiques des cristaux de magnétite se \textbf{figent} dans la direction du champ magnétique terrestre de l'époque.
\item Cette aimantation, dite \cle{thermorémanente}, est conservée durablement dans la roche, comme un enregistrement fossile de la polarité du champ au moment du refroidissement.
\end{enumerate}

\begin{aretenir}[L'aimantation thermorémanente]
Toute roche volcanique contenant des minéraux ferromagnétiques acquiert, en refroidissant sous la température de Curie, une aimantation parallèle au champ magnétique terrestre de l'époque de sa formation. Un basalte formé pendant une période de polarité normale est aimanté dans le sens actuel ; un basalte formé pendant une période de polarité inversée est aimanté en sens inverse.
\end{aretenir}

\begin{experience}[Mesurer l'aimantation d'un basalte]
\textbf{Matériel :} carottes de basalte prélevées par forage dans le plancher océanique, magnétomètre (appareil mesurant l'aimantation des roches).\par
\textbf{Protocole :} on mesure la direction et le sens de l'aimantation de chaque échantillon, puis on date la roche par les méthodes radiométriques (potassium--argon).\par
\textbf{Observations :} certains basaltes ont une aimantation dans le sens du champ actuel (polarité normale), d'autres une aimantation opposée (polarité inversée). En reportant les âges radiométriques, on constate que toutes les roches de même âge, où qu'elles soient dans le monde, présentent la même polarité.\par
\textbf{Interprétation :} les inversions magnétiques sont des événements \textbf{mondiaux et synchrones} : elles constituent donc d'excellents marqueurs chronologiques, datables par les méthodes radiométriques.
\end{experience}

\courssub{La formation des anomalies magnétiques au niveau des dorsales}

\textbf{Démarche d'investigation.} Dès les années 1950, des navires océanographiques remorquant des magnétomètres ont cartographié le champ magnétique au-dessus des océans. Le résultat fut stupéfiant : le long de profils perpendiculaires aux dorsales, on observe une alternance de zones où le champ mesuré est plus fort que le champ théorique et de zones où il est plus faible, dessinant des \textbf{bandes parallèles à la dorsale et symétriques de part et d'autre de son axe}.

\begin{definition}[Anomalie magnétique]
Une \cle{anomalie magnétique} est un écart local entre le champ magnétique mesuré et le champ magnétique théorique (champ moyen attendu en l'absence de roches aimantées).
\begin{itemize}
\item L'anomalie est \cle{positive} lorsque les roches du plancher océanique possèdent une polarité normale : leur aimantation s'\emph{ajoute} au champ actuel, qui est alors mesuré plus fort.
\item L'anomalie est \cle{négative} lorsque les roches possèdent une polarité inversée : leur aimantation se \emph{soustrait} au champ actuel, qui est alors mesuré plus faible.
\end{itemize}
\end{definition}

\textbf{Hypothèse explicative (Vine et Matthews, 1963).} Deux géophysiciens britanniques, Frederick Vine et Drummond Matthews, proposèrent en 1963 un modèle simple et élégant, qui combine l'expansion océanique et les inversions du champ magnétique :

\begin{enumerate}
\item Au niveau de la dorsale, une nouvelle croûte basaltique se forme en permanence et enregistre la polarité du champ magnétique \emph{au moment de son refroidissement}.
\item Cette croûte est ensuite \textbf{coupée en deux} par l'arrivée de magma plus récent à l'axe : chaque moitié s'écarte d'un côté de la dorsale, à la même vitesse.
\item Quand le champ magnétique s'inverse, la nouvelle croûte formée enregistre la nouvelle polarité, tandis que les bandes déjà formées continuent de s'éloigner.
\item Il en résulte, de part et d'autre de la dorsale, des bandes de polarité normale et inversée \textbf{strictement symétriques}, chaque paire de bandes symétriques ayant exactement le même âge.
\end{enumerate}

\begin{propriete}[Modèle de Vine et Matthews (admise)]
Les anomalies magnétiques du plancher océanique forment des \textbf{bandes parallèles à la dorsale et symétriques par rapport à son axe}. Cette symétrie s'explique par le fait que la croûte océanique enregistre la polarité du champ magnétique au moment de sa formation à la dorsale, puis s'écarte de part et d'autre à la même vitesse. Les anomalies magnétiques constituent donc une preuve directe de l'\cle{expansion océanique}.
\end{propriete}

\begin{popfigure}
\begin{center}
\begin{tikzpicture}[scale=1]
% Bandes magnétiques symétriques (largeurs variées, irrégulières)
% Moitié droite : largeurs depuis l'axe
\fill[popdark] (0,0) rectangle (0.7,2.2);
\fill[white] (0.7,0) rectangle (1.3,2.2);
\fill[popdark] (1.3,0) rectangle (1.8,2.2);
\fill[white] (1.8,0) rectangle (2.9,2.2);
\fill[popdark] (2.9,0) rectangle (3.4,2.2);
\fill[white] (3.4,0) rectangle (4.6,2.2);
\fill[popdark] (4.6,0) rectangle (5.6,2.2);
% Moitié gauche (miroir)
\fill[popdark] (-0.7,0) rectangle (0,2.2);
\fill[white] (-1.3,0) rectangle (-0.7,2.2);
\fill[popdark] (-1.8,0) rectangle (-1.3,2.2);
\fill[white] (-2.9,0) rectangle (-1.8,2.2);
\fill[popdark] (-3.4,0) rectangle (-2.9,2.2);
\fill[white] (-4.6,0) rectangle (-3.4,2.2);
\fill[popdark] (-5.6,0) rectangle (-4.6,2.2);
% Cadre
\draw[thick] (-5.6,0) rectangle (5.6,2.2);
% Axe de la dorsale
\draw[poporange,very thick,dashed] (0,-0.4) -- (0,2.8);
\node[poporange,above] at (0,2.8) {\textbf{axe de la dorsale}};
% Flèches d'expansion
\draw[-{Stealth},very thick,popblue] (0.3,1.1) -- (1.1,1.1);
\draw[-{Stealth},very thick,popblue] (-0.3,1.1) -- (-1.1,1.1);
\node[popblue,below] at (0.75,1.0) {\footnotesize expansion};
\node[popblue,below] at (-0.85,1.0) {\footnotesize expansion};
% Numérotation des bandes
\node[white] at (0.35,1.85) {\footnotesize 1};
\node[white] at (-0.35,1.85) {\footnotesize 1};
\node[popdark] at (1.0,1.85) {\footnotesize 2};
\node[popdark] at (-1.0,1.85) {\footnotesize 2};
\node[white] at (1.55,1.85) {\footnotesize 3};
\node[white] at (-1.55,1.85) {\footnotesize 3};
\node[popdark] at (2.35,1.85) {\footnotesize 4};
\node[popdark] at (-2.35,1.85) {\footnotesize 4};
\node[white] at (3.15,1.85) {\footnotesize 5};
\node[white] at (-3.15,1.85) {\footnotesize 5};
\node[popdark] at (4.0,1.85) {\footnotesize 6};
\node[popdark] at (-4.0,1.85) {\footnotesize 6};
% Légende des couleurs
\fill[popdark] (-5.4,-1.0) rectangle (-5.0,-0.6);
\node[right,font=\footnotesize] at (-4.9,-0.8) {polarité normale (anomalie positive)};
\draw[thick] (-0.2,-1.0) rectangle (0.2,-0.6);
\node[right,font=\footnotesize] at (0.3,-0.8) {polarité inversée (anomalie négative)};
% Échelle des âges
\draw[-{Stealth},thick] (-5.6,-0.35) -- (5.6,-0.35);
\node[below,font=\footnotesize] at (0,-0.4) {âge croissant vers les marges};
\end{tikzpicture}
\end{center}
\legende{Profil magnétique perpendiculaire à une dorsale. Les bandes noires correspondent à une polarité normale (anomalies positives), les bandes blanches à une polarité inversée (anomalies négatives). Les bandes de même numéro, symétriques par rapport à l'axe, ont le même âge. Remarque : les largeurs des bandes sont inégales car les inversions sont irrégulières.}
\end{popfigure}

\begin{attention}[Les inversions ne sont pas périodiques !]
Erreur fréquente : croire que le champ magnétique s'inverse à intervalles réguliers, comme un métronome. C'est faux : les inversions sont \textbf{totalement irrégulières}. C'est précisément cette irrégularité qui fait toute la valeur du signal : la succession de bandes de largeurs inégales constitue une véritable \emph{empreinte digitale}, unique et reconnaissable sur tous les océans du globe. Si les inversions étaient périodiques, toutes les bandes se ressembleraient et il serait impossible d'identifier l'âge d'une bande donnée.
\end{attention}

\courssub{L'échelle des inversions magnétiques : un calendrier du plancher océanique}

En datant par les méthodes radiométriques des laves continentales et océaniques de polarité connue, les géologues ont construit une \cle{échelle des inversions magnétiques} : un calendrier indiquant l'âge de chaque inversion au cours des temps géologiques. Cette échelle permet d'attribuer un âge à chaque bande d'anomalie magnétique, sans avoir besoin de forer : il suffit d'identifier la bande par sa position et sa forme dans le profil.

\begin{popfigure}
\begin{center}
\begin{tikzpicture}[scale=1]
% Échelle : 1 Ma = 0.1 cm ; 80 Ma = 8 cm
% Périodes de polarité normale (noires), âges approximatifs en Ma
% Brunhes : 0 - 0.78
\fill[popdark] (0,0) rectangle (0.078,1.4);
% Matuyama normales : 0.99-1.07 (Jaramillo), 1.77-1.95 (Olduvai)
\fill[popdark] (0.099,0) rectangle (0.107,1.4);
\fill[popdark] (0.177,0) rectangle (0.195,1.4);
% Gauss : 2.58 - 3.58 (avec Kaena 3.03-3.11, Mammoth 3.22-3.33 inversées)
\fill[popdark] (0.258,0) rectangle (0.303,1.4);
\fill[popdark] (0.311,0) rectangle (0.322,1.4);
\fill[popdark] (0.333,0) rectangle (0.358,1.4);
% Gilbert : épisode normal vers 4.18-5.0 (simplifié)
\fill[popdark] (0.418,0) rectangle (0.5,1.4);
% Anomalies C5n à C32n (âges en Ma, échelle 0.1 cm/Ma)
\fill[popdark] (0.98,0) rectangle (1.10,1.4);
\fill[popdark] (1.19,0) rectangle (1.27,1.4);
\fill[popdark] (1.36,0) rectangle (1.41,1.4);
\fill[popdark] (1.87,0) rectangle (2.01,1.4);
\fill[popdark] (2.13,0) rectangle (2.30,1.4);
\fill[popdark] (2.40,0) rectangle (2.51,1.4);
\fill[popdark] (2.58,0) rectangle (2.69,1.4);
\fill[popdark] (2.70,0) rectangle (2.83,1.4);
\fill[popdark] (2.94,0) rectangle (3.06,1.4);
\fill[popdark] (3.31,0) rectangle (3.37,1.4);
\fill[popdark] (3.47,0) rectangle (3.53,1.4);
\fill[popdark] (3.57,0) rectangle (3.63,1.4);
\fill[popdark] (3.70,0) rectangle (3.84,1.4);
\fill[popdark] (3.89,0) rectangle (4.13,1.4);
\fill[popdark] (4.25,0) rectangle (4.63,1.4);
\fill[popdark] (4.79,0) rectangle (4.97,1.4);
\fill[popdark] (5.07,0) rectangle (5.17,1.4);
\fill[popdark] (5.19,0) rectangle (5.33,1.4);
\fill[popdark] (5.39,0) rectangle (5.76,1.4);
\fill[popdark] (5.79,0) rectangle (5.92,1.4);
\fill[popdark] (6.09,0) rectangle (6.25,1.4);
\fill[popdark] (6.30,0) rectangle (6.35,1.4);
\fill[popdark] (6.47,0) rectangle (6.77,1.4);
\fill[popdark] (6.80,0) rectangle (6.93,1.4);
\fill[popdark] (6.97,0) rectangle (7.14,1.4);
\fill[popdark] (7.16,0) rectangle (7.36,1.4);
\fill[popdark] (7.39,0) rectangle (7.91,1.4);
% Cadre
\draw[thick] (0,0) rectangle (8,1.4);
% Axe des âges
\foreach \x/\a in {0/0, 1/10, 2/20, 3/30, 4/40, 5/50, 6/60, 7/70, 8/80}{
  \draw (\x,0) -- (\x,-0.12);
  \node[below,font=\footnotesize] at (\x,-0.12) {\a};
}
\node[below,font=\footnotesize] at (4,-0.55) {Âge (millions d'années)};
% Annotations des époques
\draw[poporange,thick] (0.039,1.4) -- (0.039,1.75);
\node[above,font=\footnotesize,poporange,align=center] at (0.5,1.75) {Brunhes\\(normale)};
\node[above,font=\footnotesize,align=center] at (1.7,1.5) {Matuyama\\(inversée)};
\node[above,font=\footnotesize,align=center] at (3.1,1.5) {Gauss\\(normale)};
\node[above,font=\footnotesize,align=center] at (4.0,1.5) {Gilbert\\(inversée)};
% Légende
\fill[popdark] (0.2,-1.35) rectangle (0.5,-1.05);
\node[right,font=\footnotesize] at (0.6,-1.2) {polarité normale};
\draw[thick] (3.0,-1.35) rectangle (3.3,-1.05);
\node[right,font=\footnotesize] at (3.4,-1.2) {polarité inversée};
% Inversion Brunhes-Matuyama
\draw[-{Stealth},popgreen,thick] (0.078,-0.75) -- (0.078,-0.15);
\node[left,font=\footnotesize,popgreen] at (0.06,-0.75) {inversion Brunhes--Matuyama : 0,78 Ma};
\end{tikzpicture}
\end{center}
\legende{Frise chronologique simplifiée des inversions magnétiques au cours des 80 derniers millions d'années. Les intervalles noirs correspondent aux périodes de polarité normale, les intervalles blancs aux périodes de polarité inversée. L'irrégularité des intervalles permet d'identifier chaque bande d'anomalie comme une empreinte digitale.}
\end{popfigure}

\courssub{Interpréter un profil magnétique : la méthode}

\begin{methode}[Interpréter un profil magnétique perpendiculaire à une dorsale]
\begin{enumerate}
\item \textbf{Repérer l'anomalie axiale} : c'est l'anomalie située au niveau de l'axe de la dorsale. Elle correspond à la croûte la plus récente, formée pendant l'épisode de polarité actuel (polarité normale de Brunhes depuis 0,78 Ma).
\item \textbf{Apparier les anomalies symétriques} : de part et d'autre de l'axe, identifier les anomalies de même forme et de même largeur, situées à égale distance de la dorsale. Deux anomalies symétriques ont le \textbf{même âge}.
\item \textbf{Attribuer un âge à chaque bande} : comparer la succession des bandes à l'échelle des inversions magnétiques, comme on comparerait une empreinte digitale à un fichier.
\item \textbf{En déduire le taux d'expansion} : connaissant l'âge $t$ d'une bande et sa distance $d$ à l'axe de la dorsale, on calcule la vitesse d'écartement (demi-vitesse d'expansion) :
\[ v = \dfrac{d}{t} \]
\end{enumerate}
\end{methode}

\begin{exemplebox}[L'anomalie de Brunhes--Matuyama]
La dernière inversion magnétique (Brunhes--Matuyama) date de \SI{0,78}{Ma}. Sur un profil traversant une dorsale dont le demi-taux d'expansion est de \SI{1}{cm/an}, la limite entre l'anomalie axiale (normale, de Brunhes) et la première anomalie négative (de Matuyama) se situe à environ \SI{8}{km} de part et d'autre de l'axe. Vérification :
\[ d = v \times t = \SI{1}{cm/an} \times \SI{780000}{ans} = \SI{780000}{cm} = \SI{7,8}{km} \approx \SI{8}{km}. \]
La concordance entre le calcul et l'observation confirme la validité du modèle.
\end{exemplebox}

\begin{exoresolu}[Calcul d'un taux d'expansion à partir d'un profil magnétique]
Sur un profil magnétique perpendiculaire à la dorsale médio-atlantique, la limite externe de l'anomalie correspondant à la fin de l'épisode normal daté de \SI{2,6}{Ma} se situe à \SI{52}{km} de l'axe de la dorsale.\par
\textbf{Question :} calculer le demi-taux d'expansion de cette dorsale, en cm/an.
\tcblower
\textbf{Solution détaillée.}\par
La croûte située à \SI{52}{km} de l'axe s'est formée il y a \SI{2,6}{Ma} au niveau de la dorsale, puis s'en est écartée à vitesse constante. Le demi-taux d'expansion est donc :
\begin{align*}
v &= \dfrac{d}{t} = \dfrac{\SI{52}{km}}{\SI{2,6}{Ma}} = \dfrac{\SI{5200000}{cm}}{\SI{2600000}{ans}}\\[4pt]
v &= \SI{2}{cm/an}
\end{align*}
Le demi-taux d'expansion est de \SI{2}{cm/an} de chaque côté, soit un écartement total de \SI{4}{cm/an} entre les deux plaques. C'est une valeur typique d'une dorsale à expansion moyenne.
\end{exoresolu}

\courssub{Vérification croisée : anomalies magnétiques et datations radiométriques}

Une théorie scientifique n'est acceptée que si des preuves indépendantes la confirment. Le modèle de Vine et Matthews a été testé par le programme de forages océaniques profonds (Deep Sea Drilling Project, à partir de 1968) : des carottes de basalte ont été prélevées dans le plancher océanique à différentes distances des dorsales, puis datées par les méthodes radiométriques.

\begin{experience}[Le test décisif des forages océaniques]
\textbf{Protocole :} le navire \emph{Glomar Challenger} a foré le plancher océanique de part et d'autre de la dorsale médio-atlantique. Pour chaque forage, on a mesuré la distance à l'axe de la dorsale et daté les basaltes par la méthode potassium--argon.\par
\textbf{Observations :}
\begin{itemize}
\item l'âge radiométrique des basaltes augmente régulièrement avec la distance à la dorsale ;
\item à égale distance de l'axe, de part et d'autre, les basaltes ont le même âge ;
\item les âges radiométriques coïncident avec les âges prédits par l'échelle des inversions magnétiques.
\end{itemize}
\textbf{Conclusion :} la concordance entre les âges magnétiques et les âges radiométriques valide de manière indépendante le modèle de l'expansion océanique. Les anomalies magnétiques sont bien une \emph{bande enregistreuse} fidèle de l'histoire des fonds océaniques.
\end{experience}

\begin{attention}[Ne pas confondre anomalie et inversion]
Dans les copies, deux confusions reviennent souvent :
\begin{itemize}
\item Une \textbf{anomalie magnétique} est une mesure faite \emph{aujourd'hui} au-dessus du plancher océanique ; une \textbf{inversion} est un \emph{événement passé} du champ magnétique terrestre. Les anomalies sont la trace fossile des inversions.
\item Une anomalie négative ne signifie pas « absence de champ magnétique » : le champ est simplement mesuré un peu plus faible que le champ théorique, car l'aimantation des roches inversées s'y soustrait.
\end{itemize}
\end{attention}

\begin{aretenir}[L'essentiel de la section]
\begin{itemize}
\item Le champ magnétique terrestre s'est inversé de manière \textbf{irrégulière} des centaines de fois ; la dernière inversion (Brunhes--Matuyama) date de \SI{0,78}{Ma}.
\item Les basaltes de la croûte océanique enregistrent, en refroidissant sous la température de Curie ($\approx \SI{580}{\celsius}$), la polarité du champ magnétique de l'époque : c'est l'\textbf{aimantation thermorémanente}.
\item La croûte formée à la dorsale s'écartant de part et d'autre, il en résulte des bandes d'anomalies magnétiques \textbf{parallèles et symétriques} par rapport à l'axe (modèle de Vine et Matthews, 1963).
\item L'échelle des inversions magnétiques permet d'attribuer un âge à chaque bande et de calculer le taux d'expansion : $v = d/t$.
\item Les datations radiométriques des basaltes forés confirment les âges magnétiques : double preuve de l'expansion océanique.
\end{itemize}
\end{aretenir}

Ainsi, le plancher océanique fonctionne comme une immense bande enregistreuse qui déroule, de part et d'autre des dorsales, l'histoire magnétique de la Terre. Il ne reste plus qu'à assembler tous ces marqueurs chronologiques — âges radiométriques, anomalies magnétiques, sédimentation — pour reconstituer la chronologie complète de l'histoire des fonds océaniques : ce sera l'objet de la section suivante.

\coursec{Reconstitution chronologique : synthèse des marqueurs}

Dans la section précédente, tu as découvert que les anomalies magnétiques forment une véritable \cle{bande enregistreuse} : le plancher océanique a enregistré, en bandes symétriques de part et d'autre de la dorsale, toutes les inversions du champ magnétique terrestre. Nous disposons donc désormais de \textbf{deux horloges indépendantes} : l'âge des roches et des sédiments (sections 2) et la succession des anomalies magnétiques (section 3). L'objectif de cette section est de les \textbf{combiner} pour reconstituer, étape par étape, l'histoire complète d'un océan — et en particulier celle de l'Atlantique sud, dont l'ouverture a séparé l'Afrique de l'Amérique du Sud... et façonné la façade maritime du Cameroun.

\courssub{Deux marqueurs convergents : la force d'une double preuve}

\textbf{Une intuition pour commencer.} Imagine que tu veuilles reconstituer l'itinéraire d'un chauffeur de taxi-moto entre Douala et Bafoussam. Tu disposes de deux sources indépendantes : le compteur kilométrique du véhicule et les tickets de péage horodatés. Si les deux racontent le même trajet avec les mêmes durées, ta reconstitution est quasi certaine. En géologie marine, le raisonnement est identique : deux horloges indépendantes qui donnent la même chronologie se valident mutuellement.

\begin{aretenir}[Le double marqueur chronologique]
L'histoire d'un océan repose sur deux marqueurs \textbf{indépendants} mais \textbf{convergents} :
\begin{itemize}
\item les \cle{âges radiométriques} des basaltes (datation K/Ar ou Ar/Ar) et les \cle{âges paléontologiques} des premiers sédiments déposés sur la croûte (forages des programmes DSDP, ODP, puis IODP) ;
\item les \cle{anomalies magnétiques}, datées par comparaison avec l'échelle des inversions magnétiques établie sur les coulées continentales.
\end{itemize}
Sur un même profil perpendiculaire à une dorsale, l'âge radiométrique d'une roche correspond toujours à l'âge de l'anomalie magnétique sur laquelle elle se trouve. Cette concordance est la preuve définitive de l'expansion océanique.
\end{aretenir}

\begin{experience}[Vérification croisée des deux horloges (données de référence)]
\textbf{Données.} Le long d'un profil traversant la dorsale médio-atlantique, des forages ont prélevé les basaltes sous les sédiments, et des magnétomètres ont mesuré les anomalies. On obtient, de part et d'autre de l'axe :

\begin{center}
\begin{tabularx}{\linewidth}{|l|X|X|X|X|}\hline
\rowcolor{popblueL} Distance à l'axe (km) & 0 & 200 & 600 & 1 000 \\ \hline
Âge radiométrique (Ma) & 0 & 10 & 30 & 50 \\ \hline
Anomalie magnétique (n°) & 1 & 5 & 13 & 21 \\ \hline
Âge de l'anomalie (Ma) & 0 & 10 & 33 & 50 \\ \hline
\end{tabularx}
\end{center}

\textbf{Interprétation.} Les âges radiométriques et les âges magnétiques coïncident (aux incertitudes près). Les deux horloges, fondées sur des phénomènes physiques totalement différents (désintégration radioactive d'une part, inversions du champ magnétique d'autre part), racontent la même histoire.\\
\textbf{Conclusion.} L'expansion océanique est un fait établi, et sa chronologie peut être reconstituée avec confiance.
\end{experience}

\begin{attention}[Ne pas confondre les deux marqueurs]
Une erreur classique au Probatoire consiste à écrire que « les anomalies magnétiques servent à dater les roches par radioactivité ». Non ! Les deux méthodes sont \textbf{indépendantes} : la radiochronologie date la \emph{solidification} du basalte, l'anomalie magnétique date l'\emph{inversion du champ} enregistrée au moment de cette solidification. C'est précisément parce qu'elles sont indépendantes et concordantes que la preuve est solide.
\end{attention}

\courssub{Les étapes d'une reconstitution chronologique complète}

\begin{methode}[Reconstituer l'histoire d'un océan en 4 étapes]
\begin{enumerate}
\item \textbf{Dater} les fonds : utiliser la carte des âges (isochrones) ou identifier les anomalies magnétiques sur un profil, en les rapportant à l'échelle des inversions.
\item \textbf{Mesurer} sur la carte, pour chaque isochrone, la distance qui le sépare de l'axe de la dorsale (avec l'échelle de la carte).
\item \textbf{Calculer} le taux d'expansion : $v = \dfrac{d}{t}$, où $d$ est la distance à l'axe et $t$ l'âge de l'isochrone. Vérifier la régularité du taux en répétant le calcul pour plusieurs isochrones.
\item \textbf{Extrapoler} la position passée des continents : en « refermant » l'océan le long des isochrones (on rapproche symétriquement les points de même âge), on retrouve la configuration des continents à chaque époque, jusqu'à leur jonction initiale.
\end{enumerate}
\end{methode}

\begin{propriete}[Principe de la reconstitution à rebours]
La \textbf{symétrie} des âges et des anomalies magnétiques de part et d'autre de la dorsale permet de reconstituer la position passée des continents : deux points situés sur un même isochrone, de chaque côté de l'axe, \textbf{étaient accolés au moment de leur formation}. En rapprochant virtuellement les isochrones de même âge jusqu'à l'axe, on « referme » l'océan et on retrouve la position initiale des continents avant l'ouverture. C'est le principe de la \cle{reconstitution à rebours} (ou \emph{refitting}).
\end{propriete}

\begin{exoresolu}[Calcul d'un taux d'expansion dans l'Atlantique sud]
\textbf{Énoncé.} Sur une carte des âges de l'Atlantique sud, à la latitude du golfe de Guinée, l'isochrone de 80 Ma se situe à \SI{1600}{km} de l'axe de la dorsale médio-atlantique (sur le flanc africain). (1) Calcule le taux d'expansion moyen de ce flanc. (2) Sachant que la côte camerounaise est aujourd'hui à environ \SI{5000}{km} de l'axe, estime l'âge du début de l'ouverture. (3) Compare avec l'âge des plus anciens fonds connus de cette zone ($\approx$ 125--130 Ma) et conclus.
\tcblower
\textbf{Solution détaillée.}
\begin{enumerate}
\item Le taux d'expansion (par flanc) est :
\[ v = \frac{d}{t} = \frac{\SI{1600}{km}}{\SI{80}{Ma}} = \SI{20}{km/Ma} \]
Or $\SI{20}{km/Ma} = \dfrac{20 \times 10^{5}\ \text{cm}}{10^{6}\ \text{ans}} = \SI{2}{cm/an}$. Le flanc africain s'éloigne de la dorsale à \SI{2}{cm/an} en moyenne.
\item En supposant ce taux constant :
\[ t = \frac{d}{v} = \frac{\SI{5000}{km}}{\SI{20}{km/Ma}} = \SI{250}{Ma}\ \text{?} \]
Attention : la côte n'est pas un point formé à la dorsale ! La distance côte--axe comprend la \textbf{marge continentale} (croûte continentale amincie), qui ne s'est pas formée par accrétion à la dorsale. Il faut raisonner sur la \textbf{largeur océanique totale} : environ \SI{5000}{km} de croûte océanique, soit \SI{2500}{km} par flanc :
\[ t = \frac{\SI{2500}{km}}{\SI{20}{km/Ma}} = \SI{125}{Ma}. \]
\item L'âge calculé (125 Ma) est remarquablement proche de l'âge mesuré des plus anciens fonds (125--130 Ma, Crétacé inférieur). L'hypothèse d'un taux moyen quasi constant est donc validée, et la reconstitution est cohérente : l'Atlantique sud s'est ouvert il y a environ 130 millions d'années.
\end{enumerate}
\end{exoresolu}

\courssub{Application : l'ouverture de l'Atlantique sud, du Crétacé à nos jours}

\textbf{Situation concrète.} Un pêcheur de Kribi, au sud du Cameroun, ignore que sous sa pirogue s'étendent des sédiments épais de plusieurs kilomètres, déposés sur une marge continentale née d'une déchirure : celle qui a séparé l'Afrique de l'Amérique du Sud. Reconstituons cette histoire.

\begin{itemize}
\item \textbf{Vers 130 Ma (Crétacé inférieur)} : un rift se creuse au sein du supercontinent Gondwana, entre l'Afrique et l'Amérique du Sud. La lithosphère s'amincit, se fracture, et le premier plancher océanique apparaît au niveau d'une dorsale naissante. L'océan n'est alors qu'une étroite mer.
\item \textbf{Vers 65 Ma (limite Crétacé--Cénozoïque)} : l'expansion se poursuit de part et d'autre de la dorsale médio-atlantique. L'Atlantique sud atteint déjà plus de \SI{2500}{km} de large. Les deux continents dérivent l'un de l'autre.
\item \textbf{Aujourd'hui} : l'Atlantique sud mesure environ \SI{5000}{km} de large. La dorsale médio-atlantique reste exactement à mi-chemin entre les deux continents : preuve de la symétrie de l'expansion. L'océan continue de s'élargir d'environ \SI{4}{cm/an} au total (2 cm/an par flanc).
\end{itemize}

\begin{popfigure}
\begin{tikzpicture}[scale=1]
% ===== Étape 1 : 130 Ma =====
\begin{scope}
\fill[popgreenL] (-0.9,0) rectangle (0.1,3);
\fill[popgreenL] (0.3,0) rectangle (1.3,3);
\fill[popblueL] (0.1,0) rectangle (0.3,3);
\draw[popdark,thick] (-0.9,0) rectangle (0.1,3);
\draw[popdark,thick] (0.3,0) rectangle (1.3,3);
\draw[poporange,very thick] (0.2,0) -- (0.2,3);
\node[font=\small] at (-0.4,1.5) {Afrique};
\node[font=\small] at (0.8,1.5) {Am. du Sud};
\node[font=\footnotesize\itshape, text=popblue] at (0.2,-0.35) {mer étroite};
\node[font=\bfseries] at (0.2,-0.85) {130 Ma};
\end{scope}
% ===== Étape 2 : 65 Ma =====
\begin{scope}[xshift=4.2cm]
\fill[popgreenL] (-1.2,0) rectangle (-0.2,3);
\fill[popgreenL] (1.0,0) rectangle (2.0,3);
\fill[popblueL] (-0.2,0) rectangle (1.0,3);
\draw[popdark,thick] (-1.2,0) rectangle (-0.2,3);
\draw[popdark,thick] (1.0,0) rectangle (2.0,3);
\draw[poporange,very thick] (0.4,0) -- (0.4,3);
\node[font=\small] at (-0.7,1.5) {Afrique};
\node[font=\small] at (1.5,1.5) {Am. Sud};
\node[font=\footnotesize\itshape, text=popblue] at (0.4,-0.35) {océan $\approx$ 2 500 km};
\node[font=\bfseries] at (0.4,-0.85) {65 Ma};
\end{scope}
% ===== Étape 3 : actuel =====
\begin{scope}[xshift=9.4cm]
\fill[popgreenL] (-1.6,0) rectangle (-0.6,3);
\fill[popgreenL] (1.8,0) rectangle (2.8,3);
\fill[popblueL] (-0.6,0) rectangle (1.8,3);
\draw[popdark,thick] (-1.6,0) rectangle (-0.6,3);
\draw[popdark,thick] (1.8,0) rectangle (2.8,3);
\draw[poporange,very thick] (0.6,0) -- (0.6,3);
\node[font=\small] at (-1.1,1.5) {Afrique};
\node[font=\small] at (2.3,1.5) {Am. Sud};
\node[font=\footnotesize\itshape, text=popblue] at (0.6,-0.35) {océan $\approx$ 5 000 km};
\node[font=\bfseries] at (0.6,-0.85) {Actuel};
\draw[poporange,->,thick] (0.6,3.15) -- (-0.2,3.15);
\draw[poporange,->,thick] (0.6,3.15) -- (1.4,3.15);
\node[font=\footnotesize, text=poporange] at (0.6,3.45) {expansion 2 cm/an par flanc};
\end{scope}
% légende des couleurs
\fill[popgreenL] (0,-1.6) rectangle (0.4,-1.3); \node[right,font=\footnotesize] at (0.5,-1.45) {continents};
\fill[popblueL] (3.2,-1.6) rectangle (3.6,-1.3); \node[right,font=\footnotesize] at (3.7,-1.45) {croûte océanique};
\draw[poporange,very thick] (6.4,-1.6) -- (6.8,-1.3); \node[right,font=\footnotesize] at (6.9,-1.45) {dorsale médio-atlantique};
\end{tikzpicture}
\legende{Ouverture progressive de l'Atlantique sud en trois étapes (130 Ma, 65 Ma, actuel). La dorsale (trait orange) engendre en continu une croûte océanique nouvelle (bleu) qui repousse symétriquement l'Afrique et l'Amérique du Sud (vert). La dorsale reste à mi-chemin entre les continents : signature de la symétrie de l'expansion.}
\end{popfigure}

\begin{popfigure}
\begin{center}
\begin{tikzpicture}
\begin{axis}[
width=0.85\linewidth, height=7cm,
xlabel={Temps (Ma avant le présent)},
ylabel={Largeur de l'Atlantique sud (km)},
xmin=0, xmax=140, ymin=0, ymax=5500,
x dir=reverse,
grid=major, grid style={popline!40},
xtick={0,20,40,60,80,100,120,140},
ytick={0,1000,2000,3000,4000,5000},
]
\addplot[popcurve,domain=0:130,samples=100] {5000/130*x};
\addplot[only marks, mark=*, mark size=2pt, color=poporange] coordinates {(0,0) (20,770) (40,1540) (65,2500) (80,3080) (100,3850) (130,5000)};
\node[font=\footnotesize, text=popdark, anchor=west] at (axis cs:60,1500) {pente = taux d'expansion total};
\node[font=\footnotesize, text=popdark, anchor=west] at (axis cs:60,1050) {$\approx$ 38 km/Ma $\approx$ 4 cm/an};
\end{axis}
\end{tikzpicture}
\end{center}
\legende{Largeur de l'Atlantique sud en fonction du temps, construite en reportant pour chaque isochrone sa distance cumulée à l'axe (les points correspondent aux isochrones mesurés sur la carte des âges). L'alignement quasi parfait des points montre que le taux d'expansion moyen a été presque constant ($\approx$ 4 cm/an au total, soit 2 cm/an par flanc). L'axe des temps est orienté vers la gauche : on lit le passé à droite.}
\end{popfigure}

\courssub{Le cas du golfe de Guinée : une marge passive héritée de cette ouverture}

Le littoral camerounais, de Limbé à Kribi, borde le golfe de Guinée. Or cette façade maritime n'est \textbf{pas} une zone active : il n'y a ni dorsale, ni fosse, ni séismes majeurs liés à une limite de plaque. C'est une \cle{marge passive} : un bord de continent formé lors de l'ouverture de l'Atlantique sud, puis devenu inerte, simplement recouvert par d'épais sédiments (sables, argiles, et gisements de pétrole exploités au large, par exemple dans le bassin du Rio del Rey).

\begin{aretenir}[Marge passive]
Une \cle{marge passive} est un bord continental situé \textbf{à l'intérieur} d'une plaque lithosphérique, loin de toute limite active. Elle est née de la rupture d'un continent (rift devenu océan) et s'enfonce lentement sous le poids des sédiments. Les marges du golfe de Guinée (Cameroun, Nigeria, Gabon) et celles du Brésil sont des marges passives \textbf{conjugées} : elles se faisaient face avant l'ouverture de l'Atlantique sud.
\end{aretenir}

\begin{savaistu}[Le puzzle Afrique--Amérique du Sud]
Dès 1912, Alfred Wegener avait remarqué que le contour de la côte brésilienne s'emboîte presque parfaitement dans celui du golfe de Guinée — y compris la côte camerounaise — comme deux pièces d'un puzzle. Ce n'est pas un hasard : ces deux rivages étaient \textbf{joints} avant le Crétacé. La reconstitution à rebours le long des isochrones les remet exactement en contact, et les géologues retrouvent de part et d'autre les mêmes roches anciennes, les mêmes fossiles et les mêmes ceintures de montagnes interrompues. Le puzzle de Wegener est devenu, grâce aux marqueurs océaniques, une reconstitution quantitative précise.
\end{savaistu}

\courssub{Limites de la méthode : une mémoire océanique incomplète}

La reconstitution chronologique à partir des fonds océaniques est puissante, mais elle a une limite fondamentale : \textbf{les fonds océaniques les plus anciens ont disparu}.

\begin{itemize}
\item La lithosphère océanique, en s'éloignant de la dorsale, refroidit, s'épaissit et devient plus dense. Elle finit par plonger dans le manteau au niveau des \textbf{fosses de subduction} (par exemple la fosse du Pérou--Chili), où elle est recyclée.
\item Conséquence : aucune croûte océanique actuelle ne dépasse environ \textbf{180 à 200 Ma} (Jurassique). La « bande enregistreuse » magnétique et la carte des âges ne conservent donc que les 200 derniers millions d'années de l'histoire océanique.
\item Pour les époques plus anciennes, il faut utiliser d'autres indices : les chaînes de montagnes, les fossiles et les roches des continents (qui, eux, ne subduisent pas).
\end{itemize}

\begin{attention}[Deux pièges fréquents au Probatoire]
\begin{enumerate}
\item \textbf{Oublier que le taux d'expansion peut varier.} Le calcul $t = d/v$ avec un taux unique suppose une expansion constante. En réalité, le taux a pu fluctuer au cours du temps : une reconstitution rigoureuse utilise \textbf{plusieurs isochrones} et vérifie l'alignement des points sur le graphique largeur--temps (comme dans la figure ci-dessus). Un point qui s'écarte de la droite signale une variation du taux.
\item \textbf{Croire que l'on peut reconstituer toute l'histoire de la Terre avec les fonds océaniques.} À cause de la subduction, la chronologie océanique plafonne à $\approx$ 200 Ma, alors que la Terre a 4,6 milliards d'années. Ne jamais extrapoler au-delà des données !
\end{enumerate}
\end{attention}

\begin{exercice}[Reconstitution à rebours]
Sur une carte des âges de l'Atlantique sud, on repère les isochrones 20 Ma, 50 Ma et 100 Ma, situés respectivement à 400 km, 1 050 km et 2 200 km de l'axe (flanc africain).
\begin{enumerate}
\item Calcule le taux d'expansion pour chaque isochrone. Le taux a-t-il été constant ?
\item Où se trouvaient, il y a 50 Ma, deux points aujourd'hui situés sur l'isochrone 50 Ma, de part et d'autre de l'axe ?
\item Pourquoi ne trouve-t-on pas d'isochrone de 300 Ma dans l'océan Atlantique ?
\end{enumerate}
\end{exercice}

\begin{corrige}[Exercice : reconstitution à rebours]
\begin{enumerate}
\item $v_1 = 400/20 = \SI{20}{km/Ma}$ ; $v_2 = 1050/50 = \SI{21}{km/Ma}$ ; $v_3 = 2200/100 = \SI{22}{km/Ma}$. Le taux est \textbf{quasi constant} ($\approx$ 2 cm/an par flanc), avec une très légère accélération récente. L'approximation d'un taux constant est donc acceptable.
\item Il y a 50 Ma, ces deux points \textbf{étaient accolés à l'axe de la dorsale}, car ils viennent de se former à cet instant : c'est le principe de la reconstitution à rebours le long des isochrones.
\item Un isochrone de 300 Ma n'existe pas car la croûte océanique de cet âge a été \textbf{recyclée par subduction} (ou n'a jamais existé dans l'Atlantique, océan ouvert il y a moins de 200 Ma). La mémoire des fonds océaniques ne remonte pas au-delà d'environ 200 Ma.
\end{enumerate}
\end{corrige}

En définitive, la convergence des âges radiométriques et des anomalies magnétiques transforme le plancher océanique en une \textbf{archive chronologique lisible} : en mesurant, datant et « refermant » l'océan le long des isochrones, nous reconstituons le film de la séparation des continents — y compris celui qui a donné au Cameroun sa façade atlantique. Reste une dernière question : quelle machine gigantesque met ainsi les continents en mouvement ? C'est l'objet de la section suivante : la \textbf{tectonique des plaques}.

\coursec{De l'expansion océanique à la tectonique des plaques}

Dans la section précédente, nous avons appris à reconstituer la chronologie de l'ouverture d'un océan en croisant deux marqueurs : l'\cle{âge} des basaltes et les \cle{anomalies magnétiques} symétriques de part et d'autre de la dorsale. Ces données montrent que le plancher océanique se renouvelle en permanence : il naît à la dorsale, vieillit en s'en écartant. Mais une question fondamentale demeure : si la surface des océans grandit sans cesse, pourquoi la Terre n'enfle-t-elle pas ? Et surtout : \emph{qu'est-ce qui bouge exactement} ? La réponse à ces questions constitue l'une des plus grandes révolutions scientifiques du \textsc{xx}\textsuperscript{e} siècle : la \cle{tectonique des plaques}.

\courssub{La lithosphère : une enveloppe rigide découpée en plaques}

\textbf{Une observation troublante : les continents s'emboîtent comme des pièces d'un puzzle.}\par

Regarde une carte du monde : les côtes orientales de l'Amérique du Sud semblent épouser parfaitement les côtes occidentales de l'Afrique, du golfe de Guinée jusqu'au Brésil. Dès 1912, le météorologue allemand Alfred Wegener propose une hypothèse audacieuse : les continents auraient autrefois été réunis en un supercontinent, la \cle{Pangée}, puis se seraient séparés en dérivant. Ses arguments sont solides : on retrouve les mêmes fossiles (comme le \emph{Mesosaurus}, petit reptile d'eau douce incapable de traverser un océan) de part et d'autre de l'Atlantique Sud, et les mêmes séries géologiques se prolongent d'un continent à l'autre. Pourtant, sa théorie est rejetée : il ne peut pas expliquer \emph{comment} des continents pourraient se déplacer. Il faudra attendre les années 1960, et la découverte de l'expansion océanique, pour comprendre le mécanisme.

\textbf{La structure mécanique du globe.}\par

Les études sismiques ont montré que l'enveloppe superficielle de la Terre peut être décrite selon son comportement mécanique :
\begin{itemize}
\item la \cle{lithosphère}, rigide et cassante, comprend la croûte (continentale ou océanique) \textbf{et} la partie supérieure du manteau ; son épaisseur varie d'environ \SI{100}{km} sous les océans (lithosphère âgée) à \SI{150}{--}\SI{200}{km} sous les continents ;
\item l'\cle{asthénosphère}, située en dessous (entre environ \SI{100}{km} et \SI{350}{km} de profondeur), est une zone du manteau partiellement fondue, \textbf{ductile} : elle peut s'écouler lentement, à la manière d'un mastic très visqueux.
\end{itemize}

\begin{definition}[Plaque lithosphérique]
Une \cle{plaque lithosphérique} est un fragment rigide de la lithosphère, constitué de la croûte (océanique et/ou continentale) et de la partie supérieure rigide du manteau, qui se déplace en glissant sur l'asthénosphère ductile. La surface du globe est découpée en une douzaine de grandes plaques (plaque africaine, plaque sud-américaine, plaque pacifique, plaque eurasiatique, etc.) et en plusieurs plaques plus petites.
\end{definition}

\begin{attention}[Le piège classique du « continent qui flotte »]
Il est très fréquent de lire ou d'entendre que « les continents flottent sur l'océan » ou « dérivent sur la mer ». C'est \textbf{faux}. Un continent ne se déplace pas à travers l'océan comme un bateau : continents et fonds océaniques sont \textbf{embarqués ensemble sur les mêmes plaques lithosphériques}. C'est la plaque entière — avec sa partie continentale et sa partie océanique — qui se déplace d'un seul bloc sur l'asthénosphère. Au Probatoire, cette confusion coûte cher : dis toujours que les \emph{plaques} bougent, pas les continents seuls.
\end{attention}

\courssub{Les dorsales : des frontières divergentes entre plaques}

\textbf{Observation.} Les données accumulées dans les sections précédentes convergent : la dorsale médio-atlantique sépare deux domaines de plancher océanique dont les âges et les anomalies magnétiques sont symétriques. Les séismes et les volcans sous-marins se concentrent précisément le long de l'axe de la dorsale.

\textbf{Hypothèse.} La dorsale marque la limite entre deux plaques lithosphériques qui s'écartent l'une de l'autre.

\textbf{Vérification.} Le magma basaltique qui remonte le long de l'axe de la dorsale comble en permanence l'espace libéré par l'écartement. Chaque nouvelle injection de magma pousse les bords des deux plaques vers l'extérieur : c'est exactement ce qu'enregistrent les bandes magnétiques symétriques.

\textbf{Conclusion.} La dorsale est une \cle{frontière divergente} : c'est là que deux plaques se séparent et que la lithosphère océanique est créée.

\begin{propriete}[Expansion océanique et divergence des plaques — admise à ce niveau]
L'expansion océanique au niveau des dorsales est l'expression, à la surface du globe, de la \cle{divergence} de deux plaques lithosphériques. Le magma ascendant à l'axe de la dorsale forme de la nouvelle lithosphère océanique qui s'accolde aux bords des deux plaques et les pousse en directions opposées. Le taux d'expansion mesuré de part et d'autre de la dorsale correspond donc à la vitesse d'écartement des plaques.
\end{propriete}

\begin{popfigure}
\begin{center}
\begin{tikzpicture}[scale=1.0, every node/.style={font=\small}]
% Asthénosphère
\fill[poporangeL] (-7,-2.8) rectangle (7,-1.4);
\node[popdark] at (0,-2.5) {\textbf{Asthénosphère (ductile)}};
% Plaque gauche
\fill[popblueL] (-7,0) rectangle (-0.7,-1.4);
\draw[thick, popblue] (-7,0) -- (-0.7,0) -- (-0.7,-1.4) -- (-7,-1.4);
\node[popdark] at (-3.85,-0.7) {\textbf{Plaque A}};
% Plaque droite
\fill[popblueL] (0.7,0) rectangle (7,-1.4);
\draw[thick, popblue] (0.7,0) -- (7,0) -- (7,-1.4) -- (0.7,-1.4);
\node[popdark] at (3.85,-0.7) {\textbf{Plaque B}};
% Dorsale (relief axial)
\fill[popblue!50] (-0.7,0) -- (-0.35,0.55) -- (0,0.35) -- (0.35,0.55) -- (0.7,0) -- cycle;
\draw[thick, popblue] (-0.7,0) -- (-0.35,0.55) -- (0,0.35) -- (0.35,0.55) -- (0.7,0);
% Magma ascendant
\fill[poporange] (-0.28,-2.8) rectangle (0.28,0.3);
\draw[-{Stealth[length=3mm]}, very thick, poporange] (0,-2.6) -- (0,0.15);
\node[poporange, font=\footnotesize\bfseries, align=center] at (0,-3.25) {Magma ascendant};
% Flèches de divergence
\draw[-{Stealth[length=3.5mm]}, very thick, popgreen] (-1.2,0.9) -- (-3.2,0.9);
\draw[-{Stealth[length=3.5mm]}, very thick, popgreen] (1.2,0.9) -- (3.2,0.9);
\node[popgreen, font=\footnotesize\bfseries] at (-2.2,1.25) {divergence};
\node[popgreen, font=\footnotesize\bfseries] at (2.2,1.25) {divergence};
% Anomalies magnétiques symétriques
\foreach \x in {-6.3,-5.6,-4.9,-4.2,-3.5,-2.8,-2.1,-1.4} {\fill[popmuted!60] (\x,0.06) rectangle (\x+0.35,0.22);}
\foreach \x in {1.05,1.75,2.45,3.15,3.85,4.55,5.25,5.95} {\fill[popmuted!60] (\x,0.06) rectangle (\x+0.35,0.22);}
\node[popmuted, font=\footnotesize, align=center] at (-4.5,0.55) {anomalies magnétiques\\symétriques};
\node[popmuted, font=\footnotesize, align=center] at (4.5,0.55) {anomalies magnétiques\\symétriques};
% Subduction à droite
\fill[popblueL] (7,0) -- (8.6,0) -- (8.6,-0.5) -- (7.6,-1.8) -- (7,-1.4) -- cycle;
\draw[thick, popblue] (7,0) -- (8.6,0) (7,-1.4) -- (7.6,-1.8);
\draw[-{Stealth[length=3mm]}, very thick, poppink] (7.9,-0.35) -- (7.5,-1.55);
\node[poppink, font=\footnotesize\bfseries, align=center] at (8.3,-2.2) {subduction\\(fosse)};
% Continents
\fill[popgoldL] (-7,0) rectangle (-5.2,0.5);
\draw[thick, popgold] (-7,0) rectangle (-5.2,0.5);
\node[popdark, font=\footnotesize] at (-6.1,0.25) {continent};
\fill[popgoldL] (5.6,0) rectangle (6.9,0.5);
\draw[thick, popgold] (5.6,0) rectangle (6.9,0.5);
\node[popdark, font=\footnotesize] at (6.25,0.25) {continent};
% Légendes pointées
\draw[thin, popmuted] (-0.1,0.45) -- (-1.6,1.9) node[above, popdark, font=\footnotesize\bfseries] {Dorsale (frontière divergente)};
\end{tikzpicture}
\end{center}
\legende{Schéma-bilan de l'expansion océanique : deux plaques lithosphériques divergent au niveau de la dorsale, où le magma ascendant crée de la nouvelle lithosphère enregistrant les anomalies magnétiques symétriques. Les continents sont embarqués passivement sur les plaques. À l'autre extrémité d'une plaque, la lithosphère océanique peut être consommée par subduction dans une fosse.}
\end{popfigure}

\courssub{Conséquence : la dérive des continents}

Si l'océan Atlantique s'ouvre en permanence à la dorsale médio-atlantique, alors l'Afrique et l'Amérique du Sud, solidaires chacune de leur plaque, s'éloignent l'une de l'autre. L'expansion océanique est donc le \textbf{moteur} de la dérive des continents : ce n'est pas le continent qui « fend » l'océan, c'est l'océan qui s'élargit en poussant les continents.

\begin{exemplebox}[L'ouverture de l'Atlantique Sud]
Il y a environ 130 millions d'années, l'Afrique et l'Amérique du Sud formaient un seul bloc. Une dorsale s'est mise en place entre eux ; depuis, l'expansion océanique les écarte à un rythme moyen de 3 à \SI{4}{cm/an}. Vérification par le calcul : à \SI{3,5}{cm/an} pendant 130 millions d'années, on obtient $3{,}5 \times 130\,000\,000 = 455\,000\,000$ cm, soit environ \SI{4550}{km} — très exactement l'ordre de grandeur de la largeur actuelle de l'Atlantique Sud. Les âges des basaltes, mesurés au laboratoire, confirment que les fonds les plus proches des côtes africaines et brésiliennes sont bien les plus anciens.
\end{exemplebox}

\begin{savaistu}[Le GPS confirme la dérive en direct]
Aujourd'hui, plus besoin d'attendre des millions d'années pour mesurer la dérive : des stations GPS permanentes installées en Afrique et en Amérique du Sud mesurent directement leur écartement : environ 3 à \SI{4}{cm/an}, soit la vitesse de croissance de tes ongles ! Ces mesures confirment avec une précision remarquable les vitesses calculées par les géologues à partir des anomalies magnétiques. La tectonique des plaques n'est donc plus une hypothèse : c'est un fait mesuré en temps réel.
\end{savaistu}

\courssub{Complément utile au Probatoire : la subduction équilibre le bilan}

\begin{attention}[Hors cœur de leçon, mais précieux]
Ce qui suit dépasse légèrement le cœur du programme de Première, mais il est fréquemment valorisé au Probatoire car il permet de répondre à la question piège : « si les océans s'ouvrent, la Terre grossit-elle ? »
\end{attention}

La Terre conserve une surface constante. Or les dorsales créent en permanence de la nouvelle lithosphère océanique. Il doit donc exister des zones où de la lithosphère est \textbf{détruite} : ce sont les \cle{zones de subduction}, associées aux \cle{fosses océaniques} (frontières convergentes). Là, la plaque océanique, devenue froide et dense en vieillissant, plonge sous une autre plaque et s'enfonce dans le manteau, où elle est recyclée. La fosse du Pérou-Chili, le long de la côte ouest de l'Amérique du Sud, en est un exemple célèbre : la plaque de Nazca y plonge sous la plaque sud-américaine, provoquant séismes et volcans (la cordillère des Andes).

\begin{aretenir}[Bilan de la tectonique des plaques]
La lithosphère est découpée en plaques rigides qui glissent sur l'asthénosphère ductile. La surface océanique \textbf{créée} aux dorsales (frontières divergentes) est \textbf{consommée} dans les zones de subduction (frontières convergentes) : le bilan est équilibré et la Terre garde une surface constante. L'expansion océanique est ainsi le moteur visible du mouvement des plaques et de la dérive des continents.
\end{aretenir}

\courssub{Carte des grandes plaques et méthode d'analyse}

\begin{popfigure}
\begin{center}
\begin{tikzpicture}[scale=0.92, every node/.style={font=\footnotesize}]
% Fond océanique
\fill[popblueL!60] (-7.5,-3.2) rectangle (7.5,3.2);
% Plaques (formes simplifiées)
\fill[popgreenL, draw=popgreen, thick] (-7.5,-3.2) -- (-7.5,3.2) -- (-4.2,3.2) -- (-3.6,1.2) -- (-4.0,-1.0) -- (-4.8,-3.2) -- cycle;
\fill[popgoldL, draw=popgold, thick] (-4.2,3.2) -- (-1.2,3.2) -- (-0.8,1.0) -- (-1.4,-0.8) -- (-2.2,-3.2) -- (-4.8,-3.2) -- (-4.0,-1.0) -- (-3.6,1.2) -- cycle;
\fill[poppinkL, draw=poppink, thick] (-1.2,3.2) -- (2.4,3.2) -- (2.0,0.6) -- (1.2,-1.2) -- (0.2,-3.2) -- (-2.2,-3.2) -- (-1.4,-0.8) -- (-0.8,1.0) -- cycle;
\fill[poporangeL, draw=poporange, thick] (2.4,3.2) -- (7.5,3.2) -- (7.5,-0.6) -- (5.2,-1.0) -- (3.0,-0.2) -- (2.0,0.6) -- cycle;
\fill[poppurpleL, draw=poppurple, thick] (3.0,-0.2) -- (5.2,-1.0) -- (7.5,-0.6) -- (7.5,-3.2) -- (0.2,-3.2) -- (1.2,-1.2) -- cycle;
% Noms des plaques
\node[popdark, font=\small\bfseries, align=center] at (-5.8,0.4){Plaque\\nord-américaine};
\node[popdark, font=\small\bfseries, align=center] at (-2.6,0.6){Plaque\\sud-américaine};
\node[popdark, font=\small\bfseries, align=center] at (0.6,1.4){Plaque\\africaine};
\node[popdark, font=\small\bfseries, align=center] at (4.9,1.6){Plaque\\eurasiatique};
\node[popdark, font=\small\bfseries, align=center] at (4.0,-2.0){Plaque\\indo-australienne};
% Dorsale médio-atlantique (ligne divergente)
\draw[very thick, poporange, decorate, decoration={zigzag, segment length=4, amplitude=1.2}] (-4.0,3.2) -- (-3.4,1.0) -- (-2.9,-0.6) -- (-3.4,-3.2);
\node[poporange, font=\footnotesize\bfseries, align=center] at (-3.3,-2.6) {dorsale\\médio-atlantique};
% Flèches de divergence Atlantique
\draw[-{Stealth[length=2.8mm]}, very thick, popgreen] (-3.9,0.2) -- (-4.9,0.2);
\draw[-{Stealth[length=2.8mm]}, very thick, popgreen] (-2.9,0.2) -- (-1.9,0.2);
% Dorsale océan Indien + flèches
\draw[very thick, poporange, decorate, decoration={zigzag, segment length=4, amplitude=1.2}] (1.9,0.4) -- (2.8,-0.4) -- (3.4,-1.1);
\draw[-{Stealth[length=2.8mm]}, very thick, popgreen] (2.4,-0.5) -- (1.6,-1.3);
\draw[-{Stealth[length=2.8mm]}, very thick, popgreen] (3.2,-0.7) -- (4.0,-1.5);
% Subduction Pérou-Chili
\draw[very thick, poppink] (-7.4,-0.4) -- (-4.9,-0.9);
\foreach \x in {-7.1,-6.5,-5.9,-5.3} {\draw[thick, poppink] (\x,-0.62) -- (\x+0.18,-1.0);}
\node[poppink, font=\footnotesize\bfseries, align=center] at (-6.1,-1.7) {fosse Pérou-Chili\\(subduction)};
\end{tikzpicture}
\end{center}
\legende{Carte simplifiée des grandes plaques lithosphériques. Les lignes orange en zigzag représentent les dorsales (frontières divergentes, flèches vertes d'écartement) ; la ligne rose dentée représente une zone de subduction (frontière convergente), les dents du côté de la plaque chevauchante.}
\end{popfigure}

\begin{methode}[Relier expansion océanique et mouvement des plaques]
Face à un document (carte des âges, profil d'anomalies magnétiques, carte des plaques) :
\begin{enumerate}
\item \textbf{Identifier la dorsale} sur le document : c'est la frontière divergente entre deux plaques (axe de symétrie des anomalies magnétiques, zone d'âge nul).
\item \textbf{Mesurer le taux d'expansion} : choisir une anomalie repère, mesurer sa distance $d$ à l'axe de la dorsale, et diviser par son âge $t$ : $v = d/t$ (vitesse de déplacement d'une plaque par rapport à la dorsale).
\item \textbf{En déduire la vitesse d'écartement des continents} : les deux continents bordant les deux plaques s'éloignent l'un de l'autre à la vitesse $2v$ (chaque plaque s'écartant de $v$ de part et d'autre de l'axe).
\item \textbf{Conclure} en reliant au mouvement global des plaques : expansion à la dorsale, dérive des continents embarqués, subduction éventuelle aux frontières convergentes.
\end{enumerate}
\end{methode}

\begin{exoresolu}[Vitesse de divergence entre l'Afrique et l'Amérique du Sud]
Énoncé : sur un profil perpendiculaire à la dorsale médio-atlantique, l'anomalie magnétique datée de 20 millions d'années (Ma) se situe à \SI{400}{km} de l'axe de la dorsale, de chaque côté. 1) Calculer la vitesse de déplacement de la plaque africaine par rapport à la dorsale. 2) En déduire la vitesse d'écartement entre l'Afrique et l'Amérique du Sud. 3) Expliquer pourquoi ces deux continents s'éloignent alors qu'aucun moteur ne les pousse directement.
\tcblower
Solution détaillée :
\begin{enumerate}
\item Vitesse d'une plaque par rapport à la dorsale :
\[ v = \frac{d}{t} = \frac{\SI{400}{km}}{\SI{20}{Ma}} = \SI{20}{km/Ma} = \SI{2}{cm/an}. \]
(Conversion : \SI{1}{km/Ma} = \SI{0,1}{cm/an}, car \SI{1}{km} = $10^5$ cm et 1 Ma = $10^6$ ans.)
\item Les deux plaques s'écartent de la dorsale à la même vitesse, en sens opposés ; la vitesse d'écartement Afrique--Amérique du Sud est donc :
\[ 2v = 2 \times 2 = \SI{4}{cm/an}. \]
\item Les deux continents ne sont pas poussés directement : chacun est \textbf{solidaire de sa plaque lithosphérique}. L'expansion océanique à la dorsale crée de la nouvelle lithosphère qui s'accolde aux bords des deux plaques et les écarte ; les continents, embarqués passivement sur leurs plaques respectives, dérivent avec elles. C'est le mouvement des plaques — dont l'expansion océanique est l'expression à la dorsale — qui entraîne la dérive des continents.
\end{enumerate}
\end{exoresolu}

\begin{exercice}[Pour t'entraîner]
Sur une carte des âges du Pacifique Nord, on relève des basaltes âgés de 40 Ma à \SI{1600}{km} de l'axe de la dorsale du Pacifique Est. 1) Calcule la vitesse d'expansion de chaque flanc. 2) Compare avec le taux atlantique (2 cm/an) et propose une explication à la présence d'une vaste zone de subduction tout autour du Pacifique.
\end{exercice}

\begin{corrige}[Exercice]
1) $v = d/t = 1600/40 = \SI{40}{km/Ma} = \SI{4}{cm/an}$ par flanc. 2) L'expansion pacifique est deux fois plus rapide que l'expansion atlantique. La lithosphère pacifique produite en grande quantité est consommée dans les nombreuses fosses de subduction qui bordent l'océan (Pérou-Chili, Japon, etc.) : le bilan création/destruction reste équilibré et la surface du globe demeure constante.
\end{corrige}

\begin{aretenir}[L'essentiel de la section]
\begin{itemize}
\item La lithosphère rigide est découpée en \cle{plaques} qui glissent sur l'asthénosphère ductile.
\item Les dorsales sont des \cle{frontières divergentes} : l'expansion océanique y crée la lithosphère et écarte les plaques.
\item Les continents, solidaires des plaques, dérivent avec elles (dérive des continents) — ils ne « flottent » pas sur l'océan.
\item La surface créée aux dorsales est consommée par \cle{subduction} dans les fosses : la Terre garde une surface constante.
\item Les anomalies magnétiques et les âges des fonds sont les marqueurs qui permettent de mesurer et dater ces mouvements, aujourd'hui confirmés par le GPS.
\end{itemize}
\end{aretenir}

\newpage

\coursec{Activité d'intégration}

\coursec{NCT10 : Reconstitution chronologique de l'histoire des fonds océaniques}

\courssub{Objectifs de l'activité}

À l'issue de cette activité d'intégration, tu devras être capable de :

\begin{itemize}
\item exploiter une carte (ou un profil) des âges des fonds océaniques pour reconstituer une chronologie d'ouverture océanique ;
\item repérer l'axe d'une dorsale océanique et apparier les anomalies magnétiques symétriques de part et d'autre de cet axe ;
\item attribuer un âge à une anomalie magnétique à l'aide d'une échelle des inversions du champ magnétique terrestre ;
\item tracer et exploiter un graphique âge = f(distance à l'axe de la dorsale) pour calculer un taux d'expansion océanique ;
\item relier l'expansion des fonds océaniques au mouvement des plaques lithosphériques et rédiger une conclusion argumentée, à la manière d'un rapport de mission scientifique.
\end{itemize}

\courssub{Situation}

\textbf{Mission océanographique au large de Kribi : dater l'ouverture de l'Atlantique sud.}

Dans le cadre d'un programme de recherche sur les ressources du plateau continental camerounais, un navire océanographique affrété par l'Institut de Recherches Géologiques et Minières (IRGM) quitte le port de Kribi pour une campagne d'étude de l'Atlantique sud. L'océan Atlantique sud sépare aujourd'hui l'Amérique du Sud (Brésil) de l'Afrique (golfe de Guinée). Les géologues savent depuis longtemps que ces deux continents étaient autrefois soudés : les côtes du Brésil et du golfe de Guinée s'emboîtent comme les pièces d'un puzzle, et l'on retrouve des roches et des fossiles identiques de part et d'autre de l'océan. Mais \textbf{quand cet océan s'est-il ouvert, et à quelle vitesse s'est-il élargi ?}

L'équipe embarquée dispose des données suivantes, recueillies le long d'un profil Est-Ouest joignant le Brésil au golfe de Guinée et coupant la dorsale médio-atlantique.

\textbf{Document 1 : profil magnétique simplifié de l'Atlantique sud.}

Le magnétomètre remorqué par le navire a enregistré les anomalies magnétiques du plancher océanique. De part et d'autre de l'axe de la dorsale, on observe des bandes d'anomalies alternativement positives (roches aimantées dans le sens du champ actuel) et négatives (roches aimantées en sens inverse), numérotées de 1 (la plus proche de l'axe) à 5 (la plus éloignée). La disposition est \textbf{symétrique} par rapport à l'axe de la dorsale.

\begin{popfigure}
\begin{tikzpicture}[scale=1]
% fond
\fill[popblueL] (-6.5,-0.9) rectangle (6.5,0.9);
% anomalies côté ouest
\foreach \x/\n/\c in {-0.9/1/poporange, -2.0/2/popgreen, -3.1/3/poporange, -4.3/4/popgreen, -5.4/5/poporange}{
  \fill[\c] (\x-0.45,-0.9) rectangle (\x+0.45,0.9);
  \node at (\x,1.15) {\small \n};
}
% anomalies côté est
\foreach \x/\n/\c in {0.9/1/poporange, 2.0/2/popgreen, 3.1/3/poporange, 4.3/4/popgreen, 5.4/5/poporange}{
  \fill[\c] (\x-0.45,-0.9) rectangle (\x+0.45,0.9);
  \node at (\x,1.15) {\small \n};
}
% axe de la dorsale
\draw[popdark, very thick, dashed] (0,-1.1) -- (0,1.3);
\node[popdark] at (0,1.65) {\small Axe de la dorsale médio-atlantique};
% flèches d'expansion
\draw[-{Stealth}, popink, very thick] (-0.2,-1.25) -- (-2.2,-1.25);
\draw[-{Stealth}, popink, very thick] (0.2,-1.25) -- (2.2,-1.25);
\node[popink] at (-3.1,-1.25) {\small Expansion};
\node[popink] at (3.1,-1.25) {\small Expansion};
% continents
\node at (-6.1,0) {\small \textbf{Brésil}};
\node at (6.1,0) {\small \textbf{Afrique}};
\end{tikzpicture}
\legende{Profil magnétique simplifié de l'Atlantique sud. Les bandes colorées représentent les anomalies magnétiques positives (orange) et négatives (vert), numérotées de 1 à 5 de part et d'autre de l'axe de la dorsale.}
\end{popfigure}

\textbf{Document 2 : échelle des inversions du champ magnétique terrestre (extraits).}

\begin{center}
\begin{tabularx}{\linewidth}{|c|X|}
\hline
\rowcolor{popblueL} \textbf{Anomalie n°} & \textbf{Âge de la roche (en millions d'années, Ma)} \\
\hline
1 & 10 Ma \\
\hline
2 & 30 Ma \\
\hline
3 & 50 Ma \\
\hline
4 & 70 Ma \\
\hline
5 & 90 Ma \\
\hline
\end{tabularx}
\end{center}

\textbf{Document 3 : résultats de huit forages.}

Des carottes de sédiments ont été prélevées au contact des basaltes du plancher océanique. L'âge des sédiments les plus anciens (donc l'âge du plancher sous-jacent) a été déterminé, ainsi que la distance de chaque forage à l'axe de la dorsale.

\begin{center}
\begin{tabularx}{\linewidth}{|c|X|X|}
\hline
\rowcolor{popblueL} \textbf{Forage} & \textbf{Distance à l'axe (km)} & \textbf{Âge du plancher (Ma)} \\
\hline
F1 & 250 & 10 \\
\hline
F2 & 250 & 10 \\
\hline
F3 & 750 & 30 \\
\hline
F4 & 750 & 30 \\
\hline
F5 & 1250 & 50 \\
\hline
F6 & 1750 & 70 \\
\hline
F7 & 2250 & 90 \\
\hline
F8 & 2250 & 90 \\
\hline
\end{tabularx}
\end{center}

\textbf{Donnée complémentaire :} la largeur actuelle de l'Atlantique sud, mesurée du Brésil au golfe de Guinée au niveau du profil étudié, est d'environ \SI{5000}{km} (soit \SI{2500}{km} de part et d'autre de l'axe de la dorsale).

\courssub{Consignes}

Par groupes de quatre élèves, et en rédigeant chaque étape sur ta feuille de mission :

\begin{enumerate}
\item À l'aide du document 1, repère la position de l'axe de la dorsale médio-atlantique. Justifie ton choix à partir de la disposition des anomalies magnétiques.
\item Apparie les anomalies symétriques de part et d'autre de l'axe. À l'aide du document 2, attribue un âge à chaque couple d'anomalies. Que constates-tu concernant l'âge des roches lorsqu'on s'éloigne de l'axe ?
\item Vérifie la cohérence de ces résultats avec les données des forages (document 3) : compare, pour chaque distance, l'âge des forages situés à l'est et à l'ouest de l'axe.
\item Trace le graphique représentant l'âge du plancher océanique (ordonnée, en Ma) en fonction de la distance à l'axe de la dorsale (abscisse, en km). Que remarques-tu sur l'allure de la courbe ?
\item Calcule la pente de la droite obtenue et déduis-en le \textbf{taux d'expansion} de l'Atlantique sud, exprimé en cm/an, pour un seul côté de la dorsale (taux d'expansion unilatéral), puis le taux d'expansion total (bilatéral).
\item En admettant que ce taux soit resté constant depuis l'ouverture, calcule la date du début d'ouverture de l'Atlantique sud à partir de sa largeur actuelle (\SI{5000}{km}).
\item Rédige un court rapport de mission (10 à 15 lignes) reconstituant la chronologie complète de l'ouverture de l'Atlantique sud : situation initiale, mécanisme d'expansion, vitesse, date de début d'ouverture, et lien avec la tectonique des plaques.
\end{enumerate}

\courssub{Ressources à mobiliser}

\begin{aretenir}[Boîte à outils de la mission]
\begin{itemize}
\item Au niveau d'une \cle{dorsale}, du magma basaltique s'épanche et forme en permanence un nouveau plancher océanique : c'est l'\cle{expansion océanique}. Les roches les plus jeunes sont donc à l'axe, les plus anciennes au loin.
\item En se refroidissant, les basaltes enregistrent l'orientation du champ magnétique terrestre du moment (aimantation thermorémanente). Comme le champ s'est inversé de nombreuses fois, le plancher porte des \cle{anomalies magnétiques} symétriques par rapport à l'axe : c'est une véritable \og bande enregistreuse \fg{}.
\item Relation fondamentale : $\text{vitesse} = \dfrac{\text{distance}}{\text{temps}}$, soit ici $v = \dfrac{d}{t}$ avec $d$ la distance à l'axe et $t$ l'âge de la roche. Sur le graphique âge = f(distance), la pente de la droite donne directement le taux d'expansion.
\item Conversions utiles : $1~\text{km} = 10^{5}~\text{cm}$ ; $1~\text{Ma} = 10^{6}$ ans.
\item L'expansion océanique est la conséquence de l'écartement de deux \cle{plaques lithosphériques} : c'est la base de la tectonique des plaques.
\end{itemize}
\end{aretenir}

\courssub{Démarche et corrigé commenté}

\begin{exoresolu}[Rapport de mission : corrigé détaillé]
Reconstituer la chronologie d'ouverture de l'Atlantique sud à partir des documents 1, 2 et 3.
\tcblower
\textbf{Étape 1 : repérage de l'axe de la dorsale.}

Sur le document 1, les anomalies magnétiques se répètent de manière symétrique de part et d'autre d'une ligne centrale : l'anomalie 1 encadre directement cette ligne, puis on retrouve 2, 3, 4, 5 dans le même ordre vers l'ouest (Brésil) et vers l'est (Afrique). Cette ligne de symétrie est l'\textbf{axe de la dorsale médio-atlantique}, zone où se forme actuellement le plancher le plus jeune.

\textbf{Étape 2 : appariement des anomalies et attribution des âges.}

Les anomalies de même numéro, situées à égale distance de l'axe, forment des couples symétriques : (1 ; 1), (2 ; 2), (3 ; 3), (4 ; 4), (5 ; 5). D'après le document 2, le couple 1 a 10 Ma, le couple 2 a 30 Ma, le couple 3 a 50 Ma, le couple 4 a 70 Ma et le couple 5 a 90 Ma. \textbf{Constat :} plus on s'éloigne de l'axe, plus les roches sont anciennes. C'est la signature de l'expansion océanique : le plancher se forme à l'axe puis s'écarte progressivement de part et d'autre.

\textbf{Étape 3 : vérification par les forages.}

Les forages F1 et F2, situés de part et d'autre de l'axe à 250 km, donnent le même âge (10 Ma) ; de même F3 et F4 (750 km, 30 Ma) et F7 et F8 (2250 km, 90 Ma). La symétrie des âges confirmée par les forages valide l'interprétation magnétique : les deux marqueurs chronologiques (âge des sédiments et anomalies magnétiques) concordent parfaitement.

\textbf{Étape 4 : graphique âge = f(distance à l'axe).}

En reportant les points (250 km ; 10 Ma), (750 km ; 30 Ma), (1250 km ; 50 Ma), (1750 km ; 70 Ma), (2250 km ; 90 Ma), on obtient des points \textbf{alignés sur une droite passant par l'origine} : l'âge est proportionnel à la distance à l'axe. Cela signifie que l'expansion s'est faite à vitesse constante.

\textbf{Étape 5 : calcul du taux d'expansion.}

La pente de la droite se calcule avec deux points, par exemple (250 km ; 10 Ma) et (2250 km ; 90 Ma) :
\begin{align*}
v &= \dfrac{\Delta d}{\Delta t} = \dfrac{2250 - 250~\text{km}}{90 - 10~\text{Ma}} = \dfrac{2000~\text{km}}{80~\text{Ma}} = 25~\text{km/Ma}
\end{align*}
Conversion en cm/an :
\begin{align*}
v &= \dfrac{25 \times 10^{5}~\text{cm}}{10^{6}~\text{ans}} = 2{,}5~\text{cm/an}
\end{align*}
C'est le taux d'expansion \textbf{unilatéral} (d'un seul côté de la dorsale). Le taux \textbf{bilatéral} (écartement total entre les deux continents) est donc :
\[ v_{\text{total}} = 2 \times 2{,}5 = 5~\text{cm/an} \]

\textbf{Étape 6 : date du début d'ouverture.}

L'Atlantique sud mesure aujourd'hui environ \SI{5000}{km}. À vitesse constante de 25 km/Ma de chaque côté, chaque moitié de l'océan (\SI{2500}{km}) s'est formée en :
\[ t = \dfrac{d}{v} = \dfrac{2500~\text{km}}{25~\text{km/Ma}} = 100~\text{Ma} \]
L'ouverture de l'Atlantique sud a donc débuté il y a environ \textbf{100 millions d'années} (au Crétacé), lorsque l'Amérique du Sud et l'Afrique ont commencé à se séparer. Vérification cohérente : le forage le plus ancien (90 Ma) est proche des continents, et aucun plancher plus vieux que 100 Ma n'existe sur ce profil.

\textbf{Étape 7 : rapport de mission (modèle de rédaction).}

\og Il y a environ 100 Ma, l'Amérique du Sud et l'Afrique formaient un seul continent. L'ouverture d'une dorsale au milieu de ce continent a permis l'épanchement de basaltes qui ont formé un nouveau plancher océanique : l'Atlantique sud était né. Depuis, l'expansion océanique s'est poursuivie à un rythme constant de 2,5 cm/an de chaque côté de la dorsale médio-atlantique, soit 5 cm/an d'écartement total entre les deux continents. Les basaltes, en se refroidissant, ont enregistré les inversions successives du champ magnétique terrestre, produisant des anomalies symétriques qui permettent de dater le plancher. Cette expansion traduit l'écartement de deux plaques lithosphériques (plaque sud-américaine et plaque africaine) : elle illustre la tectonique des plaques. \fg
\end{exoresolu}

\begin{attention}[Erreurs fréquentes à éviter]
\begin{itemize}
\item \textbf{Oublier le facteur 2} : la largeur totale de l'océan correspond à une expansion dans les deux sens. Si tu calcules le temps avec \SI{5000}{km} et le taux unilatéral, tu trouves 200 Ma, ce qui est faux : il faut soit diviser la largeur par 2, soit utiliser le taux bilatéral.
\item \textbf{Erreurs de conversion} : 25 km/Ma ne fait pas 25 cm/an ! Vérifie toujours tes puissances de 10 : $1~\text{km/Ma} = 0{,}1~\text{cm/an}$.
\item \textbf{Inverser le sens du raisonnement} : les roches les plus anciennes sont \emph{loin} de l'axe, pas à l'axe. À l'axe, le plancher est en formation : son âge est nul.
\end{itemize}
\end{attention}

\courssub{Bilan de l'activité}

Cette mission océanographique fictive mais réaliste t'a permis de mobiliser simultanément les trois savoir-faire de la leçon :

\begin{itemize}
\item l'\textbf{exploitation d'une carte des âges} : la symétrie des âges de part et d'autre de la dorsale et leur augmentation régulière avec la distance à l'axe constituent la preuve de l'expansion océanique ;
\item l'\textbf{interprétation des anomalies magnétiques} : la \og bande enregistreuse \fg{} des inversions du champ magnétique fournit un calendrier indépendant qui confirme les datations des forages ;
\item le \textbf{calcul d'un taux d'expansion} : la proportionnalité entre âge et distance à l'axe (droite passant par l'origine) montre que l'expansion est un processus régulier, mesurable, de l'ordre de quelques centimètres par an.
\end{itemize}

Tu as ainsi reconstitué une chronologie complète : séparation de l'Afrique et de l'Amérique du Sud il y a environ 100 Ma, formation continue du plancher à la dorsale médio-atlantique, élargissement progressif de l'Atlantique sud jusqu'à sa largeur actuelle de \SI{5000}{km}. Enfin, en reliant cette expansion à l'écartement des plaques africaine et sud-américaine, tu as posé les fondations de la \cle{tectonique des plaques}, la théorie unificatrice de la géologie moderne — la même qui explique l'activité volcanique du Mont Cameroun, situé sur une autre structure tectonique active du golfe de Guinée.

\newpage

\coursec{Méthodes et exercices résolus}

\courssub{Exploiter une carte des âges des fonds océaniques pour reconstituer une chronologie}

\begin{methode}[Lire et exploiter une carte des âges des fonds océaniques]
\begin{enumerate}
\item \textbf{Repérer la légende} : identifier la dorsale (axe de l'océan) et le code couleur des âges (en millions d'années, Ma). La couleur la plus jeune est toujours au niveau de la dorsale.
\item \textbf{Vérifier la symétrie} : les bandes d'âges doivent être \cle{symétriques} de part et d'autre de la dorsale. Si elles ne le sont pas, vérifier que tu n'as pas confondu une fosse ou un continent avec la dorsale.
\item \textbf{Classer les zones} : ordonner les régions du plus jeune (axe de la dorsale) au plus ancien (bord des continents). C'est la \cle{chronologie relative} de l'ouverture océanique.
\item \textbf{Calculer la vitesse d'expansion} : mesurer la distance $d$ entre l'axe de la dorsale et une bande d'âge $t$, puis appliquer :
\[ v = \dfrac{d}{t} \quad \text{(vitesse d'un seul flanc)} \qquad ; \qquad v_{\text{totale}} = \dfrac{2d}{t} \quad \text{(expansion totale, deux flancs)}. \]
Attention aux conversions : $1\ \text{cm/an} = 10\ \text{km/Ma}$.
\item \textbf{Conclure} : rédiger une phrase qui relie l'âge croissant des fonds à l'expansion océanique depuis la dorsale.
\end{enumerate}
\textbf{Réflexe Probatoire} : toujours préciser si la vitesse calculée est \emph{unilatérale} (un seul flanc) ou \emph{totale} (les deux flancs) ; c'est souvent la moitié des points de la question.
\end{methode}

\begin{exoresolu}[Carte des âges de l'Atlantique Nord]
Sur une carte des âges des fonds de l'océan Atlantique, on relève à la latitude des Açores : la dorsale médio-atlantique est bordée de basaltes d'âge nul ; les basaltes situés à \SI{900}{km} de part et d'autre de l'axe ont un âge de 80 Ma ; les plus anciens fonds, au contact des continents, ont 180 Ma.
\begin{enumerate}
\item Justifier que ces données confirment l'expansion des fonds océaniques.
\item Calculer la vitesse d'expansion totale de l'Atlantique à cette latitude, en cm/an.
\end{enumerate}
\tcblower
\textbf{1. Justification de l'expansion.}

Les données montrent une \cle{augmentation régulière et symétrique} de l'âge des basaltes lorsqu'on s'éloigne de la dorsale : 0 Ma à l'axe, 80 Ma à \SI{900}{km}, 180 Ma au contact des continents. Or le basalte se forme par solidification du magma au niveau de la dorsale. Chaque bande de basalte s'est donc formée à l'axe, puis s'en est éloignée au fil du temps, de manière symétrique des deux côtés. C'est la preuve que le plancher océanique \emph{s'écarte de la dorsale} : il y a \cle{expansion océanique} (ou \emph{sea-floor spreading}).

\textbf{2. Calcul de la vitesse d'expansion totale.}

Un basalte âgé de $t = 80$ Ma se trouve à $d = \SI{900}{km}$ de l'axe. Il a donc parcouru \SI{900}{km} en 80 Ma d'un seul côté de la dorsale.

Vitesse d'un flanc (unilatérale) :
\[ v_{\text{flanc}} = \dfrac{d}{t} = \dfrac{900\ \text{km}}{80\ \text{Ma}} = 11{,}25\ \text{km/Ma}. \]

Conversion : $1\ \text{cm/an} = 10\ \text{km/Ma}$, donc :
\[ v_{\text{flanc}} = \dfrac{11{,}25}{10} = 1{,}125\ \text{cm/an}. \]

L'expansion se fait des \emph{deux} côtés de la dorsale, la vitesse totale est donc :
\[ v_{\text{totale}} = 2 \times v_{\text{flanc}} = 2 \times 1{,}125 = 2{,}25\ \text{cm/an}. \]

\textbf{Conclusion} : à la latitude des Açores, l'océan Atlantique s'élargit à une vitesse totale d'environ \SI{2,25}{cm/an}, soit environ \SI{22,5}{km} par million d'années.
\end{exoresolu}

\courssub{Interpréter des anomalies magnétiques symétriques de part et d'autre d'une dorsale}

\begin{methode}[Interpréter un profil d'anomalies magnétiques]
\begin{enumerate}
\item \textbf{Identifier l'axe} : l'anomalie axiale (au centre du profil) correspond à la dorsale ; elle est d'âge nul et de polarité actuelle (normale, époque de Brunhes).
\item \textbf{Repérer les bandes} : de part et d'autre, on observe des bandes alternées d'anomalies positives (polarité normale) et négatives (polarité inverse).
\item \textbf{Mettre en correspondance les bandes symétriques} : deux bandes symétriques par rapport à l'axe ont le \emph{même âge} et la \emph{même polarité} : elles se sont formées en même temps à la dorsale puis ont migré en sens opposés.
\item \textbf{Dater les bandes} : utiliser l'échelle des inversions du champ magnétique terrestre (Brunhes : 0--0,78 Ma, normale ; Matuyama : 0,78--2,58 Ma, inverse, avec l'épisode normal de Jaramillo vers 0,99--1,07 Ma ; Gauss : époque normale de 2,58 à 3,58 Ma).
\item \textbf{Calculer une vitesse} : $v = d/t$ avec $d$ = distance d'une bande datée à l'axe. Comparer les vitesses des deux flancs : si elles sont égales, l'expansion est symétrique.
\item \textbf{Conclure} : les anomalies magnétiques constituent une \cle{bande enregistreuse} qui confirme l'expansion océanique et permet de la dater.
\end{enumerate}
\end{methode}

\begin{exoresolu}[Profil magnétique au sud de l'Islande]
Le profil ci-dessous donne les anomalies magnétiques enregistrées de part et d'autre de la dorsale Reykjanes (au sud de l'Islande). Les bandes noires correspondent aux anomalies positives (polarité normale).

\begin{popfigure}
\begin{center}
\begin{tikzpicture}[x=0.09cm,y=0.35cm]
\draw[->] (-60,0) -- (60,0) node[below left] {\small distance à l'axe (km)};
\draw[->] (0,-1) -- (0,6) node[above] {\small anomalie};
\foreach \x/\h in {0/5, 8/3.5, -8/3.5, 18/2.5, -18/2.5, 30/3, -30/3, 45/2, -45/2}{
  \fill[popblue] (\x-2,0) rectangle (\x+2,\h);
}
\foreach \x in {-45,-30,-18,-8,0,8,18,30,45}{
  \draw (\x,0) node[below]{\scriptsize \x};
}
\draw[poporange,thick,dashed] (0,0) -- (0,5.8);
\node[poporange] at (16,5.4) {\scriptsize axe de la dorsale};
\end{tikzpicture}
\end{center}
\legende{Profil schématique des anomalies magnétiques de part et d'autre de la dorsale Reykjanes. Les bandes bleues (anomalies positives, polarité normale) sont symétriques par rapport à l'axe de la dorsale (trait discontinu orange).}
\end{popfigure}

La bande positive située à \SI{30}{km} de l'axe correspond à l'épisode normal de Jaramillo, daté d'environ 1 Ma.
\begin{enumerate}
\item Montrer que ce profil apporte un argument en faveur de l'expansion océanique.
\item Calculer la vitesse d'expansion d'un flanc de cette dorsale, en cm/an.
\end{enumerate}
\tcblower
\textbf{1. Argument en faveur de l'expansion.}

Le profil montre des bandes d'anomalies \cle{symétriques} par rapport à l'axe de la dorsale : à chaque bande positive située à une distance $d$ d'un côté correspond une bande positive identique à la même distance de l'autre côté (par exemple les bandes à $+30$ km et $-30$ km, ou à $+45$ km et $-45$ km).

Interprétation : lorsque le basalte se solidifie à la dorsale et refroidit en dessous de la température de Curie, sa magnétite s'oriente dans le champ magnétique terrestre de l'époque : c'est l'\cle{aimantation thermorémanente}. Le fond océanique enregistre donc les \cle{inversions du champ magnétique} comme une bande enregistreuse. Le fait que les enregistrements soient identiques et symétriques des deux côtés prouve que les deux moitiés du plancher se sont formées ensemble à l'axe puis se sont \emph{écartées en sens opposés} : c'est l'expansion océanique.

\textbf{2. Vitesse d'expansion d'un flanc.}

La bande de Jaramillo (âge $t = 1$ Ma) se trouve à $d = \SI{30}{km}$ de l'axe :
\[ v = \dfrac{d}{t} = \dfrac{30\ \text{km}}{1\ \text{Ma}} = 30\ \text{km/Ma}. \]
Conversion en cm/an : $30\ \text{km/Ma} = \dfrac{30}{10} = 3\ \text{cm/an}$.

\textbf{Conclusion} : chaque flanc de la dorsale Reykjanes s'écarte de l'axe à environ \SI{3}{cm/an}, soit une expansion totale d'environ \SI{6}{cm/an}, valeur élevée caractéristique de cette région islandaise.
\end{exoresolu}

\courssub{Croiser les marqueurs pour reconstituer une chronologie}

\begin{methode}[Reconstituer la chronologie de l'ouverture d'un océan]
\begin{enumerate}
\item \textbf{Inventorier les marqueurs disponibles} : âges isotopiques des basaltes, anomalies magnétiques datées, âge des premiers sédiments déposés sur le basalte, fossiles contenus dans ces sédiments.
\item \textbf{Vérifier la cohérence} : l'âge du basalte et l'âge de l'anomalie magnétique au même point doivent être concordants ; les sédiments directement au contact du basalte doivent être \emph{légèrement plus jeunes} que ce basalte (ils se déposent après sa formation).
\item \textbf{Établir la chronologie} : dater le début de l'ouverture (âge des fonds les plus anciens, au contact des continents) puis les étapes de l'élargissement.
\item \textbf{Calculer les vitesses} à différentes époques pour détecter d'éventuelles variations (accélération ou ralentissement de l'expansion).
\item \textbf{Rédiger la reconstitution} sous forme d'un récit chronologique daté, du plus ancien au plus récent.
\end{enumerate}
\textbf{Réflexe} : un écart entre âge du basalte et âge du sédiment de base n'est pas une erreur : il traduit le temps écoulé entre la formation du plancher et le début de l'accumulation sédimentaire (la sédimentation est quasi nulle au voisinage de la dorsale, zone jeune et active).
\end{methode}

\begin{exoresolu}[Chronologie de l'ouverture de l'Atlantique Sud]
On dispose des données suivantes pour l'Atlantique Sud, à la latitude de Rio de Janeiro -- Walvis Bay :

\begin{center}
\begin{tabularx}{\linewidth}{|l|X|X|}
\hline
\rowcolor{popblueL} \textbf{Marqueur} & \textbf{Position} & \textbf{Âge} \\
\hline
Basalte au contact des continents & bordures de l'océan & 130 Ma \\
\hline
Anomalie magnétique M0 & \SI{1200}{km} de l'axe & 125 Ma \\
\hline
Plus anciens sédiments sur le basalte & bordures de l'océan & 120 Ma \\
\hline
Basalte de l'axe de la dorsale & axe & 0 Ma \\
\hline
\end{tabularx}
\end{center}

\begin{enumerate}
\item À l'aide de l'ensemble de ces données, reconstituer l'histoire de l'Atlantique Sud.
\item Calculer la vitesse moyenne d'expansion d'un flanc depuis 125 Ma.
\end{enumerate}
\tcblower
\textbf{1. Reconstitution de l'histoire de l'Atlantique Sud.}

\textbf{Cohérence des marqueurs} : au même endroit (bordures de l'océan), le basalte est daté à 130 Ma par datation isotopique et à 125 Ma par l'anomalie magnétique M0 : les deux méthodes donnent des âges proches et concordants. Les sédiments les plus anciens (120 Ma) sont légèrement plus jeunes que le basalte sur lequel ils reposent, ce qui est logique : ils se sont déposés \emph{après} la formation du plancher.

\textbf{Chronologie} :
\begin{itemize}
\item \textbf{Avant 130 Ma} : l'Amérique du Sud et l'Afrique forment un seul continent (Gondwana) ; il n'y a pas d'océan Atlantique Sud.
\item \textbf{Vers 130--125 Ma} (Crétacé inférieur) : rupture continentale, épanchement des premiers basaltes et naissance d'une dorsale : c'est le \cle{début de l'ouverture} de l'Atlantique Sud.
\item \textbf{De 125 Ma à aujourd'hui} : expansion océanique continue et symétrique à partir de la dorsale médio-atlantique ; les fonds s'éloignent progressivement de l'axe pendant que les sédiments s'accumulent sur le basalte vieillissant.
\item \textbf{Aujourd'hui} : l'océan atteint environ \SI{2400}{km} de large à cette latitude ($2 \times \SI{1200}{km}$ pour les fonds de 125 Ma, davantage en comptant les marges) et continue de s'élargir.
\end{itemize}

\textbf{2. Vitesse moyenne d'un flanc.}

Les fonds âgés de $t = 125$ Ma sont à $d = \SI{1200}{km}$ de l'axe :
\[ v = \dfrac{d}{t} = \dfrac{1200\ \text{km}}{125\ \text{Ma}} = 9{,}6\ \text{km/Ma} = \dfrac{9{,}6}{10} = 0{,}96\ \text{cm/an} \approx 1\ \text{cm/an}. \]

\textbf{Conclusion} : depuis sa naissance il y a environ 130 Ma, l'Atlantique Sud s'élargit à une vitesse moyenne d'environ \SI{1}{cm/an} par flanc, soit \SI{2}{cm/an} d'expansion totale : l'Amérique du Sud et l'Afrique, jadis soudées, s'éloignent l'une de l'autre d'environ \SI{2}{cm} chaque année.
\end{exoresolu}

\courssub{Relier l'expansion océanique au mouvement des plaques}

\begin{methode}[Passer de l'expansion océanique à la tectonique des plaques]
\begin{enumerate}
\item \textbf{Définir la plaque} : une \cle{plaque lithosphérique} est un fragment rigide de lithosphère (croûte + manteau supérieur rigide, environ 100 km d'épaisseur) qui se déplace sur l'asthénosphère ductile.
\item \textbf{Identifier les limites de plaques} : dorsales (limites \emph{divergentes}), fosses de subduction et chaînes de collision (limites \emph{convergentes}), failles transformantes (limites \emph{coulissantes}).
\item \textbf{Relier expansion et divergence} : à la dorsale, l'accrétion de nouveau plancher océanique s'accompagne de l'écartement des deux plaques : l'expansion océanique \emph{est} l'expression du mouvement relatif de divergence de deux plaques.
\item \textbf{Appliquer le principe de conservation} : la surface terrestre étant constante, la création de lithosphère aux dorsales est compensée par sa disparition dans les zones de \cle{subduction}.
\item \textbf{Conclure} : la vitesse d'expansion totale mesurée à la dorsale donne directement la vitesse relative de divergence des deux plaques.
\end{enumerate}
\end{methode}

\begin{exoresolu}[De la dorsale à la dérive des continents]
La dorsale médio-atlantique sépare la plaque sud-américaine de la plaque africaine. L'expansion totale mesurée à la latitude de Dakar est de \SI{2,5}{cm/an}.
\begin{enumerate}
\item Définir ce qu'est une plaque lithosphérique et préciser la nature de la limite de plaque au niveau de cette dorsale.
\item Calculer de quelle distance l'Afrique et l'Amérique du Sud se seront éloignées dans 10 millions d'années si la vitesse reste constante.
\item La surface totale de la Terre restant constante, que doit-on nécessairement observer ailleurs à la surface du globe ?
\end{enumerate}
\tcblower
\textbf{1. Plaque et nature de la limite.}

Une \cle{plaque lithosphérique} est un fragment rigide de la lithosphère (croûte terrestre + partie supérieure rigide du manteau, épaisseur d'environ 100 km) qui se déplace à la surface du globe en glissant sur l'asthénosphère, couche ductile du manteau supérieur.

Au niveau de la dorsale médio-atlantique, la plaque sud-américaine et la plaque africaine \emph{s'écartent l'une de l'autre} : il s'agit d'une limite de plaque \cle{divergente} (ou constructive), car de la nouvelle lithosphère océanique y est créée en permanence par solidification du magma.

\textbf{2. Distance d'éloignement dans 10 Ma.}

La vitesse d'expansion totale $v = \SI{2,5}{cm/an}$ est précisément la vitesse relative d'éloignement des deux plaques. Sur une durée $t = 10$ Ma $= 10^{7}$ ans :
\[ d = v \times t = 2{,}5\ \text{cm/an} \times 10^{7}\ \text{ans} = 2{,}5 \times 10^{7}\ \text{cm}. \]
Conversion : $2{,}5 \times 10^{7}\ \text{cm} = 2{,}5 \times 10^{5}\ \text{m} = \SI{250}{km}$.

\textbf{3. Conséquence à l'échelle du globe.}

La création permanente de lithosphère océanique aux dorsales augmenterait la surface de la Terre si rien ne la compensait. Or le rayon terrestre est constant : il doit donc exister des zones où la lithosphère \emph{disparaît} en quantité équivalente. Ce sont les zones de \cle{subduction} (fosses océaniques, par exemple la fosse du Pérou--Chili), où une plaque plonge dans le manteau et y est recyclée.

\textbf{Conclusion} : l'expansion de 2,5 cm/an de l'Atlantique traduit la divergence des plaques africaine et sud-américaine, qui s'éloigneront de \SI{250}{km} en 10 Ma ; cette création de lithosphère est compensée ailleurs par la subduction, principe fondamental de la \cle{tectonique des plaques}.
\end{exoresolu}

\courssub{Construire et exploiter un graphique âge--distance}

\begin{methode}[Tracer et exploiter le graphique âge = f(distance à l'axe)]
\begin{enumerate}
\item \textbf{Placer les axes} : abscisse = distance à l'axe de la dorsale (km), ordonnée = âge du fond (Ma). Placer l'origine (0 km ; 0 Ma) au niveau de la dorsale.
\item \textbf{Reporter les points} mesurés sur un flanc (ou les deux flancs en valeurs absolues).
\item \textbf{Tracer la droite moyenne} passant par l'origine : si les points sont alignés, l'expansion s'est faite à \emph{vitesse constante}.
\item \textbf{Exploiter la pente} : la pente de la droite âge = $f$(distance) vaut $t/d$, donc :
\[ v = \dfrac{1}{\text{pente}} = \dfrac{d}{t}. \]
Une pente qui change brutalement indique un \cle{changement de vitesse} d'expansion au cours du temps.
\item \textbf{Conclure} sur la régularité ou les variations de l'expansion.
\end{enumerate}
\end{methode}

\begin{exoresolu}[Vitesse d'expansion du Pacifique Nord]
On a mesuré, sur le flanc est de la dorsale Est-Pacifique, les âges des fonds en fonction de leur distance à l'axe :

\begin{center}
\begin{tabular}{|l|c|c|c|c|}
\hline
\rowcolor{popblueL} Distance à l'axe (km) & 0 & 200 & 500 & 800 \\
\hline
Âge du fond (Ma) & 0 & 4 & 10 & 16 \\
\hline
\end{tabular}
\end{center}

\begin{enumerate}
\item Tracer le graphique âge en fonction de la distance à l'axe.
\item En déduire la vitesse d'expansion de ce flanc, en cm/an, et commenter la régularité de l'expansion.
\end{enumerate}
\tcblower
\textbf{1. Graphique.}

\begin{popfigure}
\begin{center}
\begin{tikzpicture}
\begin{axis}[
  width=0.75\linewidth, height=6cm,
  xlabel={distance à l'axe (km)}, ylabel={âge du fond (Ma)},
  xmin=0, xmax=850, ymin=0, ymax=18,
  xtick={0,200,400,600,800}, ytick={0,4,8,12,16},
  grid=both, grid style={popline!40},
  axis line style={popdark},
]
\addplot[only marks, mark=*, mark size=2pt, poporange] coordinates {(0,0) (200,4) (500,10) (800,16)};
\addplot[popcurve, domain=0:850, samples=2] {0.02*x};
\end{axis}
\end{tikzpicture}
\end{center}
\legende{Âge des fonds océaniques en fonction de la distance à l'axe de la dorsale Est-Pacifique. Les points sont alignés sur une droite passant par l'origine : l'expansion s'est faite à vitesse constante.}
\end{popfigure}

\textbf{2. Vitesse d'expansion.}

Les quatre points sont \emph{alignés sur une droite passant par l'origine} : l'âge est proportionnel à la distance, donc l'expansion s'est faite à \cle{vitesse constante}.

Calcul de la vitesse à partir du point le plus éloigné (le plus précis) :
\[ v = \dfrac{d}{t} = \dfrac{800\ \text{km}}{16\ \text{Ma}} = 50\ \text{km/Ma} = \dfrac{50}{10} = 5\ \text{cm/an}. \]
Vérification avec les autres points : $\dfrac{200}{4} = 50$ km/Ma ; $\dfrac{500}{10} = 50$ km/Ma. Les trois calculs donnent le même résultat, ce qui confirme la constance de la vitesse.

\textbf{Conclusion} : le flanc est de la dorsale Est-Pacifique s'écarte à \SI{5}{cm/an} (expansion totale d'environ \SI{10}{cm/an}), vitesse constante depuis au moins 16 Ma : les dorsales rapides du Pacifique s'écartent environ trois fois plus vite que la dorsale médio-atlantique.
\end{exoresolu}

\begin{attention}[Les 5 erreurs qui coûtent des points au Probatoire]
\begin{enumerate}
\item \textbf{Confondre vitesse d'un flanc et vitesse totale.} Si la distance mesurée est celle d'une bande \emph{d'un seul côté} de l'axe, $v = d/t$ est la vitesse d'\emph{un flanc} ; l'expansion totale vaut le \emph{double}. Précise toujours laquelle tu calcules.
\item \textbf{Se tromper dans la conversion km/Ma $\leftrightarrow$ cm/an.} Retiens : $1\ \text{cm/an} = 10\ \text{km/Ma}$. Pour passer des km/Ma aux cm/an, on \emph{divise} par 10 (et non l'inverse). Vérifie l'ordre de grandeur : les vitesses de plaques sont de l'ordre de 1 à 10 cm/an, jamais 100 cm/an.
\item \textbf{Affirmer que « la dorsale pousse les continents » sans mécanisme.} Il faut rédiger le mécanisme : le magma solidifié à la dorsale forme du nouveau plancher qui s'écarte symétriquement ; ce sont les \emph{plaques} qui divergent, entraînant les continents solidaires d'elles.
\item \textbf{Oublier la symétrie dans l'interprétation des anomalies magnétiques.} L'argument décisif est la \emph{symétrie} des bandes par rapport à l'axe : une simple alternance de polarités ne prouve rien sans cette symétrie. Cite-la explicitement.
\item \textbf{Dater le début de l'ouverture par les sédiments au lieu du basalte.} Le début de l'ouverture océanique correspond à l'âge du \emph{basalte le plus ancien} (au contact des continents) ; les sédiments de base sont nécessairement un peu plus jeunes. Ne conclus jamais une chronologie sans recouper au moins deux marqueurs (âge isotopique + anomalie magnétique).
\end{enumerate}
\end{attention}

\newpage

\coursec{Exercices}

\courssub{Je vérifie mes connaissances}

\begin{exercice}[QCM : les fonds océaniques]
Pour chaque affirmation, choisis la (ou les) bonne(s) réponse(s) en justifiant brièvement ton choix.
\begin{enumerate}
\item Les fonds océaniques les plus jeunes se trouvent :
\begin{enumerate}
\item près des continents ;
\item au niveau de l'axe des dorsales ;
\item dans les fosses océaniques.
\end{enumerate}
\item Les anomalies magnétiques de part et d'autre d'une dorsale sont :
\begin{enumerate}
\item disposées au hasard ;
\item symétriques par rapport à l'axe de la dorsale ;
\item toutes positives.
\end{enumerate}
\item L'âge maximal des fonds océaniques actuels est d'environ :
\begin{enumerate}
\item 4,6 Ga ;
\item 180 Ma ;
\item 65 Ma.
\end{enumerate}
\item Le taux d'expansion d'un océan s'exprime généralement en :
\begin{enumerate}
\item km/Ma ;
\item cm/an ;
\item Ma/km.
\end{enumerate}
\end{enumerate}
\end{exercice}

\begin{exercice}[Vrai ou faux]
Indique si chaque affirmation est vraie ou fausse, puis justifie ta réponse en une ou deux phrases.
\begin{enumerate}
\item Les roches basaltiques des fonds océaniques enregistrent le champ magnétique terrestre au moment de leur refroidissement.
\item Plus on s'éloigne de l'axe d'une dorsale, plus les fonds océaniques sont jeunes.
\item Le champ magnétique terrestre a toujours eu la même polarité au cours des temps géologiques.
\item La dorsale est-Pacifique s'ouvre plus rapidement que la dorsale médio-atlantique.
\item Les sédiments sont d'autant plus épais que l'on se rapproche de l'axe de la dorsale.
\end{enumerate}
\end{exercice}

\begin{exercice}[Définitions]
Définis précisément, en une ou deux phrases chacune, les termes suivants : \cle{dorsale océanique}, \cle{anomalie magnétique}, \cle{isochrone}, \cle{expansion océanique}, \cle{inversion du champ magnétique terrestre}.
\end{exercice}

\begin{exercice}[Calcul direct]
Un échantillon de basalte prélevé par forage à \SI{600}{km} de l'axe d'une dorsale a un âge de \SI{30}{Ma}.
\begin{enumerate}
\item Calcule le taux d'expansion (demi-vitesse) de ce fond océanique en cm/an.
\item En supposant ce taux constant, à quelle distance de l'axe se trouve un fond âgé de \SI{60}{Ma} ?
\item Quel serait l'âge d'un basalte foré à \SI{900}{km} de l'axe ?
\end{enumerate}
\end{exercice}

\courssub{Je m'entraîne}

\begin{exercice}[Lecture d'une carte des âges]
Sur une carte des âges des fonds de l'océan Atlantique, on relève les distances suivantes entre l'axe de la dorsale médio-atlantique et les isochrones, mesurées perpendiculairement à l'axe :
\begin{center}
\begin{tabularx}{\linewidth}{|l|X|X|X|X|}
\hline
\rowcolor{popblueL} Isochrone (Ma) & 10 & 20 & 40 & 80 \\
\hline
Distance à l'axe (km) & 125 & 250 & 500 & 1000 \\
\hline
\end{tabularx}
\end{center}
\begin{enumerate}
\item Trace le graphique donnant la distance à l'axe en fonction de l'âge des fonds.
\item Que peux-tu conclure de l'allure de la courbe obtenue ?
\item Calcule le taux d'expansion de l'Atlantique à cette latitude, en cm/an.
\end{enumerate}
\end{exercice}

\begin{exercice}[Appariement des anomalies magnétiques]
Un profil magnétique perpendiculaire à une dorsale montre, de part et d'autre de l'axe, cinq bandes d'anomalies notées A, B, C, D, E d'un côté et A', B', C', D', E' de l'autre. L'échelle des inversions magnétiques attribue les âges suivants aux limites de bandes : A : \SI{0}{Ma} ; B : \SI{5}{Ma} ; C : \SI{10}{Ma} ; D : \SI{20}{Ma} ; E : \SI{35}{Ma}.
\begin{enumerate}
\item Apparie les bandes symétriques et justifie ton appariement.
\item La bande D' se situe à \SI{300}{km} de l'axe. Calcule le taux d'expansion.
\item À quelle distance de l'axe se trouve la bande E ?
\item Pourquoi dit-on que le fond océanique fonctionne comme une \og bande enregistreuse \fg{} ?
\end{enumerate}
\end{exercice}

\begin{exercice}[Comparaison de deux dorsales]
On donne les taux d'expansion moyens de deux dorsales : dorsale médio-atlantique : \SI{2}{cm/an} ; dorsale est-Pacifique : \SI{10}{cm/an}.
\begin{enumerate}
\item Calcule, pour chaque dorsale, la distance parcourue par un point du fond océanique en \SI{10}{Ma}.
\item La dorsale médio-atlantique présente un relief accidenté avec un rift axial profond, tandis que la dorsale est-Pacifique présente un profil plus bombé et lisse. Propose une explication reliant la morphologie des dorsales à leur vitesse d'expansion.
\item Calcule le temps nécessaire pour que l'Atlantique s'élargisse de \SI{100}{km} de chaque côté de sa dorsale.
\end{enumerate}
\end{exercice}

\begin{exercice}[Datation d'un basalte foré]
Au cours d'une campagne de forages océaniques, un basalte est prélevé à \SI{900}{km} de l'axe d'une dorsale ; son âge, déterminé par une méthode radiométrique, est de \SI{45}{Ma}.
\begin{enumerate}
\item Calcule le taux d'expansion du fond océanique en cm/an.
\item En supposant ce taux constant depuis l'ouverture de l'océan, détermine la distance à l'axe d'un fond âgé de \SI{90}{Ma}.
\item Un autre forage, réalisé à \SI{1500}{km} de l'axe, atteint des sédiments dont la base repose sur un basalte de \SI{75}{Ma}. Ce résultat est-il cohérent avec le taux calculé à la question 1 ? Justifie par un calcul.
\end{enumerate}
\end{exercice}

\courssub{Je me prépare au Probatoire}

\begin{exercice}[Type Probatoire : reconstitution de l'ouverture de l'Atlantique]
\textit{(D'après un sujet de Probatoire C, session 2019, adapté.)}

Au large du golfe de Guinée, l'océan Atlantique sud sépare l'Afrique de l'Amérique du Sud. Une carte des âges des fonds océaniques fournit les isochrones suivantes, mesurées de part et d'autre de l'axe de la dorsale médio-atlantique, à la latitude du Cameroun :
\begin{center}
\begin{tabularx}{\linewidth}{|l|X|X|X|X|}
\hline
\rowcolor{popblueL} Âge de l'isochrone (Ma) & 0 & 40 & 80 & 120 \\
\hline
Distance à l'axe, côté africain (km) & 0 & 500 & 1000 & 1500 \\
\hline
Distance à l'axe, côté sud-américain (km) & 0 & 500 & 1000 & 1500 \\
\hline
\end{tabularx}
\end{center}
\begin{enumerate}
\item Définis les termes : \cle{dorsale océanique}, \cle{isochrone}.
\item À partir des données du tableau, situe l'axe de la dorsale et montre que la répartition des âges est symétrique de part et d'autre de cet axe.
\item Calcule le taux d'expansion (demi-vitesse) de l'Atlantique sud, en cm/an, puis la vitesse totale d'écartement entre l'Afrique et l'Amérique du Sud.
\item Construis le graphique donnant la largeur totale de l'Atlantique sud en fonction du temps, de \SI{0}{Ma} à \SI{120}{Ma}.
\item Déduis de ce graphique l'époque à laquelle l'Afrique (et donc la région du Cameroun actuel) et l'Amérique du Sud ont commencé à se séparer. À quelle ère géologique cette époque appartient-elle ?
\end{enumerate}
\end{exercice}

\begin{exercice}[Type Probatoire : les anomalies magnétiques, témoins de l'expansion]
\textit{(D'après un sujet de Probatoire TI, session 2021, adapté.)}

Un navire océanographique enregistre, le long d'un profil Est--Ouest de \SI{700}{km} traversant la dorsale médio-atlantique, les variations de l'intensité du champ magnétique mesuré à la surface des fonds. On observe une succession de bandes d'anomalies positives et négatives, symétriques par rapport à l'axe de la dorsale. L'échelle de datation des inversions magnétiques donne : fin de la période normale actuelle (Brunhes) : \SI{0,78}{Ma} ; limite d'une bande ancienne repérée à \SI{350}{km} de l'axe : \SI{20}{Ma}.
\begin{enumerate}
\item Rappelle ce qu'est une anomalie magnétique et explique son origine au niveau d'une dorsale.
\item Schématise, par un profil légendé, la disposition des bandes d'anomalies de part et d'autre de l'axe de la dorsale.
\item Calcule le taux d'expansion de ce fond océanique à partir de la bande située à \SI{350}{km} de l'axe.
\item Calcule la largeur de la bande correspondant à la période normale actuelle (Brunhes), de part et d'autre de l'axe.
\item Explique, en deux arguments tirés de cet exercice, en quoi les anomalies magnétiques constituent une preuve de l'expansion des fonds océaniques.
\end{enumerate}
\end{exercice}

\begin{exercice}[Type Probatoire : synthèse — de l'expansion océanique à la tectonique des plaques]
\textit{(Sujet de synthèse, type épreuve de SVTEEHB au Probatoire.)}

Les forages profonds réalisés dans les océans ont montré que : (1) l'âge des basaltes augmente régulièrement lorsqu'on s'éloigne de l'axe des dorsales ; (2) aucun fond océanique n'a un âge supérieur à environ \SI{180}{Ma}, alors que les continents présentent des roches âgées de plus de \SI{3}{Ga} ; (3) l'épaisseur des sédiments augmente avec la distance à l'axe des dorsales ; (4) les anomalies magnétiques dessinent des bandes symétriques par rapport aux dorsales.
\begin{enumerate}
\item Pour chacune des quatre observations ci-dessus, indique ce qu'elle apporte à la compréhension de la dynamique des fonds océaniques.
\item Explique pourquoi l'absence de fonds océaniques très anciens impose l'existence de zones de disparition de la lithosphère océanique. Comment appelle-t-on ces zones ?
\item En t'appuyant sur les âges des fonds et les anomalies magnétiques, montre que l'existence de fonds océaniques jeunes au niveau des dorsales prouve que les plaques lithosphériques sont en mouvement.
\item Rédige une conclusion de dix à quinze lignes présentant la notion de \cle{tectonique des plaques} : tu y définiras la lithosphère, tu citeras les trois types de limites de plaques et tu illustreras chacune par un exemple.
\end{enumerate}
\end{exercice}

\newpage

\coursec{Corrigés des exercices}

\coursec{Corrigés des exercices — Reconstitution chronologique de l'histoire des fonds océaniques}

\begin{corrige}[Exercice 1 — QCM : les fonds océaniques]
\textbf{Énoncé :} Se reporter à l'exercice 1 du manuel (QCM sur l'âge aux dorsales, symétrie des anomalies, âge maximal, unité de vitesse).
\tcblower
\textbf{Solution détaillée :}
\begin{enumerate}
\item \textbf{Réponse b} : au niveau de l'axe des dorsales. C'est là que le magma basaltique s'épanche et se solidifie en permanence : la croûte océanique vient de se former, elle est donc d'âge nul (ou très récent). En s'éloignant de l'axe, les fonds sont de plus en plus anciens.
\item \textbf{Réponse b} : symétriques par rapport à l'axe de la dorsale. Le basalte enregistre la polarité du champ magnétique au moment de son refroidissement ; comme l'expansion s'effectue des deux côtés de l'axe, chaque bande d'anomalie possède une bande jumelle de même âge de l'autre côté.
\item \textbf{Réponse b} : environ \SI{180}{Ma} (Jurassique). Les fonds plus anciens ont disparu par subduction ; seuls les continents conservent des roches de plusieurs milliards d'années.
\item \textbf{Réponse b} : en cm/an. Il s'agit d'une vitesse (distance/temps) ; l'ordre de grandeur est de quelques centimètres par an. (Remarque : \SI{1}{km/Ma} = \SI{1}{mm/an} = \SI{0,1}{cm/an} ; l'unité usuelle au lycée est le cm/an.)
\end{enumerate}
\end{corrige}

\begin{corrige}[Exercice 2 — Vrai ou faux]
\textbf{Énoncé :} Se reporter à l'exercice 2 du manuel (affirmations sur l'aimantation thermorémanente, âge vs distance, inversions magnétiques, vitesses des dorsales, épaisseur des sédiments).
\tcblower
\textbf{Solution détaillée :}
\begin{enumerate}
\item \textbf{Vrai.} En refroidissant sous la température de Curie (environ \SI{580}{\celsius} pour la magnétite), les minéraux ferromagnétiques du basalte s'aimantent dans le sens du champ magnétique terrestre du moment : c'est l'aimantation thermorémanente.
\item \textbf{Faux.} C'est l'inverse : plus on s'éloigne de l'axe, plus les fonds sont \emph{anciens}, car ils se sont formés à l'axe puis ont été écartés par l'expansion océanique.
\item \textbf{Faux.} Le champ magnétique terrestre s'est inversé de nombreuses fois au cours des temps géologiques ; on connaît ainsi des périodes de polarité normale (comme l'actuelle période de Brunhes) et des périodes de polarité inverse.
\item \textbf{Vrai.} La dorsale est-Pacifique s'ouvre à environ \SI{10}{cm/an} (dorsale rapide), contre environ \SI{2}{cm/an} pour la dorsale médio-atlantique (dorsale lente).
\item \textbf{Faux.} Les sédiments sont d'autant plus \emph{épais} que l'on s'\emph{éloigne} de l'axe : un fond ancien a eu plus de temps pour accumuler des sédiments. À l'axe, le basalte est pratiquement nu.
\end{enumerate}
\end{corrige}

\begin{corrige}[Exercice 3 — Définitions]
\textbf{Énoncé :} Définir les termes suivants : Dorsale océanique, Anomalie magnétique, Isochrone, Expansion océanique, Inversion du champ magnétique terrestre.
\tcblower
\textbf{Solution détaillée :}
\begin{itemize}
\item \cle{Dorsale océanique} : ride montagneuse sous-marine, longue de plusieurs milliers de kilomètres, située au milieu des océans, au niveau de laquelle le magma basaltique s'épanche et forme en permanence une nouvelle croûte océanique.
\item \cle{Anomalie magnétique} : écart entre l'intensité du champ magnétique mesuré localement au-dessus des fonds océaniques et la valeur théorique du champ terrestre ; elle est positive si les roches renforcent le champ actuel (aimantation normale), négative si elles l'atténuent (aimantation inverse).
\item \cle{Isochrone} : ligne joignant, sur une carte, les points des fonds océaniques de même âge.
\item \cle{Expansion océanique} : écartement progressif, de part et d'autre de l'axe d'une dorsale, des fonds océaniques au fur et à mesure de la formation de nouvelle croûte ; elle se mesure en cm/an.
\item \cle{Inversion du champ magnétique terrestre} : renversement de la polarité du champ magnétique (le pôle Nord magnétique prenant la place du pôle Sud et inversement), phénomène survenu de nombreuses fois au cours de l'histoire de la Terre.
\end{itemize}
\end{corrige}

\begin{corrige}[Exercice 4 — Calcul direct]
\textbf{Énoncé :} Un fond océanique s'est écarté de \SI{600}{km} de l'axe de la dorsale en \SI{30}{Ma}. Calculer le taux d'expansion (demi-vitesse) en km/Ma et en cm/an. Quelle distance aura parcouru ce fond dans \SI{60}{Ma} ? À quelle distance se trouve un fond âgé de \SI{45}{Ma} ?
\tcblower
\textbf{Solution détaillée :}
\begin{enumerate}
\item Le taux d'expansion (demi-vitesse) est $v = \dfrac{d}{t} = \dfrac{\SI{600}{km}}{\SI{30}{Ma}} = \SI{20}{km/Ma}$.
Conversion : $\SI{20}{km/Ma} = \dfrac{20 \times 10^{5}\ \text{cm}}{10^{6}\ \text{an}} = \SI{2}{cm/an}$.
\item À taux constant, $d = v \times t = \SI{20}{km/Ma} \times \SI{60}{Ma} = \SI{1200}{km}$.
\item $t = \dfrac{d}{v} = \dfrac{\SI{900}{km}}{\SI{20}{km/Ma}} = \SI{45}{Ma}$. (Vérification : $900 = 20 \times 45$. Correct.)
\end{enumerate}
\end{corrige}

\begin{corrige}[Exercice 5 — Lecture d'une carte des âges]
\textbf{Énoncé :} On donne les couples (âge en Ma ; distance à l'axe en km) : (0;0), (10;125), (20;250), (40;500), (80;1000). Tracer le graphique, interpréter et calculer la vitesse.
\tcblower
\textbf{Solution détaillée :}
\begin{enumerate}
\item Graphique : on place l'âge en abscisse (Ma) et la distance en ordonnée (km). Les points $(0\,;\,0)$, $(10\,;\,125)$, $(20\,;\,250)$, $(40\,;\,500)$, $(80\,;\,1000)$ sont alignés sur une droite passant par l'origine.
\begin{popfigure}
\begin{center}
\begin{tikzpicture}
\begin{axis}[width=0.8\linewidth,height=6cm,xlabel={Âge (Ma)},ylabel={Distance à l'axe (km)},xmin=0,xmax=90,ymin=0,ymax=1100,grid=both,grid style={popline!40}]
\addplot[popcurve,domain=0:80,samples=2]{12.5*x};
\addplot[only marks,mark=*,color=poporange] coordinates {(10,125) (20,250) (40,500) (80,1000)};
\end{axis}
\end{tikzpicture}
\end{center}
\legende{Distance à l'axe de la dorsale en fonction de l'âge des fonds : les points sont alignés sur une droite passant par l'origine.}
\end{popfigure}
\item L'alignement des points sur une droite passant par l'origine montre que la distance est \textbf{proportionnelle à l'âge} : le taux d'expansion est \textbf{constant} depuis au moins \SI{80}{Ma}.
\item $v = \dfrac{\SI{1000}{km}}{\SI{80}{Ma}} = \SI{12,5}{km/Ma} = \dfrac{12{,}5 \times 10^{5}\ \text{cm}}{10^{6}\ \text{an}} = \SI{1,25}{cm/an}$.
\end{enumerate}
\end{corrige}

\begin{corrige}[Exercice 6 — Appariement des anomalies magnétiques]
\textbf{Énoncé :} Sur un profil magnétique symétrique, les bandes A, B, C, D, E (côté Ouest) et A', B', C', D', E' (côté Est) sont observées. La bande A (et A') a \SI{20}{Ma} et se trouve à \SI{300}{km} de l'axe. Apparier les bandes, calculer la vitesse, déterminer la position de la bande E (âge \SI{35}{Ma}) et expliquer l'expression « bande enregistreuse ».
\tcblower
\textbf{Solution détaillée :}
\begin{enumerate}
\item Appariement : A avec A', B avec B', C avec C', D avec D', E avec E'. Justification : chaque bande s'est formée à l'axe à une époque donnée, puis s'est scindée en deux moitiés écartées symétriquement ; deux bandes symétriques ont donc le même âge et la même polarité.
\item $v = \dfrac{d}{t} = \dfrac{\SI{300}{km}}{\SI{20}{Ma}} = \SI{15}{km/Ma} = \SI{1,5}{cm/an}$.
\item $d = v \times t = \SI{15}{km/Ma} \times \SI{35}{Ma} = \SI{525}{km}$. La bande E se trouve à \SI{525}{km} de l'axe (et E' symétriquement de l'autre côté).
\item Le fond océanique se forme à l'axe en enregistrant la polarité du champ magnétique du moment, puis défile de part et d'autre comme la bande d'un magnétophone : la succession chronologique des inversions reste \emph{enregistrée} dans les basaltes, en double exemplaire symétrique. D'où l'expression \og bande enregistreuse \fg{}.
\end{enumerate}
\end{corrige}

\begin{corrige}[Exercice 7 — Comparaison de deux dorsales]
\textbf{Énoncé :} Comparer la dorsale médio-atlantique (vitesse \SI{2}{cm/an}) et la dorsale est-Pacifique (vitesse \SI{10}{cm/an}) sur \SI{10}{Ma}. Expliquer la différence de morphologie (rift axial, relief). Calculer l'âge d'un fond à \SI{100}{km} de l'axe de la dorsale médio-atlantique.
\tcblower
\textbf{Solution détaillée :}
\begin{enumerate}
\item Dorsale médio-atlantique : $d = \SI{2}{cm/an} \times 10^{7}\ \text{an} = 2 \times 10^{7}\ \text{cm} = \SI{200}{km}$. Dorsale est-Pacifique : $d = \SI{10}{cm/an} \times 10^{7}\ \text{an} = 10^{8}\ \text{cm} = \SI{1000}{km}$.
\item Une dorsale \textbf{lente} (Atlantique) laisse le temps à la lithosphère de se refroidir, de s'épaissir et de se fracturer près de l'axe : d'où un rift axial profond et un relief accidenté. Une dorsale \textbf{rapide} (Pacifique) écarte vite les fonds encore chauds et peu épais : le rift est peu marqué et le profil bombé et lisse. La morphologie dépend donc de la vitesse d'expansion.
\item $t = \dfrac{d}{v} = \dfrac{\SI{100}{km}}{\SI{2}{cm/an}} = \dfrac{10^{7}\ \text{cm}}{2\ \text{cm/an}} = 5 \times 10^{6}\ \text{an} = \SI{5}{Ma}$.
\end{enumerate}
\end{corrige}

\begin{corrige}[Exercice 8 — Datation d'un basalte foré]
\textbf{Énoncé :} Un basalte foré à \SI{900}{km} de l'axe a \SI{45}{Ma}. Calculer la vitesse. Prédire la distance d'un basalte de \SI{90}{Ma}. Vérifier la cohérence avec un basalte de \SI{75}{Ma} foré à \SI{1500}{km}.
\tcblower
\textbf{Solution détaillée :}
\begin{enumerate}
\item $v = \dfrac{\SI{900}{km}}{\SI{45}{Ma}} = \SI{20}{km/Ma} = \SI{2}{cm/an}$.
\item $d = v \times t = \SI{20}{km/Ma} \times \SI{90}{Ma} = \SI{1800}{km}$.
\item Âge attendu à \SI{1500}{km} : $t = \dfrac{\SI{1500}{km}}{\SI{20}{km/Ma}} = \SI{75}{Ma}$. L'âge mesuré (\SI{75}{Ma}) est exactement celui prévu par le taux calculé : le résultat est donc \textbf{parfaitement cohérent}, ce qui confirme la constance du taux d'expansion.
\end{enumerate}
\end{corrige}

\begin{corrige}[Exercice 9 — Type Probatoire : reconstitution de l'ouverture de l'Atlantique]
\textbf{Énoncé :} Définir dorsale océanique et isochrone. Montrer la symétrie des âges (isochrones 0, 40, 80, 120 Ma à distances 0, 500, 1000, 1500 km de part et d'autre). Calculer la demi-vitesse et la vitesse totale d'écartement Afrique–Amérique du Sud. Tracer la largeur totale en fonction du temps. Dater l'ouverture et situer l'ère géologique.
\tcblower
\textbf{Solution détaillée :}
\begin{enumerate}
\item \cle{Dorsale océanique} : ride volcanique sous-marine au niveau de laquelle se forme en permanence une nouvelle croûte océanique par épanchement de basaltes. \cle{Isochrone} : ligne joignant les points des fonds océaniques de même âge.
\item L'axe de la dorsale correspond à l'isochrone \SI{0}{Ma} (fonds d'âge nul). À âge égal, les distances à l'axe sont identiques des deux côtés (\SI{500}{km} pour \SI{40}{Ma}, \SI{1000}{km} pour \SI{80}{Ma}, \SI{1500}{km} pour \SI{120}{Ma}) : la répartition des âges est donc \textbf{symétrique} par rapport à l'axe.
\item Demi-vitesse : $v = \dfrac{\SI{500}{km}}{\SI{40}{Ma}} = \SI{12,5}{km/Ma} = \SI{1,25}{cm/an}$. Vitesse totale d'écartement Afrique--Amérique du Sud : $2 \times 1,25 = \SI{2,5}{cm/an}$.
\item La largeur totale vaut $L = 2 \times 12{,}5 \times t = 25\,t$ (en km, $t$ en Ma) : \SI{0}{km} à \SI{0}{Ma}, \SI{1000}{km} à \SI{40}{Ma}, \SI{2000}{km} à \SI{80}{Ma}, \SI{3000}{km} à \SI{120}{Ma} — droite passant par l'origine.
\begin{popfigure}
\begin{center}
\begin{tikzpicture}
\begin{axis}[width=0.8\linewidth,height=6cm,xlabel={Temps (Ma)},ylabel={Largeur totale (km)},xmin=0,xmax=130,ymin=0,ymax=3300,grid=both,grid style={popline!40}]
\addplot[popcurve,domain=0:120,samples=2]{25*x};
\addplot[only marks,mark=*,color=poporange] coordinates {(40,1000) (80,2000) (120,3000)};
\end{axis}
\end{tikzpicture}
\end{center}
\legende{Élargissement de l'Atlantique sud en fonction du temps : croissance linéaire de \SI{25}{km/Ma}.}
\end{popfigure}
\item Le graphique passe par l'origine : la largeur s'annule il y a environ \SI{120}{Ma}. C'est donc à cette époque que l'Afrique (et la région du Cameroun actuel) et l'Amérique du Sud ont commencé à se séparer. \SI{120}{Ma} correspond au \textbf{Crétacé}, ère \textbf{Mésozoïque} (ère secondaire).
\end{enumerate}
\textbf{Barème indicatif (sur 10)} : définitions 2 pts ; symétrie justifiée 2 pts ; calculs des vitesses 2 pts ; graphique (axes, échelle, droite) 2 pts ; datation et ère 2 pts.
\textbf{Points de vigilance} : ne pas confondre demi-vitesse et vitesse totale (facteur 2) ; convertir correctement km/Ma en cm/an ($\SI{1}{km/Ma} = \SI{0,1}{cm/an}$) ; citer l'ère géologique, pas seulement la période.
\end{corrige}

\begin{corrige}[Exercice 10 — Type Probatoire : les anomalies magnétiques, témoins de l'expansion]
\textbf{Énoncé :} Définir l'anomalie magnétique et expliquer son origine à la dorsale. Tracer le profil d'anomalies symétriques sur \SI{700}{km} (axe au centre, bande de Brunhes \SI{0,78}{Ma} au centre). Calculer la demi-vitesse si la bande de \SI{20}{Ma} est à \SI{350}{km}. Calculer la largeur totale de l'anomalie de Brunhes. Donner deux arguments prouvant l'expansion.
\tcblower
\textbf{Solution détaillée :}
\begin{enumerate}
\item Une \cle{anomalie magnétique} est l'écart entre le champ magnétique mesuré au-dessus des fonds et la valeur théorique du champ terrestre. À la dorsale, le basalte en refroidissant sous sa température de Curie s'aimante dans le sens du champ de l'époque : polarité normale $\Rightarrow$ anomalie positive ; polarité inverse $\Rightarrow$ anomalie négative.
\item Schéma :
\begin{popfigure}
\begin{center}
\begin{tikzpicture}[x=0.011cm,y=1cm]
\fill[popblueL] (-350,0) rectangle (-262.5,1);
\fill[popcream] (-262.5,0) rectangle (-175,1);
\fill[popblueL] (-175,0) rectangle (-87.5,1);
\fill[popcream] (-87.5,0) rectangle (-13.65,1);
\fill[popgreenL] (-13.65,0) rectangle (13.65,1);
\fill[popcream] (13.65,0) rectangle (87.5,1);
\fill[popblueL] (87.5,0) rectangle (175,1);
\fill[popcream] (175,0) rectangle (262.5,1);
\fill[popblueL] (262.5,0) rectangle (350,1);
\draw[thick] (-350,0) rectangle (350,1);
\draw[->,thick,poporange] (0,1.1) -- (0,1.8) node[above]{axe de la dorsale (\SI{0}{Ma})};
\draw[<->] (-350,-0.25) -- (350,-0.25) node[midway,below]{profil de \SI{700}{km}};
\node at (-306,0.5){\small $-$}; \node at (-219,0.5){\small $+$}; \node at (-131,0.5){\small $-$}; \node at (-50,0.5){\small $+$};
\node at (0,0.5){\small $+$ Brunhes}; \node at (50,0.5){\small $+$}; \node at (131,0.5){\small $-$}; \node at (219,0.5){\small $+$}; \node at (306,0.5){\small $-$};
\end{tikzpicture}
\end{center}
\legende{Bandes d'anomalies positives ($+$) et négatives ($-$) symétriques par rapport à l'axe de la dorsale ; la bande axiale correspond à la période normale actuelle (Brunhes).}
\end{popfigure}
\item $v = \dfrac{\SI{350}{km}}{\SI{20}{Ma}} = \SI{17,5}{km/Ma} = \SI{1,75}{cm/an}$.
\item Largeur de la bande de Brunhes d'un côté : $d = v \times t = \SI{17,5}{km/Ma} \times \SI{0,78}{Ma} = \SI{13,65}{km}$. Largeur totale de part et d'autre : $2 \times 13{,}65 = \SI{27,3}{km}$.
\item Argument 1 : la \textbf{symétrie} des bandes par rapport à l'axe montre que la croûte se forme à l'axe puis s'écarte des deux côtés. Argument 2 : la succession des bandes correspond exactement à l'échelle des inversions magnétiques datée sur les continents, et leur largeur croît avec l'âge selon un taux constant : le fond océanique a bien \emph{défilé} comme une bande enregistreuse, ce qui prouve l'expansion.
\end{enumerate}
\textbf{Barème indicatif (sur 10)} : définition et origine 2 pts ; schéma légendé 2 pts ; calcul du taux 2 pts ; largeur de Brunhes 2 pts ; deux arguments 2 pts.
\textbf{Points de vigilance} : exiger la légende complète du schéma (axe, anomalies $+$/$-$, symétrie) ; pour la question 4, préciser si l'on donne la demi-largeur ou la largeur totale ; les arguments doivent s'appuyer sur les \emph{données de l'exercice}, pas sur des généralités.
\end{corrige}

\begin{corrige}[Exercice 11 — Type Probatoire : synthèse — de l'expansion océanique à la tectonique des plaques]
\textbf{Énoncé :} Analyser 4 observations (âge croissant, âge max 180 Ma, sédiments, symétrie anomalies). En déduire l'existence de zones de disparition (subduction). Prouver que la lithosphère est mobile et découpée en plaques. Rédiger une conclusion définissant la tectonique des plaques, la lithosphère, et les 3 types de limites avec exemples.
\tcblower
\textbf{Solution détaillée :}
\begin{enumerate}
\item (1) L'augmentation régulière de l'âge avec la distance à l'axe montre que la croûte océanique se forme à la dorsale puis s'en écarte : c'est l'expansion océanique. (2) L'absence de fonds plus vieux que \SI{180}{Ma} montre que la lithosphère océanique est recyclée : elle disparaît quelque part, contrairement à la lithosphère continentale, très ancienne et conservée. (3) L'épaisseur croissante des sédiments confirme le vieillissement des fonds avec la distance à l'axe (plus un fond est ancien, plus il a accumulé de sédiments). (4) La symétrie des anomalies magnétiques prouve que l'écartement se fait de manière égale des deux côtés de l'axe et fournit une chronologie de l'expansion.
\item Si de la croûte se crée en permanence aux dorsales alors que la surface de la Terre reste constante et qu'aucun fond n'excède \SI{180}{Ma}, il doit exister des zones où la lithosphère océanique \textbf{disparaît} en plongeant dans le manteau : ce sont les zones de \textbf{subduction} (matérialisées par les fosses océaniques, par exemple la fosse du Pérou--Chili).
\item Les fonds d'âge nul (\SI{0}{Ma}) se trouvent aujourd'hui à l'axe des dorsales, tandis que des fonds de même âge se retrouvent à égale distance de part et d'autre (isochrones et anomalies symétriques). Une même bande de croûte, formée à l'axe, est aujourd'hui séparée en deux moitiés éloignées : les portions de lithosphère ont donc \textbf{bougé}. La lithosphère est ainsi découpée en \textbf{plaques en mouvement}, emportées de part et d'autre des dorsales.
\item \textbf{Conclusion rédigée.} La \cle{tectonique des plaques} est le modèle qui décrit la couche externe rigide de la Terre, la \cle{lithosphère} (croûte + partie supérieure du manteau, épaisse d'environ \SI{100}{km}), comme un puzzle de plaques rigides dérivant sur l'asthénosphère, zone ductile du manteau. Les mouvements relatifs de ces plaques définissent trois types de limites. Les \textbf{limites divergentes}, où deux plaques s'écartent, correspondent aux dorsales océaniques (ex. : dorsale médio-atlantique, qui sépare l'Afrique de l'Amérique du Sud à \SI{2,5}{cm/an}). Les \textbf{limites convergentes}, où une plaque plonge sous une autre, correspondent aux zones de subduction (ex. : fosse du Pérou--Chili, où la plaque de Nazca plonge sous la plaque sud-américaine). Les \textbf{limites transformantes}, où deux plaques coulissent l'une contre l'autre, correspondent aux failles transformantes (ex. : faille de San Andreas en Californie). L'expansion océanique, démontrée par l'âge des fonds et les anomalies magnétiques, est ainsi le moteur visible de la dérive des plaques : elle explique l'ouverture de l'Atlantique, la séparation de l'Afrique et de l'Amérique du Sud il y a environ \SI{120}{Ma}, et plus généralement la mobilité de la surface terrestre.
\end{enumerate}
\textbf{Barème indicatif (sur 20)} : analyse des 4 observations 6 pts ; subduction justifiée et nommée 4 pts ; démonstration du mouvement des plaques 4 pts ; conclusion rédigée (définition de la lithosphère 2 pts, trois limites + exemples 4 pts) 6 pts.
\textbf{Points de vigilance} : la question 2 exige le raisonnement complet (création $\Rightarrow$ disparition obligatoire), pas seulement le nom des zones ; dans la conclusion, chaque type de limite doit être illustré par \emph{un exemple précis} ; veiller à définir la lithosphère en la distinguant de la croûte seule.
\end{corrige}

\newpage

\coursec{Fiche bilan}

\begin{popfigure}
\begin{tikzpicture}[mindmap, grow cyclic, every node/.style={concept, text=white, font=\small\bfseries},
  root concept/.append style={concept color=popblue, font=\normalsize\bfseries},
  level 1/.append style={level distance=3.6cm, sibling angle=72},
  level 1 concept/.append style={font=\footnotesize\bfseries}]
\node[root concept]{Fonds océaniques : reconstitution chronologique}
  child[concept color=popgreen]{ node {Dorsales : zones d'expansion} }
  child[concept color=poporange]{ node {Âge des fonds : croît en s'éloignant de l'axe} }
  child[concept color=poppurple]{ node {Anomalies magnétiques symétriques} }
  child[concept color=popgold]{ node {Inversions du champ magnétique terrestre} }
  child[concept color=poppink]{ node {Vitesse d'expansion : $v = d/t$} }
  child[concept color=popblue!70]{ node {Tectonique des plaques} };
\end{tikzpicture}
\legende{Carte mentale de la leçon : les marqueurs chronologiques des fonds océaniques (dorsales, âge, anomalies magnétiques) permettent de reconstituer l'expansion océanique, fondement de la tectonique des plaques.}
\end{popfigure}

\begin{bilan}
\textbf{Définitions clés}
\begin{itemize}
  \item \cle{Dorsale océanique} : chaîne de montagnes sous-marines où le magma basaltique remonte et forme en permanence une nouvelle croûte océanique (lithosphère jeune).
  \item \cle{Expansion océanique} : écartement progressif des fonds de part et d'autre de l'axe de la dorsale, à raison de quelques centimètres par an.
  \item \cle{Anomalie magnétique} : écart entre le champ magnétique mesuré et le champ théorique ; elle est \textbf{positive} si les basaltes ont enregistré un champ de même sens que l'actuel, \textbf{négative} en cas de champ inversé.
  \item \cle{Inversion du champ magnétique terrestre} : renversement des pôles magnétiques survenu de nombreuses fois au cours de l'histoire de la Terre, enregistré par les basaltes lors de leur refroidissement (aimantation rémanente).
  \item \cle{Plaque lithosphérique} : fragment rigide de la lithosphère, en déplacement à la surface du globe sur l'asthénosphère ductile.
\end{itemize}

\textbf{Formules et faits à connaître par cœur}
\begin{center}
\begin{tabularx}{\linewidth}{|l|X|}
\hline
\rowcolor{popblueL} \textbf{Notion} & \textbf{À retenir} \\
\hline
Vitesse d'expansion & $v = \dfrac{d}{t}$ avec $d$ = distance à l'axe de la dorsale et $t$ = âge du basalte (en cm/an ou km/Ma) \\
\hline
Âge des fonds & Nul à l'axe de la dorsale, croissant symétriquement vers les continents ; jamais plus de $\approx 200$ Ma (Jurassique) \\
\hline
Anomalies magnétiques & Bandes alternées $+$/$-$, parallèles à la dorsale et \textbf{symétriques} par rapport à son axe \\
\hline
Double marqueur & Concordance entre carte des âges et carte des anomalies : deux preuves indépendantes de l'expansion \\
\hline
\end{tabularx}
\end{center}

\textbf{Méthodes express}
\begin{itemize}
  \item \emph{Reconstituer une chronologie} : repérer l'axe de la dorsale $\rightarrow$ classer les bandes de la plus proche (la plus jeune) à la plus éloignée (la plus ancienne) $\rightarrow$ vérifier la symétrie.
  \item \emph{Calculer une vitesse} : mesurer la distance d'une anomalie (ou d'un basalte daté) à l'axe, diviser par son âge ; attention aux conversions (\SI{1}{cm/an} = \SI{10}{km/Ma}).
  \item \emph{Conclure} : toujours relier l'expansion océanique au mouvement divergent de deux plaques lithosphériques.
\end{itemize}
\end{bilan}

\begin{aretenir}[Checklist avant le Probatoire]
\begin{itemize}
  \item[$\square$] Je sais situer une dorsale sur une carte et expliquer son rôle dans la formation de la croûte océanique.
  \item[$\square$] Je sais que l'âge des fonds océaniques augmente en s'éloignant de l'axe de la dorsale.
  \item[$\square$] Je sais que les fonds océaniques actuels sont tous plus jeunes que $\approx 200$ Ma.
  \item[$\square$] Je sais définir une anomalie magnétique positive et négative.
  \item[$\square$] Je sais que les basaltes enregistrent le champ magnétique terrestre lors de leur refroidissement.
  \item[$\square$] Je sais que les anomalies magnétiques sont symétriques de part et d'autre de la dorsale.
  \item[$\square$] Je sais calculer une vitesse d'expansion avec $v = d/t$ et convertir les unités.
  \item[$\square$] Je sais exploiter une carte des âges des fonds océaniques pour reconstituer une chronologie.
  \item[$\square$] Je sais que la concordance âges/anomalies constitue une double preuve de l'expansion.
  \item[$\square$] Je sais relier l'expansion océanique à la divergence de deux plaques lithosphériques.
  \item[$\square$] Je sais rédiger une conclusion structurée : observation $\rightarrow$ interprétation $\rightarrow$ mécanisme.
\end{itemize}
\end{aretenir}

\findecours

\end{document}
